% !TEX root = poisson-cubic.tex
% !TeX root = ../../poisson-cubic.tex
\RequirePackage{fix-cm}
\documentclass[a4paper,11pt]{ltjsarticle}
\usepackage[margin=25mm,headheight=16pt]{geometry}
\usepackage{amsmath,amssymb,amsthm,mathtools,bm,mathrsfs}
\usepackage{booktabs,longtable,array,enumitem}
\usepackage{xcolor}
\usepackage[unicode,colorlinks=true,linkcolor=blue!40!black,citecolor=blue!40!black,urlcolor=blue!45!black]{hyperref}
\usepackage{fancyhdr}
\setlength{\parindent}{1em}
\setlength{\parskip}{0.25em}
\setlist{itemsep=0.2em,topsep=0.4em}
\allowdisplaybreaks[2]
\clubpenalty=10000
\widowpenalty=10000
\pagestyle{fancy}
\fancyhf{}
\fancyhead[L]{\small 3次超曲面のPoisson構造}
\fancyhead[R]{}
\fancyfoot[C]{\thepage}
\newcommand{\CC}{\mathbb C}
\newcommand{\ZZ}{\mathbb Z}
\newcommand{\PP}{\mathbb P}
\newcommand{\OO}{\mathcal O}
\newcommand{\cK}{\mathcal K}
\newcommand{\cP}{\mathcal P}
\newcommand{\cD}{\mathcal D}
\newcommand{\cC}{\mathcal C}
\newcommand{\cE}{\mathcal E}
\newcommand{\cF}{\mathcal F}
\newcommand{\cM}{\mathcal M}
\newcommand{\fm}{\mathfrak m}
\newcommand{\Perf}{\operatorname{Perf}}
\newcommand{\Ku}{\operatorname{Ku}}
\newcommand{\MF}{\operatorname{MF}}
\newcommand{\HH}{\operatorname{HH}}
\newcommand{\QGr}{\operatorname{QGr}}
\newcommand{\qgr}{\operatorname{qgr}}
\newcommand{\Dsg}{D_{\mathrm{sg}}}
\newcommand{\RHom}{\operatorname{RHom}}
\newcommand{\REnd}{\operatorname{REnd}}
\newcommand{\Hom}{\operatorname{Hom}}
\newcommand{\Ext}{\operatorname{Ext}}
\newcommand{\Ind}{\operatorname{Ind}}
\newcommand{\Def}{\operatorname{Def}}
\newcommand{\thick}{\operatorname{thick}}
\newcommand{\Spec}{\operatorname{Spec}}
\newcommand{\Spf}{\operatorname{Spf}}
\newcommand{\Id}{\operatorname{Id}}
\newcommand{\sym}{\operatorname{sym}}
\newcommand{\Cone}{\operatorname{Cone}}
\newcommand{\rank}{\operatorname{rank}}
\newcommand{\Pic}{\operatorname{Pic}}
\newcommand{\otL}{\mathbin{\otimes^{\mathbf L}}}
\theoremstyle{plain}
\newtheorem{theorem}{定理}[section]
\newtheorem{proposition}[theorem]{命題}
\newtheorem{lemma}[theorem]{補題}
\newtheorem{corollary}[theorem]{系}
\theoremstyle{definition}
\newtheorem{definition}[theorem]{定義}
\newtheorem{problem}[theorem]{問題}
\newtheorem{example}[theorem]{例}
\theoremstyle{remark}
\newtheorem{remark}[theorem]{注意}
\renewcommand{\proofname}{証明}
\hypersetup{pdfauthor={},pdfsubject={3次超曲面のPoisson正規形・退化軌跡・モジュライ}}

\newif\ifbuildmetadata
\InputIfFileExists{papers/build-metadata.tex}{\buildmetadatatrue}{\buildmetadatafalse}
\numberwithin{equation}{section}
\hypersetup{pdftitle={滑らかな3次超曲面のPoisson構造：正規形・退化軌跡・モジュライ}}
\fancyhead[L]{\small 3次超曲面のPoisson構造}
\ifbuildmetadata
  \fancyhead[R]{\small \RevisionDate}
\fi
\begin{document}
\hypersetup{pageanchor=false}
\begin{titlepage}
\centering
\vspace*{22mm}
{\LARGE\bfseries 滑らかな3次超曲面のPoisson構造\par}
\vspace{5mm}
{\Large 正規形・退化軌跡・モジュライ\par}
\vspace{16mm}
\begin{minipage}{.89\textwidth}
\begin{abstract}
本稿では、複素数体上の滑らかな射影3次超曲面 $X\subset\mathbb P^{n+1}$（$n\ge3$）上の代数的Poisson構造を全次元・全階数にわたって分類し、その退化軌跡とモジュライ空間の構造を決定する。

出発点として、$X$ 上の任意の大域双ベクトル場が定数3重ベクトル $P\in\bigwedge^3 V$ を用いて一意に表示されることを示し、非線形なJacobi恒等式を $X$ の定義方程式のHessianに対する線形条件へと帰着させる。さらに、$X$ の滑らかさから従うHessian系の基点自由性とLie代数の構造論を組み合わせることにより、全ての非零Poisson構造が3変数Jacobianブロックの直和に分解するという正規形定理を証明する。

この正規形を基盤として、以下の幾何学的・モジュライ的帰結を導く。
第一に、全階数の退化スキームの定義イデアルを明示的に記述し、滑らかな3次超曲面に対するBondal予想の次元下界を全次元・全退化階数で証明する。
第二に、3次四様体における一般階数4のPoisson構造は2つの滑らかな平面3次曲線の分離型に限られ、すべて純楕円型の対数symplectic構造となることを示す。あわせて、分離型四様体の非連結な零点曲線に沿って3重に消える3次切断について、消失次数を保った周囲の射影空間への持上げ障害を計算する。
第三に、Fermat型3次超曲面の固定Poissonスキームの既約成分を大きさ3以上の集合分割によって完全に記述し、相対モジュライの既約成分数が $\lfloor(n+2)/3\rfloor$ であることを示す。さらに、5次元において節点型・尖点型の特異ブロックを一般点に持つ固定Poisson成分を構成する。
\end{abstract}
\end{minipage}
\vfill
\ifbuildmetadata
  {\small 作成日時：\CreationTime\par}
  \vspace{3mm}
  {\small PDFファイル：\texttt{\PDFFileName}\par}
\fi
\vspace*{8mm}
\end{titlepage}
\pagenumbering{roman}
\hypersetup{pageanchor=true}
\tableofcontents
\clearpage
\pagenumbering{arabic}

\section{序論}\label{sec:paperA-introduction}
本稿の目的は、滑らかな複素射影3次超曲面上の代数的Poisson構造を完全に分類し、そこから従うsymplectic leaf、退化軌跡、およびモジュライ空間の幾何学を明らかにすることである。
滑らかな複素射影多様体 $Y$ 上の代数的Poisson構造とは、大域双ベクトル場 $\pi\in H^0(Y,\bigwedge^2T_Y)$ であって、Schouten括弧について可積分条件 $[\pi,\pi]=0$ を満たすものをいう。

\subsection*{射影完全交叉における3次超曲面の位置づけ}
射影完全交叉の中で非零Poisson構造を許容するものは極めて限られている。
次元 $n\ge3$ の滑らかな射影完全交叉 $Y=Y_{d_1,\ldots,d_c}\subset\mathbb P^{n+c}$ について、冗長な1次方程式を除いて $d_i\ge2$ と仮定すると、非零Poisson構造を持ち得る型は次の4種類にちょうど限られる：
\begin{itemize}
  \item 射影空間 $\mathbb P^n$、
  \item 2次超曲面 $Q^n\subset\mathbb P^{n+1}$、
  \item 3次超曲面 $Y_3\subset\mathbb P^{n+1}$、
  \item $(2,2)$型完全交叉 $Y_{2,2}\subset\mathbb P^{n+2}$。
\end{itemize}
この制限は、$s=\sum_i(d_i-1)$ とおいたときの随伴公式 $\bigwedge^2T_Y\simeq\Omega_Y^{n-2}(n+1-s)$ から直ちに見通せる。実際、$s\ge3$ のときには微分形式の低次数消滅定理 \cite[Lemma~4.1]{FuHwangCI} により大域双ベクトル場そのものが存在しない。これら4型それぞれにおける非零Poisson構造の存在と、ここで用いる消滅定理の初等的な証明は、付録の定理\ref{thm:ci-poisson-existence}および補題\ref{lem:ci-low-twist-vanishing}にまとめる。

これら4型のうち、射影空間 $\mathbb P^n$ や2次超曲面 $Q^n$ では、次元が上がるにつれてトーラス作用や高次の葉層構造に由来する多様なPoisson構造が現れ、全次元・全階数にわたる統一的分類は極めて複雑になる。これに対し、本稿の主結果は、滑らかな3次超曲面の全ての非零Poisson構造が、次元 $n\ge3$ と一般階数によらず「3変数Jacobianブロックの直和」という単一の正規形で尽くされることを示すものである。著者の知る限り、射影完全交叉の1つの型について、全ての滑らかなメンバーを全次元・全階数にわたって同一の正規形で分類した最初の結果である。

以下、$V$ を次元 $N=n+2\ge5$ の複素ベクトル空間とし、斉次3次式 $F\in\operatorname{Sym}^3V^*$ が滑らかな3次超曲面 $X=V(F)\subset\mathbb P(V)$ を定めるとする。

\subsection*{主結果1：全次元・全階数の正規形}
$F$ を固定すると、$X$ 上の任意の大域双ベクトル場 $\pi\in H^0(X,\bigwedge^2T_X)$ は、定数3重ベクトル $P\in\bigwedge^3V$ を用いて一意に
\begin{equation}\label{eq:intro-bivector-representation}
\pi_P=(\iota_{dF}P)|_X,\qquad P\in\bigwedge^3V
\end{equation}
と表示される（命題\ref{prop:pg-represent}）。ここで右辺は、アフィン錐 $V$ 上の2次係数双ベクトル場 $q=\iota_{dF}P$ を射影化して $X$ へ制限したものである。重要な点として、$\pi_P$ が $X$ 上でPoisson構造であることと、錐上の双ベクトル場 $q$ が $V$ 全体でPoissonであること（$[q,q]=0$）は同値となる（命題\ref{prop:pg-jacobi-extension}）。

この表示の下で、全ての非零Poisson構造は次の正規形を持つ（定理\ref{thm:pg-normal}）：ある整数 $k\ge1$、$m\ge0$（$N=3k+m$）と、$V^*$ の適切な線形座標 $x_s=(x_{s1},x_{s2},x_{s3})$（$1\le s\le k$）および $u=(u_1,\ldots,u_m)$ に対して、
\begin{equation}\label{eq:intro-block-normal-form}
P=\sum_{s=1}^k\partial_{x_{s1}}\wedge\partial_{x_{s2}}\wedge\partial_{x_{s3}},
\qquad F=\sum_{s=1}^k f_s(x_s,u)
\end{equation}
が成り立つ。ここで各 $f_s$ は $(x_s,u)$ に関する次数3の斉次多項式であり、中心変数 $u_\alpha$ は錐上の括弧 $\iota_{dF}P$ に関してCasimirとなる。逆に、式\eqref{eq:intro-block-normal-form}の形を持つ滑らかな3次式 $F$ は常にPoisson構造を定め、その一般階数はちょうど $2k$ である。なお、各ブロックの3変数制限 $f_s(x_s,0)$ には滑らかさを仮定せず、特異な平面3次曲線や恒等的な零多項式も許容される。

証明の鍵は、非線形なJacobi恒等式を「2次形式の線形系とLie代数の構造論」へ翻訳する点にある。具体的には、次の4段階で議論を進める。
\begin{enumerate}
  \item 勾配 $dF$ のKoszul複体の完全性とEuler--divergence恒等式を用いて、大域双ベクトル場の表示\eqref{eq:intro-bivector-representation}とJacobi恒等式のアフィン錐への持上げを確立する。
  \item 勾配の1次syzygyの消滅により、Jacobi恒等式が「任意の点 $p\in V$ における $F$ のHessianが、式\eqref{eq:pg-K}の4変数交代恒等式で定まる2次形式の線形空間 $K(P)$ に属する」という線形条件と同値になることを示す。
  \item $X$ の滑らかさによりHessianの線形系（従って $K(P)$ も）は $\mathbb P(V)$ 上に基点を持たない。$P$ の最小支持部分空間 $W\subset V$ 上で $K(P)$ から非退化な2次形式 $B$ を1つ選び、$P$ と $B$ を組み合わせて $W$ 上に完全（$[W,W]=W$）かつ中心が自明なLie代数構造を定義する。
  \item 可解根基内の極小可換イデアルに対する $\mathfrak{sl}_2$-加群と、階数2以上の単純因子の最高根ベクトルを式\eqref{eq:pg-K}に代入すると、いずれも $K(P)$ の共通零点（基点）を生じるため矛盾する。残るLie代数は $\mathfrak{sl}_2^{\oplus k}$ に限られ、その各3次元単純因子が正規形\eqref{eq:intro-block-normal-form}のブロックを与える。
\end{enumerate}

この正規形の直截な帰結として、固定した滑らかな3次超曲面 $X$ 上で一般階数2の非零Poisson構造が張る射影パラメータ空間は、被約構造を入れればGrassmann多様体 $\operatorname{Gr}(3,V)$ に一致する。また、一般の滑らかな3次超曲面では非零Poisson構造はこの一般階数2のものに限られ、複数ブロック（$k\ge2$）を持つ構造は特別な3次超曲面にのみ現れる（系\ref{cor:pg-generic}）。

\subsection*{主結果2：退化軌跡、Bondal予想、および3次四様体}
Poisson構造 $\pi$ の階数が $2j$ 以下となる退化スキームを $D_{2j}(\pi)=V(\pi^{j+1})$ と書く。正規形\eqref{eq:intro-block-normal-form}を用いると、任意の $j$ に対して $D_{2j}(\pi)$ の定義イデアルが、各ブロックの偏微分とEuler場を用いて明示的に記述できる（命題\ref{prop:pg-deg-ideal}）。

この記述から、滑らかな3次超曲面に対するBondal予想が全次元・全階数で従う。すなわち、$j<k$ の場合には $j+1$ 個のブロックから選んだ平面曲線の射影的joinを構成し、$j\ge k$ の場合には $D_{2j}(\pi)=X$ であることを用いると、全ての $0\le2j<n$ について次元下界
\begin{equation}\label{eq:intro-bondal}
\dim D_{2j}(\pi)\ge2j+1
\end{equation}
が成り立つ（定理\ref{thm:pg-bondal}）。また、各3変数制限 $f_s(x_s,0)$ が滑らかな平面3次曲線 $E_s$ を定める場合には、零点スキーム $D_0(\pi_X)$ は互いに交わらない楕円曲線 $\coprod_s E_s$ と残余スキーム $Z$ の非交和に分解する。残余スキーム $Z$ は $m\le1$ ならば空であり、$m\ge2$ ならば中心空間 $\mathbb P^{m-1}$ 上の次数 $8^k$ の有限平坦被覆から得られる完全交叉となる（命題\ref{prop:pg-deg-support}、\ref{prop:pg-zero-ci}）。

4次元の3次四様体（$n=4$、$N=6$）では、一般階数4（すなわち一般点で非退化）の構造は $k=2$、$m=0$ に対応するため、方程式が
\[
F=f(x_0,x_1,x_2)+g(y_0,y_1,y_2)
\]
と分離型に書けることと必要十分になる。このとき $f,g$ はともに滑らかな平面3次曲線を定め、退化因子 $D=V(f,g)$ は2つの楕円曲線のjoinとなる。局所Pfaffianとmodularベクトル場を直接計算することにより、3次四様体上の一般階数4のPoisson構造は例外なく purely elliptic log symplectic であることを証明する（命題\ref{prop:pg-split-elliptic}、系\ref{cor:pg-cubic-fourfold-elliptic}）。さらに、Bricalli--Favale--Pirola \cite[Theorem~3.4]{PBricalliFavalePirola} の結果と合わせると、この $(3,3)$ 分離条件は、周囲の $\mathbb P^5$ における $F$ のHessianの階数3以下の軌跡が線形平面 $\mathbb P^2$ を含むこととも同値である。なお、Hessian超曲面の可約性だけでは階数4の存在を判定できないことも具体例で確認する（例\ref{ex:pg-reducible-hessian-no-rank-four}）。

\subsection*{主結果3：固定・相対モジュライと非楕円ブロック}
モジュライの記述では、Poisson構造をスカラー倍で割らずに零構造も含めて扱う。固定した $X$ に対し、Jacobi方程式 $[\pi,\pi]=0$ の零スキームを $\operatorname{Pois}(X)$ とし、自己同型群による商スタックを $\mathfrak{Pois}(X)=[\operatorname{Pois}(X)/\operatorname{Aut}(X)]$ と定める。同様に、滑らかな3次式のパラメータ空間 $U_N\subset\mathbb P(\operatorname{Sym}^3V^*)$ 上で組 $([F],\pi)$ をパラメータ化する相対Poissonスキームを $\mathscr P_n$ と書き、相対モジュライスタックを $\mathfrak{Pois}^{\mathrm{rel}}_n=[\mathscr P_n/\mathrm{PGL}(V)]$ と定める。これらはいずれも有限型Deligne--Mumfordスタックとなる（第\ref{sec:pg-poisson-stacks}節）。

特に、Fermat型3次超曲面 $X_N=V(\sum_{i=1}^N z_i^3)$ に対しては、$\operatorname{Pois}(X_N)$ の既約成分分解を完全に決定できる。$\{1,\ldots,N\}$ を大きさ3以上の部分に分ける非順序集合分割 $\mathcal B=\{B_1,\ldots,B_k\}$ を考え、$z_i$ の双対基底 $e_i$ に対して $V_B=\langle e_i:i\in B\rangle$ とおくと、Pl\"ucker錐 $\widehat{\operatorname{Gr}}(3,V_B)$ の和によって
\begin{equation}\label{eq:intro-fermat-components}
Z_{\mathcal B}=\sum_{B\in\mathcal B}\widehat{\operatorname{Gr}}(3,V_B),
\qquad\operatorname{Pois}(X_N)_{\mathrm{red}}=\bigcup_{\mathcal B}Z_{\mathcal B}
\end{equation}
が既約成分分解を与える。各成分 $Z_{\mathcal B}$ は次元 $3N-8k$、一般Poisson階数 $2k$ を持つ。$\operatorname{Aut}(X_N)$ の作用は同じ大きさの構成を持つ集合分割を同一視するため、固定スタックの既約成分は $N$ の各部分が3以上の整数分割で添字付けられる。一方、相対モジュライでは一般階数 $2k$ ごとに既約成分がちょうど1つ存在し、既約成分数はそれぞれ
\begin{equation}\label{eq:intro-poisson-stack-counts}
\#\operatorname{Irr}(\mathfrak{Pois}(X_N))
=[q^N]\prod_{b\ge3}(1-q^b)^{-1},\qquad
\#\operatorname{Irr}(\mathfrak{Pois}^{\mathrm{rel}}_n)
=\left\lfloor\frac{n+2}{3}\right\rfloor
\end{equation}
と求まる（定理\ref{thm:pg-fermat-relative-count}）。

また、正規形\eqref{eq:intro-block-normal-form}に現れる3変数制限 $f_s(x_s,0)$ が特異になり得る現象を具体化するため、5次元の滑らかな3次超曲面上に、節点型または尖点型の平面3次曲線を一般点のブロックとして持つ固定Poisson成分を構成する（定理\ref{ne:main}）。この構成では、一般点におけるJacobi方程式の線形化の階数が33であることを整数上の小行列式によって厳密に検証し、その成分の滑らかさとブロックの内在性を示す。

\subsection*{先行研究との関係}
代数的Poisson幾何の体系的基礎はPolishchuk \cite{PPolishchuk} によって築かれ、Poissonスキーム・Poisson加群・退化軌跡が代数幾何学的に定式化された。曲面では $\bigwedge^2T_Y=K_Y^{-1}$ のためJacobi条件が自動的に成り立ち、反標準因子と双有理幾何が分類の主軸となる。曲面から高次元にわたる構成と分類の全体像はPymの概説 \cite[\S\S2--5]{PPymSurvey} にまとめられている。3次元では、Druel \cite[p.~230]{DruelThreefold} が有限個の点を除いて階数2である非零Poisson構造を持つ滑らかな射影三様体を分類し、実際には至る所階数2であることを示した。

3次元射影空間 $\mathbb P^3$ では、同型 $\bigwedge^2T_{\mathbb P^3}\simeq\Omega^1_{\mathbb P^3}(4)$ を通じてPoisson構造の分類が余次元1・次数2の葉層構造の分類に帰着する。Cerveau--Lins Neto \cite{CerveauLinsNeto} は、$\mathbb P^n$（$n\ge3$）上の次数2・余次元1葉層構造の空間が6個の既約成分を持つことを証明した。零点に因子成分を許すPoisson構造への適用およびPicard数1の滑らかなFano三様体への拡張はLoray--Pereira--Touzet \cite[\S9, Theorem~9.1]{LPTPoissonCubic} によって与えられ、非零構造をスカラー倍で割った空間の既約成分数は $\mathbb P^3$ で6、$Q^3$ で3である。滑らかな3次三様体の分類は同論文のTheorem~7.1に含まれており、Mayanskiy \cite[Theorem~1]{PMayanskiy} はSchouten括弧の直接計算によって、その射影PoissonスキームがPl\"ucker埋め込みされた $\operatorname{Gr}(2,5)$ に一致することを記述した。

高次元の射影空間や2次超曲面では、分類すべき構造の多様性が急激に増す。例えば両者のトーラス作用に由来する可換な大域ベクトル場 $v_1,\ldots,v_r$ から、任意の定数 $c_{ij}$ に対して $\sum_{i<j}c_{ij}v_i\wedge v_j$ というPoisson構造が得られ、トーラスの次元とともに可能な一般階数が増大する。さらに射影空間の階数2の問題には、次数2葉層構造の有理型・対数型・Lie代数作用型・線形引戻し型が混在し、Corr\^ea--Muniz \cite[Theorem~A, Corollary~C]{CorreaMuniz} は任意次元の次数2葉層の構造定理を与えて $\mathbb P^4$ の一般階数2のPoisson構造に適用している。

任意次元における大域的結果としては、Lima--Pereira \cite[Theorems~1--2]{LimaPereiraDiagonal} が全ての $n\ge3$ で $\mathbb P^n$ 上の対角Poisson構造の閉包が1つの既約成分をなすことを示し、さらに偶数次元4以上でPicard群が $\mathbb Z$ のFano多様体上で、退化因子が被約な単純正規交叉因子となる一般点非退化Poisson構造は射影空間の対角型に限ることを証明した。また、Druel--Pereira--Pym--Touzet \cite[Theorem~1.1]{GlobalWeinstein} は、コンパクトK\"ahler Poisson多様体が有限基本群を持つコンパクトなsymplectic leafを持つとき、有限\'etale被覆がその葉の普遍被覆と別のPoisson多様体の積に分解することを示した。

4次元射影空間 $\mathbb P^4$ については、Matviichuk--Pym--Schedler \cite[Theorems~1.2--1.4]{MPSLocalTorelli} が単純正規交叉の退化因子を持つlog symplectic構造の変形理論を任意次元で構成し、$\mathbb P^4$ のPoisson構造の空間でこの軌跡と交わる40個の既約成分を記述した。さらに公開草稿 \cite{PoissonP4Draft} では、$\mathrm{PGL}_5$ の商を取る前の非零Poisson構造の射影パラメータ空間が、被約既約成分147個（一般階数4の141成分と一般階数2の6成分）からなるという分類が提案されている\footnote{この草稿の証明に対する人間の検証は未了である。}。

手法面では、射影空間のPoisson構造をdivergence-freeな斉次双ベクトル場へ持ち上げる議論がBondal \cite[\S4]{PBondal93} に始まる。本稿ではこのEuler--divergenceの方法を3次超曲面上で消えるJacobiatorへ適用する。また、半単純Lie代数上の4変数恒等式の解空間はDerdzinski--Gal \cite[(1.4), Theorem~B]{PDerdzinskiGal} にも現れるが、第\ref{sec:pg-normal}節ではHessianの基点自由性から直接に可解根基と階数2以上の単純因子を排除する。

Bondal予想については、Fano四様体の場合は \cite[Theorem~7.13]{BondalGualtieriPym}、Fano五様体の場合は \cite[Corollary~1.3]{BondalDruelPereiraPymTouzet} で証明されている。本稿のjoin構成による証明は、3次超曲面に対して任意次元・全退化階数を一挙に扱う。一方で、分離型3次四様体が純楕円型となる事実（命題\ref{prop:pg-split-elliptic}）は、Pym \cite[Theorem~1.3]{PPymElliptic} のarXiv v1における3次四様体上の非存在主張とは両立しない。命題\ref{prop:pg-triple-lifting-obstruction}では、非連結な零点曲線に沿って高次消失を保った持上げを行う際のコホモロジー障害を計算し、同版の持上げ補題（Lemma~5.5）がこの非連結曲線に適用できない理由を注意\ref{rem:pg-split-elliptic-lift-obstruction}で明らかにする。

\subsection*{本稿の構成}
第\ref{sec:pg-extension}節と第\ref{sec:pg-normal}節で本稿の中心であるPoisson正規形定理を証明する。第\ref{sec:pg-leaves}節と第\ref{sec:pg-degeneracy}節では、正規形からCasimir・対数形式・symplectic leafの閉包・全退化スキームの定義イデアルを導出する。第\ref{sec:pg-fourfold}節では3次四様体の一般階数4の構造、Hessian軌跡との同値性、および非連結零点曲線に沿う持上げ障害を詳しく調べる。第\ref{sec:pg-bondal}節では滑らかな3次超曲面に対するBondal予想を証明する。第\ref{sec:pg-poisson-stacks}節から第\ref{sec:pg-special}節では、固定および相対モジュライスタック、Fermat型超曲面の成分分解、非楕円ブロックを持つ固定成分、ならびに特殊なブロック制限を扱う。最後に付録\ref{sec:general-poisson-existence}で、射影完全交叉・トーリック多様体・射影斉次多様体におけるPoisson構造の存在判定をまとめる。なお、本稿の正規形に基づく代数的量子化については \cite{CubicQuantization} を参照されたい。

% !TeX root = ../../poisson-cubic.tex
\section{大域的双ベクトル場と射影空間への延長}\label{sec:pg-extension}

本節では、滑らかな3次超曲面 $X=V(F)\subset\PP(V)$ 上のPoisson構造を線形代数と可換環論の土台に乗せるための2つの基本結果を示す。第一に、$X$ 上の任意の大域双ベクトル場が周囲の射影空間 $\PP(V)$ へ一意に延長され、さらに定数3重ベクトル $P\in\bigwedge^3V$ との縮約 $\iota_{dF}P$ として一意に表示されることを示す（命題\ref{prop:pg-represent}）。第二に、$X$ 上でのJacobi恒等式 $[\pi_P,\pi_P]=0$ がアフィン錐 $V$ 上でのJacobi恒等式 $[\iota_{dF}P,\iota_{dF}P]=0$ へそのまま持ち上がることを示す（命題\ref{prop:pg-jacobi-extension}）。

\paragraph{記号と規約}
$V$ を次元 $N=n+2\ge5$ の複素ベクトル空間とし、多項式環を $S=\operatorname{Sym}(V^*)$、射影空間を $\mathbb P=\PP(V)$（$V$ の1次元部分空間全体の空間）と書く。斉次3次式 $F\in S_3$ は滑らかな3次超曲面 $X=V(F)\subset\mathbb P$ を定めるものとし、$V^*$ の線形座標を $z_1,\ldots,z_N$、Eulerベクトル場を $E=\sum_{i=1}^N z_i\partial_{z_i}$ と書く。また偏微分を $F_i=\partial_i F$、$F_{ij}=\partial_i\partial_j F$ と略記する。

1-形式 $\eta\in V^*$ による多重ベクトルの縮約は
\begin{equation}\label{eq:pg-contraction-convention}
\iota_\eta(v_1\wedge\cdots\wedge v_r)
=\sum_{a=1}^r(-1)^{a-1}\eta(v_a)\,
v_1\wedge\cdots\wedge\widehat v_a\wedge\cdots\wedge v_r
\end{equation}
で定める。Poisson構造とは、Schouten括弧について $[\pi,\pi]=0$ を満たす大域双ベクトル場 $\pi\in H^0(X,\bigwedge^2T_X)$ のことである。以下では、アフィン錐 $V$ 上の2次係数双ベクトル場を $q$、その $\mathbb P$ への射影化を $\pi$、さらに $X$ への制限を $\pi_X$ と区別して書く。

なお、$X$ の滑らかさから $F$ は自動的に既約である。実際、もし正次数の因子分解 $F=GH$ があれば、$\dim\mathbb P=N-1\ge4$ より交叉 $V(G)\cap V(H)\subset\mathbb P$ は非空となり、その各点で $dF=GdH+HdG=0$ となって滑らかさに反する。

\subsection{勾配のKoszul複体とEuler列}
大域双ベクトル場の表示に必要となる2つの補題——勾配 $F_1,\ldots,F_N$ の正則列性および射影空間のEuler列からのコホモロジー計算——を準備する。

\begin{lemma}\label{lem:pg-gradient}
偏微分 $F_1,\ldots,F_N$ は多項式環 $S$ の正則列をなし、縮約 $\iota_{dF}$ を微分とするKoszul複体は正のホモロジー次数で完全である。特に、$\sum_{i=1}^N a_i F_i=0$ を満たす各 $a_i\in S$ が定数または1次斉次式ならば、すべての $a_i$ は零である。
\end{lemma}
\begin{proof}
Euler等式 $\sum_{i=1}^N z_i F_i=3F$ と $X$ の滑らかさから、$F_1,\ldots,F_N$ の共通零点はアフィン原点のみである。従ってイデアル $(F_1,\ldots,F_N)$ は極大イデアル $\mathfrak m=(z_1,\ldots,z_N)$-準素であり、$F_1,\ldots,F_N$ は $N$ 次元Cohen--Macaulay局所環 $S_{\mathfrak m}$ のパラメータ系、ゆえに正則列をなす \cite[Tags~00ND, 02JN]{PStacks}。

この正則性は多項式環 $S$ 全体でも成り立つ。実際、$S/(F_1,\ldots,F_{i-1})$ 上の $F_i$ 倍写像の核は斉次加群である。任意の非零斉次元は、定数項が非零の多項式を掛けても零にならないため、$\mathfrak m$ での局所化によって零にならない。局所化した核が零であることから、元の核も零である。

正則列のKoszul複体は正の次数で完全である \cite[Tag~062F]{PStacks}。各 $F_i$ の次数は $2$ であるから、第1 syzygy は自明な関係式 $F_j e_i - F_i e_j$（係数次数 $2$）で生成される。従って係数次数 $0$ または $1$ の非自明な斉次syzygyは存在せず、最後の主張を得る。
\end{proof}

\begin{lemma}\label{lem:pg-euler}
射影空間 $\mathbb P=\PP(V)$（$\dim\mathbb P=N-1\ge4$）について、次が成り立つ：
\begin{equation}\label{eq:pg-euler-vanishing}
H^0(\mathbb P,\textstyle\bigwedge^2T_{\mathbb P}(-3))
=H^1(\mathbb P,\textstyle\bigwedge^2T_{\mathbb P}(-3))=0.
\end{equation}
また、アフィン錐上の2次係数双ベクトル場の射影化によって、自然な同型
\begin{equation}\label{eq:pg-quadratic-euler-quotient}
H^0(\mathbb P,\textstyle\bigwedge^2T_{\mathbb P})
\simeq
\frac{S_2\otimes\bigwedge^2V}
{E\wedge(S_1\otimes V)}
\end{equation}
を得る。
\end{lemma}
\begin{proof}
Euler列とその2次外積を $\OO(-3)$ で捻ると、2つの短完全列
\begin{equation}\label{eq:pg-euler-sequences}
\begin{aligned}
0&\longrightarrow\OO(-3)\longrightarrow V\otimes\OO(-2)
\longrightarrow T_{\mathbb P}(-3)\longrightarrow0,\\
0&\longrightarrow T_{\mathbb P}(-3)
\longrightarrow\bigwedge^2V\otimes\OO(-1)
\longrightarrow\bigwedge^2T_{\mathbb P}(-3)\longrightarrow0
\end{aligned}
\end{equation}
が得られる。$\dim\mathbb P=N-1\ge4$ であるため、射影空間の線束の中間コホモロジー消滅 \cite[Tag~01XT]{PStacks} により、第1列から $i=0,1,2$ に対して $H^i(\mathbb P,T_{\mathbb P}(-3))=0$ が従う。さらに第2列の長完全列から消滅\eqref{eq:pg-euler-vanishing}を得る。

同様に、捻らないEuler列からは $H^1(\mathbb P,T_{\mathbb P})=0$ および全射 $S_1\otimes V\twoheadrightarrow H^0(\mathbb P,T_{\mathbb P})$ が従う。よって外積Euler列
\[
  0\longrightarrow T_{\mathbb P}
  \longrightarrow\bigwedge^2V\otimes\OO(2)
  \longrightarrow\bigwedge^2T_{\mathbb P}\longrightarrow0
\]
の大域切断を取れば同型\eqref{eq:pg-quadratic-euler-quotient}を得る。ここで左側の単射は、アフィン錐上ではベクトル場 $Z\in S_1\otimes V$ に対する外積 $Z\mapsto E\wedge Z$ で与えられる。
\end{proof}

以上の2補題を組み合わせることで、$X$ 上の全ての大域双ベクトル場が定数3重ベクトルによって一意に表示されることが従う。

\begin{proposition}\label{prop:pg-represent}
滑らかな3次式 $F$ を固定すると、自然な線形同型
\begin{equation}\label{eq:pg-representation-isomorphism}
\bigwedge^3V\xrightarrow{\sim}H^0(X,\bigwedge^2T_X),
\qquad P\longmapsto\pi_P=(\iota_{dF}P)|_X
\end{equation}
が存在する。ここで右辺は、アフィン錐上の2次係数双ベクトル場 $q=\iota_{dF}P$ を $\mathbb P$ へ射影化して $X$ へ制限したものである。さらに、$X$ 上の任意の大域双ベクトル場は周囲の射影空間 $\mathbb P$ 上の大域双ベクトル場へ一意に延長される。
\end{proposition}
\begin{proof}
\paragraph{ステップ1：$\mathbb P$ への一意な延長と接性条件}
超曲面 $X=V(F)$ の制限短完全列に $\bigwedge^2T_{\mathbb P}$ をテンソルし、補題\ref{lem:pg-euler}の消滅\eqref{eq:pg-euler-vanishing}を適用すると、制限写像
\begin{equation}\label{eq:pg-ambient-restriction-isomorphism}
H^0(\mathbb P,\bigwedge^2T_{\mathbb P})
\xrightarrow{\sim}
H^0(X,(\bigwedge^2T_{\mathbb P})|_X)
\end{equation}
が同型であることが分かる。接束の包含 $T_X\hookrightarrow T_{\mathbb P}|_X$ の外積により、与えられた $\pi_X\in H^0(X,\bigwedge^2T_X)$ を右辺の元と見なすと、その一意な延長 $\pi\in H^0(\mathbb P,\bigwedge^2T_{\mathbb P})$ が定まる。式\eqref{eq:pg-quadratic-euler-quotient}により、$\pi$ は2次係数双ベクトル場 $B\in S_2\otimes\bigwedge^2V$ で代表される。

次に $v=\iota_{dF}B\in S_4\otimes V$ とおく。法線写像 $T_{\mathbb P}|_X\to\OO_X(3)$ は $dF$ との縮約で与えられ、その外積への縮約の核がちょうど $\bigwedge^2T_X$ である。従って $\pi_X$ が $X$ に接する条件は、$v|_X$ が $T_{\mathbb P}(3)|_X$ の切断として零になることと同値である。制限したEuler列
\begin{equation}\label{eq:pg-twisted-euler-on-X}
0\longrightarrow\OO_X(3)
\xrightarrow{\ E\ }V\otimes\OO_X(4)
\longrightarrow T_{\mathbb P}(3)|_X\longrightarrow0
\end{equation}
より、$v|_X=h_X E$ を満たす $h_X\in H^0(X,\OO_X(3))$ が存在する。$H^1(\mathbb P,\OO)=0$ なので $h_X$ は $h\in S_3$ へ持ち上がり、ある1次係数ベクトル場 $A\in S_1\otimes V$ を用いて
\begin{equation}\label{eq:pg-tangency-correction}
\iota_{dF}B=FA+hE
\end{equation}
と書ける。この両辺をさらに $dF$ で縮約すると、$\iota_{dF}^2=0$ とEuler等式 $\iota_{dF}E=3F$ から $0=FA(F)+3Fh$、ゆえに $h=-A(F)/3$ を得る。そこで代表元を
\begin{equation}\label{eq:pg-horizontal-representative}
q=B-\frac13E\wedge A,\qquad\iota_{dF}q=0
\end{equation}
と取り替える。$E\wedge A$ は射影化で消えるため、$q$ は同じ延長 $\pi$ を代表し、かつアフィン錐上で厳密に水平（$\iota_{dF}q=0$）である。

\paragraph{ステップ2：Koszul完全性による全射性と単射性}
補題\ref{lem:pg-gradient}のKoszul複体の適切な斉次部分を取り出すと、完全列
\begin{equation}\label{eq:pg-graded-koszul-representation}
0\longrightarrow S_0\otimes\bigwedge^3V
\xrightarrow{\iota_{dF}}S_2\otimes\bigwedge^2V
\xrightarrow{\iota_{dF}}S_4\otimes V
\end{equation}
を得る（左端の直前は $S_{-2}\otimes\bigwedge^4V=0$ なので、左の写像は単射である）。$\iota_{dF}q=0$ より $q=\iota_{dF}P$ を満たす定数3重ベクトル $P\in\bigwedge^3V$ が一意に存在し、写像\eqref{eq:pg-representation-isomorphism}の全射性が従う。

逆に、任意の定数 $P\in\bigwedge^3V$ に対して $q=\iota_{dF}P$ は2次係数を持ち $\iota_{dF}q=0$ を満たすため、その射影化は $X$ に接する双ベクトル場 $\pi_P$ を線形に定める。もし $\pi_P=0$ ならば、同型\eqref{eq:pg-ambient-restriction-isomorphism}により $q$ の射影化は $\mathbb P$ 上で零、すなわちアフィン錐上で $E\wedge q=0$ である。これを $dF$ で縮約すると $0=3Fq-E\wedge\iota_{dF}q=3Fq$ より $q=0$ となり、完全列\eqref{eq:pg-graded-koszul-representation}の単射性から $P=0$ を得る。これで同型性と延長の一意性が完全に証明された。
\end{proof}

なお、$a\in\CC^\times$ に対して定義方程式を $F$ から $aF$ へスカラー倍すると、命題\ref{prop:pg-represent}の同型は $a$ 倍される。

\paragraph{構造層上のPoisson括弧の明示式}
アフィン錐上の2次係数双ベクトル場 $q=\sum_{i<j}q^{ij}\partial_i\wedge\partial_j$ に対し、多項式上の双導分を $\{a,b\}_q=\sum_{i,j}q^{ij}(\partial_i a)(\partial_j b)$ で定める。係数の2次斉次性により、この双導分は局所化環 $S[z_i^{-1}]$ の次数0部分（アフィンチャート $U_i=D_+(z_i)$ の座標環）を保つ。実際、アフィン座標 $t_j=z_j/z_i$ の間の括弧は
\begin{equation}\label{eq:pg-projective-ratio-bracket}
\{t_j,t_k\}_q
=\frac{z_iq^{jk}-z_jq^{ik}+z_kq^{ij}}{z_i^3}
\end{equation}
で与えられる。重なり上の括弧は同一の局所化双導分の制限であるため自然に貼り合わさる。また $\{F,-\}_q=0$ より、$U_i$ 上で $X$ を定める方程式 $F/z_i^3$ が生成するイデアルは括弧で保たれるため、括弧は $X$ の構造層へ降下し、式\eqref{eq:pg-projective-ratio-bracket}が $\pi_P$ の誘導する括弧を与える。特に、アフィン錐上でJacobi恒等式 $[q,q]=0$ が成り立てば、射影化と $X$ への制限後もJacobi恒等式が成り立つ。

\subsection{Jacobi恒等式のアフィン錐への持上げ}
次に、逆向き——「$X$ 上で $\pi_P$ がPoissonならば、アフィン錐上で $q=\iota_{dF}P$ 自体がPoissonである」——を証明する。ここでBondal \cite[\S4, pp.~11--12]{PBondal93} に由来するEuler--divergence恒等式が鍵となる。

固定した体積形式 $dz_1\wedge\cdots\wedge dz_N$ に関する多重ベクトル場のdivergence作用素 $\delta$ を
\begin{equation}\label{eq:pg-divergence-definition}
\delta(a\partial_{i_1}\wedge\cdots\wedge\partial_{i_r})
=\sum_{b=1}^r(-1)^{b-1}(\partial_{i_b}a)\,
\partial_{i_1}\wedge\cdots\wedge\widehat{\partial_{i_b}}\wedge\cdots
\wedge\partial_{i_r}
\end{equation}
と定める。係数次数 $d$、外積次数 $r$ の斉次多重ベクトル場 $C$ に対して、次の基本恒等式が成り立つ：
\begin{equation}\label{eq:pg-divergence-identities}
\begin{aligned}
\delta(FC)&=\iota_{dF}C+F\delta C,\qquad \delta^2=0,\\
\delta(E\wedge C)&=(N+d-r)C-E\wedge\delta C,\\
\delta\iota_{dF}C&=-\iota_{dF}\delta C.
\end{aligned}
\end{equation}
第2式は、$E$ への微分が $N-r$ を、$C$ の係数へのEuler微分が $d$ を生じることから従う。第3式では、$F$ の2階偏微分を含む項がHessianの対称性と $C$ の交代性により相殺し、残りの項は縮約の順序交換により符号が反転する。

\begin{proposition}\label{prop:pg-jacobi-extension}
定数3重ベクトル $P\in\bigwedge^3V$ と $q=\iota_{dF}P$ について、$\pi_P$ が $X$ 上のPoisson構造であることと、アフィン錐上で $[q,q]=0$ が成り立つことは同値である。従って、$X$ 上の任意のPoisson構造の $\mathbb P$ への一意な延長もまたPoisson構造である。
\end{proposition}
\begin{proof}
逆向き（$[q,q]=0\Rightarrow[\pi_P,\pi_P]=0$）は上で確認した括弧の降下から直ちに従う。以下、$\pi_P$ が $X$ 上でPoissonであると仮定し、錐上のJacobiator $J=\tfrac12[q,q]$ が恒等的に零であることを示す。

\paragraph{ステップ1：$J$ の基本性質（divergence-freeかつ水平）}
$J$ の成分表示は
\begin{equation}\label{eq:pg-jacobiator-coefficients}
J^{ijk}=\sum_a\left(
q^{ia}\partial_aq^{jk}+q^{ja}\partial_aq^{ki}
+q^{ka}\partial_aq^{ij}\right)
\end{equation}
であり、係数は3次斉次多項式である。$q^{ij}=\sum_a P^{ija}F_a$ に対して $P^{ija}$ の交代性と $F_{ai}$ の対称性から $\delta q=0$ が成り立つ。Schouten括弧とdivergenceの積の微分則より $J=(\delta q)\wedge q-\tfrac12\delta(q\wedge q)$ であるため、$\delta^2=0$ から $\delta J=0$ を得る。また $F$ は $q$ のCasimirなので、Jacobiatorに $F$ を代入した各項がすべて消え、$\iota_{dF}J=0$ が成り立つ。

\paragraph{ステップ2：超曲面上の消滅から $E\wedge J=FQ$ を導く}
Jacobiatorは交代3重導分であり、アフィンチャート $U_i=D_+(z_i)$ 上の座標 $t_j=z_j/z_i$ に式\eqref{eq:pg-projective-ratio-bracket}を適用すると
\begin{equation}\label{eq:pg-projective-ratio-jacobi}
J(t_j,t_k,t_l)=
\frac{z_iJ^{jkl}-z_jJ^{ikl}+z_kJ^{ijl}-z_lJ^{ijk}}{z_i^4}
=\frac{(E\wedge J)^{ijkl}}{z_i^4}
\end{equation}
を得る。左辺は射影化された $J$ の局所成分全体であるから、$[\pi_P,\pi_P]=0$ は $(E\wedge J)|_X=0$ を意味する（実際、$U_i$ 上ではEuler直線が局所直和因子をなすため、商の外積からの $E\wedge-$ は単射である）。従って $E\wedge J$ の各4次斉次係数は、制限写像 $S_4\to H^0(X,\OO_X(4))$ の核 $FS_1$ に属する。次数を比較すれば、ある1次係数の4重ベクトル場 $Q\in S_1\otimes\bigwedge^4V$ が存在して
\begin{equation}\label{eq:pg-euler-jacobi-divisibility}
E\wedge J=FQ
\end{equation}
と書ける。

\paragraph{ステップ3：Euler--divergence恒等式による $J=0$ の導出}
式\eqref{eq:pg-euler-jacobi-divisibility}の両辺を $dF$ で縮約し、$\iota_{dF}J=0$ と $\iota_{dF}E=3F$ を用いると $F\iota_{dF}Q=3FJ$ を得る。$S$ は整域なので $F$ で割って
\begin{equation}\label{eq:pg-jacobi-quotient-horizontal}
\iota_{dF}Q=3J,\qquad
\iota_{dF}(\delta Q)=-\delta(\iota_{dF}Q)=-3\delta J=0
\end{equation}
が従う（ここで式\eqref{eq:pg-divergence-identities}の第3式と $\delta J=0$ を用いた）。$Q$ の係数は1次式なので $\delta Q$ は定数3重ベクトルであり、補題\ref{lem:pg-gradient}の定数係数syzygyの消滅から $\delta Q=0$ を得る。

最後に式\eqref{eq:pg-euler-jacobi-divisibility}へ $\delta$ を作用させる。$J$ は係数次数 $d=3$、外積次数 $r=3$ かつ $\delta J=0$ なので、式\eqref{eq:pg-divergence-identities}より
\begin{equation}
NJ=\delta(E\wedge J)=\delta(FQ)
=\iota_{dF}Q+F\delta Q=3J
\end{equation}
となる。仮定 $N\ge5$ より $N-3\neq 0$ であり、ゆえに $J=0$ を得る。
\end{proof}

\section{Hessian による Poisson 構造の正規形}\label{sec:pg-normal}

前節の結果により、滑らかな3次超曲面 $X=V(F)$ 上のPoisson構造の分類は、「定数3重ベクトル $P\in\bigwedge^3V$ に対し、アフィン錐上で $[\iota_{dF}P,\iota_{dF}P]=0$ が成り立つための条件を決定する」という純粋に代数的な問題へ帰着された。

本節では、この非線形条件を $F$ のHessianに対する線形条件へと変換し（第\ref{sec:pg-normal}.1節）、$X$ の滑らかさから従う「2次形式の系の基点自由性」とLie代数の構造論を組み合わせて $P$ が3次元ブロックの直和に分解することを証明する（第\ref{sec:pg-normal}.2節）。最後に、$F$ 自体も同じブロック構造に沿って分離することを示し、本稿の主定理であるPoisson正規形（定理\ref{thm:pg-normal}）を完成させる。

\subsection{Jacobi 条件の線形化}
$P\in\bigwedge^3V$ と1-形式 $\alpha,\beta\in V^*$ に対し、任意の $\gamma\in V^*$ について $\gamma(P(\alpha,\beta))=P(\alpha,\beta,\gamma)$ となるベクトル $P(\alpha,\beta)\in V$ を定める。$P$ に付随する対称2次形式の線形空間 $K(P)\subset\operatorname{Sym}^2V^*$ を、次の4変数交代恒等式によって定義する：
\begin{align}
K(P)=\{H\in\operatorname{Sym}^2V^*\mid{}
&H(P(\alpha,\beta),P(\gamma,\eta))\\
+H(P(\beta,\gamma),P(\alpha,\eta))\notag\\
&+H(P(\gamma,\alpha),P(\beta,\eta))=0
\quad(\alpha,\beta,\gamma,\eta\in V^*)\}.\label{eq:pg-K}
\end{align}

\begin{lemma}\label{lem:pg-linearize}
$q=\iota_{dF}P$ がPoisson構造であることと、全ての点 $p\in V$ について $F$ のHessian $\operatorname{Hess}_F(p)$ が $K(P)$ に属することは同値である。さらに $X=V(F)$ の滑らかさの仮定の下で、Hessianが張る2次形式の線形系 $\{\operatorname{Hess}_F(p):p\in V\}$ は $\PP(V)$ 上に基点を持たない。
\end{lemma}
\begin{proof}
$P^{ija}=P(dz_i,dz_j,dz_a)$ とおくと $q^{ij}=\sum_a P^{ija}F_a$ である。これをJacobiatorの成分表示\eqref{eq:pg-jacobiator-coefficients}に代入して整理すると、
\begin{equation}
J^{ijk}=\sum_a F_a L_a^{ijk},\qquad
L_a^{ijk}=\sum_{r,b}(P^{ira}P^{jkb}+P^{jra}P^{kib}+P^{kra}P^{ijb})F_{br}
\end{equation}
となる。各 $L_a^{ijk}$ は座標の1次斉次式であるため、補題\ref{lem:pg-gradient}（1次syzygyの消滅）により、$J^{ijk}=0$ は全ての $a,i,j,k$ について $L_a^{ijk}=0$ であることと同値になる。式\eqref{eq:pg-K}の左辺を $\mathcal B_H(\alpha,\beta,\gamma,\eta)$ と書けば、$P$ の交代性から
\begin{equation}
L_a^{ijk}(p)=-\mathcal B_{\operatorname{Hess}_F(p)}(dz_i,dz_j,dz_k,dz_a)
\end{equation}
が成り立つため、多重線形性より第1の主張を得る。

第2の主張（基点自由性）を確認する。3次式の3階偏微分の対称性により、任意の $p,w\in V$ に対して $\operatorname{Hess}_F(p)(w,w)=2\,dF_w(p)$（ここで $dF_w$ は点 $w$ における $dF$ の値）が成り立つ。もし非零ベクトル $w\in V$ が全ての $p\in V$ に対して $\operatorname{Hess}_F(p)(w,w)=0$ を満たすならば、$dF_w=0$ かつEuler等式より $F(w)=0$ となり、$[w]$ が $X$ の特異点となって滑らかさに矛盾する。
\end{proof}

\subsection{基点をもたない2次形式と3重ベクトル}
補題\ref{lem:pg-linearize}により、$q=\iota_{dF}P$ がPoissonならば線形空間 $K(P)$ は $\PP(V)$ 上に基点を持たない。本小節では、基点自由な2次形式の系から非退化形式 $B$ を1つ選び、$P$ と $B$ から $P$ の支持部分空間上にLie代数構造を誘導することで、$P$ が3次元ブロックの直和に分解することを証明する。

まず、基点自由性から非退化形式の存在を導く初等的な線形代数の補題を示す。

\begin{lemma}\label{lem:pg-nondegenerate-pencil}
$W$ を有限次元複素ベクトル空間とし、線形部分空間 $\mathcal L\subset\operatorname{Sym}^2W^*$ が $\PP(W)$ 上に基点を持たない（すなわち任意の非零 $v\in W$ に対し $H(v,v)\ne0$ となる $H\in\mathcal L$ が存在する）とする。このとき $\mathcal L$ は非退化な2次形式を含む。
\end{lemma}
\begin{proof}
$\mathcal L$ の中で階数が最大となる元 $B$ を選び、その階数を $r$ とする。背理法により $B$ が退化している（$r<\dim W$）と仮定し、根基から非零ベクトル $0\ne v\in\operatorname{rad}B$ を取り、$\operatorname{rad}B$ の補空間 $U\subset W$ を選ぶ。このとき制限 $A=B|_U$ は非退化で $\dim U=r$ である。任意の $H\in\mathcal L$ とパラメータ $t\in\CC$ に対し、部分空間 $U\oplus\CC v$ 上の $B+tH$ の表現行列は
\begin{equation}\label{eq:pg-bordered-pencil}
(B+tH)|_{U\oplus\CC v}=
\begin{pmatrix}A+tC&tb\\t b^{\mathsf T}&tH(v,v)\end{pmatrix}
\end{equation}
の形を取る。$B$ の階数最大性により、全ての $t\in\CC$ で $\operatorname{rank}(B+tH)\le r$ であるため、この $(r+1)\times(r+1)$ 行列の行列式は $t$ の多項式として恒等的に零である。その1次の係数は $\det(A)\,H(v,v)$ であり、$\det(A)\ne0$ から $H(v,v)=0$ を得る。これは全ての $H\in\mathcal L$ で成り立つため、$[v]$ が $\mathcal L$ の基点となり仮定に反する。
\end{proof}

次に、式\eqref{eq:pg-K}から誘導されるLie代数上の4変数恒等式が、基点自由性の下でLie代数を $\mathfrak{sl}_2$ の直和へと完全に制限することを示す。これが正規形定理の核心である。

\begin{lemma}[基点自由性によるLie代数の制約]
\label{lem:pg-basepoint-free-lie}
$\mathfrak a$ を完全（$[\mathfrak a,\mathfrak a]=\mathfrak a$）かつ中心が自明（$Z(\mathfrak a)=0$）な有限次元複素Lie代数とする。線形部分空間 $\mathcal L\subset\operatorname{Sym}^2\mathfrak a^*$ は、全ての $H\in\mathcal L$ と $a,b,c,d\in\mathfrak a$ に対して
\begin{equation}
H([a,b],[c,d])+H([b,c],[a,d])+H([c,a],[b,d])=0
\label{eq:pg-lie-H}
\end{equation}
を満たすとする。もし $\mathcal L$ が $\PP(\mathfrak a)$ 上に基点を持たなければ、ある整数 $k\ge0$ に対して $\mathfrak a\simeq\mathfrak{sl}_2^{\oplus k}$ である（$\mathfrak a=0$ のとき $k=0$ とする）。
\end{lemma}
\begin{proof}
証明は2段階で行う。可解根基が非零の場合と、半単純で階数2以上の単純因子を持つ場合のそれぞれについて、式\eqref{eq:pg-lie-H}を巧妙に適用して $\mathcal L$ の共通零点（基点）を明示的に構成し、矛盾を導く。

\paragraph{ステップ1：可解根基の排除（$\mathfrak a$ の半単純性）}
背理法により、可解根基 $\mathfrak r=\operatorname{rad}(\mathfrak a)$ が非零であると仮定する。$\mathfrak r$ の導来列の最後の非零項の中に、極小な非零 $\mathfrak a$-イデアル $I$ を選ぶ。$I$ は可換イデアルであり、随伴表現 $\rho:\mathfrak a\to\operatorname{End}(I)$ に関して既約である。複素数体上の有限次元既約表現では可解根基はスカラーで作用する \cite[Proposition~16.9]{PEtingofLie}。一方、完全性 $\mathfrak a=[\mathfrak a,\mathfrak a]$ により全ての $x\in\mathfrak a$ で $\operatorname{tr}\rho(x)=0$ であるから、$\mathfrak r$ のスカラー作用は零でなければならない。よって $\rho$ は半単純Lie代数 $\mathfrak a/\mathfrak r$ の表現へ降下する。しかもこの作用は非自明である（自明ならば $I\subset Z(\mathfrak a)=0$ となり矛盾する）。

半単純Lie代数 $\mathfrak a/\mathfrak r$ は根に対応する $\mathfrak{sl}_2$ 部分代数で生成されるため、$I$ に非自明に作用する $\mathfrak{sl}_2$ 部分代数を1つ選べる \cite[Lemma~19.16, Theorem~19.19(i)]{PEtingofLie}。有限次元 $\mathfrak{sl}_2$-加群の既約分解 \cite[Theorem~11.16(iii),(iv)]{PEtingofLie} により、この部分代数の標準三元 $\bar h,\bar e,\bar f$ と最高ウェイトベクトル $0\ne v\in I$ を
\begin{equation}\label{eq:pg-primitive-vector}
\bar h v=\ell v,\qquad \bar e v=0,\qquad
\ell\in\mathbb Z_{>0}
\end{equation}
を満たすように取れる。$\bar h,\bar e,\bar f$ の $\mathfrak a$ への任意の持上げを $h,e,f\in\mathfrak a$ とし、$w=\ell^{-1}[f,v]\in I$ とおく。$\mathfrak r$ が $I$ に自明に作用するため商の $\mathfrak{sl}_2$ 関係式を $I$ 上で用いることができ、
\begin{equation}\label{eq:pg-radical-witness}
[h,v]=\ell v,\qquad [e,w]=v,\qquad [e,v]=0,
\qquad [v,w]=0
\end{equation}
が成り立つ（第2式は $e(fv)=[e,f]v+f(ev)=\ell v$ から、最後の式は $I$ の可換性から従う）。ここで式\eqref{eq:pg-lie-H}に $(a,b,c,d)=(h,v,e,w)$ を代入すると、第2項は $[v,e]=0$ により、第3項は $[v,w]=0$ により消滅し、残る第1項から
\[
  \ell\,H(v,v)=H([h,v],[e,w])=0
\]
を得る。$\ell>0$ なので全ての $H\in\mathcal L$ で $H(v,v)=0$ となり、$[v]$ が $\mathcal L$ の基点となって矛盾する。従って $\mathfrak r=0$、すなわち $\mathfrak a$ は半単純である。

\paragraph{ステップ2：階数2以上の単純因子の排除}
次に、半単純Lie代数 $\mathfrak a$ が階数2以上の単純因子 $\mathfrak s$ を持つと仮定する。$\mathfrak s$ の正根系を固定し、その最高根を $\theta$ とおく。$\theta$ は単純根ではあり得ない。実際、もし $\theta$ が単純根なら、階数2以上の既約根系のDynkin図の連結性より $\theta$ に隣接する単純根 $\gamma$ が存在し、根列の性質から $\theta+\gamma$ も根となって最高根性に反する。従って $\theta=\alpha+\beta$ を満たす正根 $\alpha,\beta$ が存在する \cite[Definition~21.13, \S32.2]{PEtingofLie}。

根空間の交換子が非零であること \cite[Theorem~19.19(iv)]{PEtingofLie} より、$a\in\mathfrak s_\alpha$、$b\in\mathfrak s_\beta$ を $v=[a,b]\ne0$ となるように取れる。$\theta$ の最高根性から $\theta+\alpha$ と $\theta+\beta$ は根ではないため、$[a,v]=[b,v]=0$ である。またCartan部分代数から $c=\theta(h)\ne0$ となる元 $h$ を選べば $[h,v]=cv$ である。式\eqref{eq:pg-lie-H}に $(a,b,h,v)$ を代入すると、$[b,h]$ と $[a,v]=0$ および $[h,a]$ と $[b,v]=0$ の組が消え、
\[
  c\,H(v,v)=H([a,b],[h,v])=0
\]
を得る。$c\ne0$ より $[v]$ が再び $\mathcal L$ の基点となって矛盾する。

以上より、$\mathfrak a$ の全ての単純因子の階数は $1$ である。階数1の根系は $\{\alpha,-\alpha\}$ であり各根空間は1次元なので、各単純因子は $\CC h_\alpha\oplus\mathfrak s_\alpha\oplus\mathfrak s_{-\alpha}\simeq\mathfrak{sl}_2$ に同型である \cite[Lemma~19.16, Corollary~19.18]{PEtingofLie}。よって $\mathfrak a\simeq\mathfrak{sl}_2^{\oplus k}$ を得る。
\end{proof}

なお、半単純Lie代数上における恒等式\eqref{eq:pg-lie-H}の解空間は、Derdzinski--Gal \cite[(1.4), Theorem~B]{PDerdzinskiGal} でも別の観点から記述されている。

\begin{lemma}\label{lem:pg-quadratic-lie}
定数3重ベクトル $P\in\bigwedge^3V$ に対し、$K(P)$ が $\PP(V)$ 上に基点を持たなければ、互いに直和をなす3次元部分空間 $W_1,\ldots,W_k\subset V$ と非零元 $P_s\in\bigwedge^3W_s$ が存在して
\[
  P=\sum_{s=1}^k P_s
\]
と分解される（$P=0$ のときは $k=0$ とする）。
\end{lemma}
\begin{proof}
$P\ne0$ とし、$P$ の最小支持部分空間を
\[
  W=\operatorname{span}\{P(\alpha,\beta):\alpha,\beta\in V^*\}\subset V
\]
とおく。$W$ 上で消える任意の共変ベクトル $\eta\in W^\perp$ に対し $P(\alpha,\beta,\eta)=\eta(P(\alpha,\beta))=0$ なので $P\in\bigwedge^3W$ である。従って $P$ の評価は共変ベクトルの $W$ への制限のみに依存し、制限系 $\mathcal L=\{H|_{W\times W}:H\in K(P)\}\subset\operatorname{Sym}^2W^*$ は $P\in\bigwedge^3W$ に対する式\eqref{eq:pg-K}を満たす。$\mathcal L$ は $\PP(W)$ 上に基点を持たないため、補題\ref{lem:pg-nondegenerate-pencil}により非退化な2次形式 $B\in\mathcal L$ が存在する。

$B$ が誘導する同型 $B^\flat:W\xrightarrow{\sim}W^*$（$B^\flat(x)=B(x,-)$）を用いて、$W$ 上に双線形交代括弧を
\begin{equation}\label{eq:pg-support-bracket}
[x,y]=P(B^\flat x,B^\flat y)
\end{equation}
で定める。$B([x,y],z)=P(B^\flat x,B^\flat y,B^\flat z)$ は $x,y,z$ に関して完全交代なので、不変性 $B([x,y],z)=B(x,[y,z])$ が成り立つ。これを用いると任意の $x,y,z,t\in W$ に対し
\begin{equation}\label{eq:pg-support-jacobi}
\begin{aligned}
B([[x,y],z]+[[y,z],x]+[[z,x],y],t)
&=B([x,y],[z,t])+B([y,z],[x,t])\\
&\quad+B([z,x],[y,t])=0
\end{aligned}
\end{equation}
となる（最後の等式は $B\in\mathcal L$ と式\eqref{eq:pg-K}による）。$B$ の非退化性からJacobi恒等式が従い、$W$ はLie代数となる。さらに、$B^\flat$ と制限 $V^*\to W^*$ が全射であるため括弧の値は $W$ 全体を張り $[W,W]=W$ であり、また中心元 $z\in Z(W)$ については $B(z,[x,y])=B([z,x],y)=0$ かつ $[W,W]=W$ より $B$ の非退化性から $Z(W)=0$ である。

任意の $H\in\mathcal L$ に対して式\eqref{eq:pg-K}に $B^\flat a,B^\flat b,B^\flat c,B^\flat d$ を代入すれば式\eqref{eq:pg-lie-H}を得るため、補題\ref{lem:pg-basepoint-free-lie}を適用して直和分解 $W=W_1\oplus\cdots\oplus W_k$（各 $W_s\simeq\mathfrak{sl}_2$、従って $\dim W_s=3$）を得る。異なる単純因子 $W_s,W_t$（$s\ne t$）は $B$ に関して直交する。実際、任意の $x=[a,b]\in W_s$ と $y\in W_t$ について $B(x,y)=B(a,[b,y])=0$ であり、$[W_s,W_s]=W_s$ だからである。従って各制限 $B_s=B|_{W_s}$ は非退化となり、Cartan 3-形式 $C(x,y,z)=B([x,y],z)$ は各因子上の3-形式 $C_s$ の直和に分解する。同型 $(B^\flat)^{-1}$ で3つの添字を上げれば $P=\sum_{s=1}^k P_s$（$P_s\in\bigwedge^3W_s$）を得る。各因子で括弧は非零かつ $B_s$ は非退化なので $C_s\ne0$、ゆえに $P_s\ne0$ である。
\end{proof}

\subsection{3次式の分解と一般階数}
正規形定理の証明において、アフィン錐から射影空間・超曲面へ降りた際のPoisson階数の変化を正確に捉えるため、次の初等的な線形代数の補題を用意する。

\begin{lemma}[直線による商と双ベクトルの階数]\label{lem:pg-projective-rank}
$T$ を有限次元ベクトル空間、$b\in\bigwedge^2T$ を階数 $2r$ の双ベクトル、$L\subset T$ を1次元部分空間とする。$I_b=\operatorname{im}(b^\sharp:T^*\to T)$ とおき、商空間 $T/L$ における $b$ の像を $\overline b$ とすると、
\begin{equation}
\operatorname{rank}\overline b=
\begin{cases}
2r-2,&L\subset I_b,\\
2r,&L\not\subset I_b
\end{cases}
\end{equation}
が成り立つ。
\end{lemma}
\begin{proof}
$L\not\subset I_b$ ならば射影 $I_b\to T/L$ は単射であるため、双ベクトルの階数は不変である。逆に $L\subset I_b$ ならば、$I_b$ のDarboux基底 $e_1,f_1,\ldots,e_r,f_r$ を $L=\CC e_1$ かつ $b=e_1\wedge f_1+\cdots+e_r\wedge f_r$ となるように取れる。商空間 $T/L$ へ移ると最初の項 $e_1\wedge f_1$ のみが消え、残る $2r-2$ 本の基底ベクトルは1次独立であるため階数は $2r-2$ となる。
\end{proof}

以上の準備がすべて揃ったので、本稿の主定理であるPoisson正規形を証明する。

\begin{theorem}[Poisson 正規形]\label{thm:pg-normal}\label{pthm:normal-form}
滑らかな3次超曲面 $X=V(F)\subset\PP(V)$（$\dim V=N=n+2\ge5$）上の任意の非零Poisson構造 $\pi_X$ に対して、$N=3k+m$ を満たす整数 $k\ge1$、$m\ge0$ と、$V^*$ の適切な線形座標 $x_s=(x_{s1},x_{s2},x_{s3})$（$1\le s\le k$）および $u=(u_1,\ldots,u_m)$ が存在して、
\begin{equation}
P=\sum_{s=1}^k\partial_{x_{s1}}\wedge\partial_{x_{s2}}\wedge\partial_{x_{s3}},
\qquad F=\sum_{s=1}^k f_s(x_s,u)\label{eq:pg-normal}
\end{equation}
となる。このときアフィン錐上のPoisson括弧 $q=\iota_{dF}P$ は、各ブロック内の3変数Jacobian括弧
\begin{equation}
\{x_{s1},x_{s2}\}=(f_s)_{x_{s3}},\quad
\{x_{s2},x_{s3}\}=(f_s)_{x_{s1}},\quad
\{x_{s3},x_{s1}\}=(f_s)_{x_{s2}}\label{eq:pg-bracket}
\end{equation}
および異なるブロック間と中心変数の括弧 $\{x_{si},x_{tj}\}=0$（$s\ne t$）、$\{u_\alpha,-\}=0$ によって定まる。

逆に、式\eqref{eq:pg-normal}の形を持つ任意の滑らかな3次式 $F$ は、この括弧によって $\PP(V)$ および $X$ 上に非零Poisson構造 $\pi,\pi_X$ を定める。さらに、$\pi$ と $\pi_X$ の一般階数はいずれも $2k$ である。
\end{theorem}
\begin{proof}
\paragraph{ステップ1：正規形\eqref{eq:pg-normal}の導出}
命題\ref{prop:pg-represent}および命題\ref{prop:pg-jacobi-extension}により、非零Poisson構造 $\pi_X$ は非零定数3重ベクトル $P\in\bigwedge^3V$ と錐上のPoisson双ベクトル場 $q=\iota_{dF}P$ を一意に定める。補題\ref{lem:pg-linearize}および補題\ref{lem:pg-quadratic-lie}を適用して、直和分解 $P=\sum_{s=1}^k P_s$（$0\ne P_s\in\bigwedge^3W_s$、$\dim W_s=3$）を得る。各 $W_s$ の基底を1本スカラー倍して $P_s$ をその3本の外積に規格化し、さらに補空間 $U$（$V=(\bigoplus_{s=1}^k W_s)\oplus U$、$\dim U=m$）の基底を選んで双対座標を $x_s,u$ と書けば、式\eqref{eq:pg-normal}の $P$ の表示を得る。

次に、3次式 $F$ もこのブロック構造に沿って分離することを示す。補題\ref{lem:pg-quadratic-lie}で構成したLie代数構造について、異なる単純因子から $a,b\in W_s$、$c,d\in W_t$（$s\ne t$）を選んで式\eqref{eq:pg-lie-H}に代入すると、$[b,c]=[c,a]=0$ より $H([a,b],[c,d])=0$ を得る。各因子は完全（$[W_s,W_s]=W_s$）なので、任意の $H\in K(P)$ に対して $H(W_s,W_t)=0$（$s\ne t$）が成り立つ。特に全ての点における $F$ のHessianについてこの直交性が成り立つため、異なるブロック間の混合2階偏微分 $\partial_{x_{si}}\partial_{x_{tj}}F$ は恒等的に零である。標数0なので、異なる2つの $x$-ブロックの変数を同時に含む単項式は $F$ に現れない。従って各単項式をそれが含む $x$-ブロックへ振り分け、中心変数 $u$ のみからなる単項式をいずれかのブロックへまとめれば、$F=\sum_{s=1}^k f_s(x_s,u)$ を得る。

\paragraph{ステップ2：逆向きの確認}
逆に式\eqref{eq:pg-normal}を仮定すると、縮約 $q=\iota_{dF}P$ は式\eqref{eq:pg-bracket}の括弧を与え、異なるブロック間および $u_\alpha$ との括弧はすべて零となる。単一ブロック内で $f=f_s$、$f_i=\partial_{x_{si}}f$ と略記すれば、3座標 $(x_{s1},x_{s2},x_{s3})$ のJacobiatorは
\begin{equation}\label{eq:pg-block-jacobi-calculation}
\{x_{s1},f_1\}+\{x_{s2},f_2\}+\{x_{s3},f_3\}
=f_3(f_{12}-f_{21})+f_2(f_{31}-f_{13})
 +f_1(f_{23}-f_{32})=0
\end{equation}
である。異なるブロックや中心変数を含む組のJacobiatorも、各ブロックの係数が他ブロックの変数に依存せず $u$ がCasimirであることから直ちに零となる。Jacobiatorは各引数について導分なので、座標生成元上での消滅から多項式環全体でのJacobi恒等式 $[q,q]=0$ が従う。さらに $\iota_{dF}q=0$ かつ2次斉次なので $\PP(V)$ および $X$ へ降下し、$P\ne0$ と命題\ref{prop:pg-represent}の単射性から $X$ 上で非零なPoisson構造を得る。

\paragraph{ステップ3：$\pi$ と $\pi_X$ の一般階数の計算}
各ブロックの勾配 $\nabla_{x_s}f_s$ は恒等的に零ではない（もし零なら $F$ は $W_s$ の変数に依存せず、射影平面 $\PP(W_s)$ が丸ごと $X$ の特異点集合に含まれて滑らかさに反する）。従ってアフィン空間 $V$ の一般点では各3変数ブロックが階数 $2$ を持ち、直和である $q$ の階数は $2k$ となる。

$\PP(V)$ 上の一般点（$F\ne0$）では、$dF(E)=3F\ne0$ である一方 $\operatorname{im}q^\sharp\subset\ker dF$ なので、Euler方向 $E$ は $\operatorname{im}q^\sharp$ に属さない。ゆえに補題\ref{lem:pg-projective-rank}により射影化で階数は下がらず、$\pi$ の一般階数は $2k$ である。

最後に $X$ 上の一般点を調べる。非零の2次式である各 $x_s$-偏微分は既約な3次式 $F$ で割り切れないため、$X$ の一般点でも各ブロックは階数 $2$ を持ち、錐上の階数は $2k$ である。ここで $E\notin\operatorname{im}q^\sharp$ が $X$ の一般点で成り立つことを $m>0$ と $m=0$ に分けて確認する：
\begin{itemize}
  \item $m>0$ のとき：$X$ の一般点では $u\ne0$ なので $E$ は中心方向 $U$ に非零成分を持つが、$\operatorname{im}q^\sharp\subset\bigoplus_{s=1}^k W_s$ なので $E\notin\operatorname{im}q^\sharp$ である。
  \item $m=0$ のとき：$N=3k\ge5$ より $k\ge2$ である。各ブロックの像は $\ker df_s$ なので、$E=\sum_s E_s$ が $\operatorname{im}q^\sharp$ に属する条件は全ての $s$ について $df_s(E_s)=3f_s=0$ となることである。各 $f_s$ は非零で互いに異なる変数に依存するため $F$ に比例せず、共通零点 $\{f_1=\cdots=f_k=0\}$ は $X$ の真の閉部分集合である。
\end{itemize}
いずれの場合も $X$ の一般点で $E\notin\operatorname{im}q^\sharp$ が成り立ち、補題\ref{lem:pg-projective-rank}より $\pi_X$ の一般階数も $2k$ である。
\end{proof}

\subsection{ブロックの内在性と一般の3次超曲面}
正規形\eqref{eq:pg-normal}に現れる3次元部分空間の組 $\{W_1,\ldots,W_k\}$ は座標の選び方によらない内在的な不変量である。なお、Cartan 3-形式からの単純因子の復元はDerdzinski--Gal \cite[Lemma~7.1(a)]{PDerdzinskiGal} でも述べられているが、次の補題は縮約双ベクトルの階数条件を用いてこれをスキーム論的に復元する。

\begin{lemma}[縮約双ベクトルによるブロックの復元]\label{ne:blocks}
$W=\bigoplus_{s=1}^k W_s$（各 $\dim W_s=3$）とし、$P=\sum_{s=1}^k P_s$（$0\ne P_s\in\Lambda^3 W_s$）とする。このとき $W^*$ 内の閉部分スキーム
\begin{equation}
D_P=\{\alpha\in W^*:(\iota_\alpha P)\wedge(\iota_\alpha P)=0\}
\end{equation}
は、$k$ 個の3次元線形部分空間 $W_1^*,\ldots,W_k^*$ の被約な和に一致する。従って $P$ は順序を除いてブロックの組 $\{W_1,\ldots,W_k\}$ を一意に決定する。
\end{lemma}
\begin{proof}
$\alpha=(\alpha_1,\ldots,\alpha_k)\in\bigoplus_s W_s^*$ に対して $\beta_s=\iota_{\alpha_s}P_s\in\Lambda^2 W_s$ とおく。$\dim W_s=3$ なので $\beta_s\wedge\beta_s=0$ であり、異なる $s\ne t$ の積 $\beta_s\wedge\beta_t$ は互いに独立な直和成分 $\Lambda^2 W_s\otimes\Lambda^2 W_t$ に属する。各 $W_s^*\to\Lambda^2 W_s$（$\alpha_s\mapsto\beta_s$）は線形同型であるため、$D_P$ の定義イデアルは異なるブロックの座標の積全体で生成される。これは各座標部分空間 $W_s^*$ の定義イデアルの交叉にほかならず、特に根基イデアルである。よって $D_P$ の既約成分として $\{W_1^*,\ldots,W_k^*\}$ が復元され、他の成分の和の零化空間として各 $W_s$ も一意に復元される。
\end{proof}

正規形\eqref{eq:pg-normal}では、この内在的な組 $\{W_1,\ldots,W_k\}$ に対して補空間 $U$ と各空間の基底を選んでいる。また座標を固定した後でも、$\sum_{s=1}^k h_s(u)=0$ を満たす中心変数のみの3次式 $h_s(u)$ を各 $f_s$ に加える自由度が残る。

正規形定理の直截な応用として、一般の滑らかな3次超曲面では一般階数2のPoisson構造しか存在しないことが分かる。

\begin{corollary}\label{cor:pg-generic}
固定した滑らかな3次超曲面 $X$ 上で一般階数 $2$ のPoisson構造が張る被約な射影パラメータ空間は、Pl\"ucker埋込みされたGrassmann多様体 $\operatorname{Gr}(3,V)\subset\PP(\bigwedge^3V)$ に一致する。また、一般の滑らかな3次超曲面においては、全ての非零Poisson構造が一般階数 $2$ を持つ。
\end{corollary}
\begin{proof}
命題\ref{prop:pg-represent}の線形同型の下で、一般階数 $2$（すなわち $k=1$）であることは非零3重ベクトル $P$ が分解可能であることと同値なので、対応する射影パラメータ空間は $\operatorname{Gr}(3,V)$ に一致する。

第2の主張を示すため、各 $k\ge2$ に対して一般階数 $2k$ のPoisson構造を許容する3次式の族の余次元を評価する。互いに直和をなす順序付き3次元部分空間の組 $(W_1,\ldots,W_k)$ の空間は $\operatorname{Gr}(3,V)^k$ の非空開集合であり、その次元は $3k(N-3)$ である。この組を固定したとき、許される3次式 $F$ は全ての $v\in V$ に対して $\operatorname{Hess}F(v)(W_s,W_t)=0$（$s\ne t$）を満たす線形部分空間に属する。この条件は $W_s$ のみに依存するため、次元を数えるために任意の補空間 $U$（$\dim U=m=N-3k$）を選んでよい。その座標で許される単項式を数えると、この線形空間の次元は
\begin{equation}
\binom{m+2}{3}+k\left(10+6m+3\binom{m+1}{2}\right)
\end{equation}
である。従って3次式全体の空間 $\operatorname{Sym}^3V^*$ におけるこの族の余次元は少なくとも
\begin{equation}
\binom{N+2}{3}-\binom{m+2}{3}
-k\left(10+6m+3\binom{m+1}{2}\right)-3k(N-3)
=\frac{3k}{2}\bigl((3k-5)m+3k(k-1)\bigr)>0\label{eq:pg-codim}
\end{equation}
となる（$k\ge2$、$m\ge0$ より右辺は真に正である）。可能な $k\le\lfloor N/3\rfloor$ は有限個しかないので、一般の滑らかな3次式はこれらすべての族を避ける。
\end{proof}


\section{Casimir、対数形式および symplectic leaf}\label{sec:pg-leaves}

以下では定理\ref{thm:pg-normal}の正規形\eqref{eq:pg-normal}を固定し、アフィン錐上の双ベクトル場 $q=\iota_{dF}P$、その $\PP(V)$ への射影化 $\pi$、および $X$ への制限 $\pi_X$ の幾何学的構造を詳しく調べる。多項式 $G\in S$ が\emph{Casimir}であるとは $\{G,S\}_q=0$ を満たすことをいい、射影超曲面が\emph{Poisson超曲面}であるとはその定義イデアル層がPoisson括弧で保たれることをいう。

\subsection{Poisson超曲面と1次Casimir}
\begin{proposition}\label{prop:pg-casimir}
既約な斉次多項式 $G\in S$ に対し、超曲面 $V(G)\subset\PP(V)$ がPoisson超曲面であることと、$G$ がアフィン錐上のCasimirであることは同値である。特に、1次Casimir全体の空間は中心変数が張る空間 $\operatorname{span}(u_1,\ldots,u_m)$ に一致し、Poisson超平面のパラメータ空間は $m>0$ ならば $\PP^{m-1}$、$m=0$ ならば空である。
\end{proposition}
\begin{proof}
$G$ がCasimirならばイデアル $(G)$ は括弧で保たれるため、$V(G)$ はPoisson超曲面である。

逆に $V(G)$ がPoisson超曲面であるとする。$G$ が $F$ の定数倍のときは既知なので、$\gcd(F,G)=1$ と仮定する。$d=\deg G$、$Y=V(G)$ とおく。$V(G)$ がPoissonであるという射影的接性条件により、$\iota_{dG}q|_Y$ は制限したEuler列の核 $\OO_Y(d)E$ に属する。$H^1(\PP(V),\OO)=0$ より $H^0(Y,\OO_Y(d))$ の切断は $S_d$ へ持ち上がるため、ある $A\in S_1\otimes V$ と $h\in S_d$ を用いて
\begin{equation}
\iota_{dG}q=GA+hE, \qquad A\in S_1\otimes V, \quad h\in S_d
\end{equation}
と書ける。この両辺を $dG$ で縮約すると $0=GA(G)+dhG$ より $h=-A(G)/d$ を得る。次に $dF$ で縮約し $\iota_{dF}q=0$ を用いると
\[
  GA(F)-\frac{3}{d}FA(G)=0
\]
が成り立つ。$F$ と $G$ は互いに素であり、$\deg A(F)=3$、$\deg A(G)=d$ なので、ある定数 $c\in\CC$ に対して $A(F)=cF$ かつ $A(G)=cdG/3$ となる。従って $A-(c/3)E$ は勾配 $dF$ の1次syzygyを与え、補題\ref{lem:pg-gradient}により $A=(c/3)E$ である。このとき $h=-A(G)/d=-(c/3)G$ となるため、$\iota_{dG}q=GA+hE=0$、すなわち $G$ はCasimirである。

最後に1次式 $\ell\in S_1$ について考える。$\iota_{d\ell}q=\iota_{d\ell}\iota_{dF}P=-\iota_{dF}\iota_{d\ell}P$ であり、$\iota_{d\ell}P$ は定数双ベクトルなので、補題\ref{lem:pg-gradient}（定数係数syzygyの消滅）により $\iota_{d\ell}q=0\iff\iota_{d\ell}P=0$ である。正規形\eqref{eq:pg-normal}の下ではこれは $\ell$ がどのブロック変数 $x_{si}$ にも依存しないこと、すなわち $\ell\in\operatorname{span}(u_1,\ldots,u_m)$ と同値である。
\end{proof}

正規形\eqref{eq:pg-normal}からは、各中心変数 $u_\alpha$、各ブロックの3次式 $f_s$、および $F=\sum_s f_s$ とそれらの多項式結合がすべてCasimirであることが直ちに従う。

\subsection{最高外積の対数微分形式表示}
Poisson構造の一般階数が $2k$ であることから、最高次の非零外積 $\pi^k$ および $\pi_X^k$ は、余次元がそれぞれ $m+k-1$ および $n-2k=m+k-2$ の微分形式と対応する。正規形を用いると、これらがCasimirの対数微分で極めて簡潔に表される。

\begin{proposition}\label{prop:pg-log}
$V$ の体積形式を1つ固定すると、$\PP(V)$ 上の $\pi^k$ に対応する微分形式は非零定数倍を除いて
\begin{equation}
\beta=\iota_E(du_1\wedge\cdots\wedge du_m\wedge df_1\wedge\cdots\wedge df_k)
\in H^0(\PP^{N-1},\Omega^{m+k-1}(N))\label{eq:pg-log-main}
\end{equation}
である。特に、分母が消えない開集合上では次のように対数形式で書ける：
\begin{align}
m>0:\quad &\beta=\left(\prod_{\alpha=1}^m u_\alpha\prod_{s=1}^kf_s\right)
\bigwedge_{\alpha=2}^m d\log(u_\alpha/u_1)
\wedge\bigwedge_{s=1}^k d\log(f_s/u_1^3),\\
m=0:\quad &\beta=3\left(\prod_{s=1}^kf_s\right)
\bigwedge_{s=2}^k d\log(f_s/f_1).
\end{align}
さらに、$X$ 上の $\pi_X^k$ に対応する微分形式は
\begin{equation}
\alpha=\iota_E(du_1\wedge\cdots\wedge du_m\wedge df_1\wedge\cdots\wedge df_{k-1})|_X
\in H^0(X,\Omega_X^{m+k-2}(N-3))\label{eq:pg-log-X}
\end{equation}
で与えられる。
\end{proposition}
\begin{proof}
第 $s$ ブロックの双ベクトル場を $q_s$ と書く。各ブロックは3次元空間上の双ベクトルなので $q_s^2=0$ であり、従って $q^k=k!\,q_1\wedge\cdots\wedge q_k$ である。これをアフィン体積形式へ縮約すると、各 $q_s=\iota_{df_s}(\partial_{x_{s1}}\wedge\partial_{x_{s2}}\wedge\partial_{x_{s3}})$ の寄与により非零定数倍を除いて $du_1\wedge\cdots\wedge du_m\wedge df_1\wedge\cdots\wedge df_k$ を得る（ここで $df_s$ の中心方向 $du_\alpha$ 成分は先頭の $du_1\wedge\cdots\wedge du_m$ との外積で自動的に消える）。さらにEuler場 $E$ で縮約すれば射影空間上の水平形式が得られ、全スケーリング重み $m+3k=N$ がtwistを与える。

対数表示は、$m>0$ では $\iota_E d\log u_1=1$ および $\iota_E d\log(u_\alpha/u_1)=\iota_E d\log(f_s/u_1^3)=0$ を、$m=0$ では $\iota_E d\log f_1=3$ および $\iota_E d\log(f_s/f_1)=0$ を用いて外積を展開すれば得られる。計算は分母が非零の開集合で行うが、左辺 $\beta$ は大域的に正則であるため大域表示として意味を持つ。

最後に $X$ 上への制限を考える。$F=\sum_{s=1}^k f_s$ なので外積の中で $df_k$ を $dF$ に置き換えても形式は変わらない。$X$ 上で法線方向 $dF$ を除いて随伴公式を適用すれば式\eqref{eq:pg-log-X}を得る（局所的に $dF$ が座標微分となる表示で縮約すればよく、法線束 $\OO_X(3)$ によりtwistは $N-3$、形式次数は $m+k-2=n-2k$ となる）。
\end{proof}

なお、上の対数表示に現れる因子の次数の総和は $m+3k=N$ であり、$V(u_1\cdots u_m f_1\cdots f_{k-1}F)$ は $\PP(V)$ の反標準因子を定める。

\subsection{第一積分とsymplectic leafの閉包}
Casimir $u_\alpha,f_s$ はそのまま $\pi$ および $\pi_X$ の第一積分を与え、一般のsymplectic leafの幾何学を完全に記述する。

\begin{proposition}\label{prop:pg-first-integrals}
重み付き射影空間への有理写像
\begin{equation}
\Phi:\PP^{N-1}\dashrightarrow\PP(1^m,3^k),\qquad
[x,u]\longmapsto[u_1:\cdots:u_m:f_1:\cdots:f_k]
\end{equation}
は一般点で1次独立な第一積分を与え、一般ファイバーの滑らかな最大階数部分の連結成分が $\pi$ のsymplectic leafとなる。$X=V(F)$ へ制限したときの底空間は、$m>0$ ならば $\PP(1^m,3^{k-1})$、$m=0$ ならば $\PP^{k-2}$ であり、いずれの場合も底の次元は $n-2k$ である。
\end{proposition}
\begin{proof}
全てのブロックの勾配 $\nabla_{x_s}f_s$ が非零となる開集合上では、$du_\alpha$ を法とした $df_s$ たちが互いに独立なブロック空間に非零成分を持つため、$(u_1,\ldots,u_m,f_1,\ldots,f_k)$ の微分の階数は $m+k$ である。$X$ への制限（条件 $\sum_s f_s=0$）と射影化（スケーリングの1方向）がそれぞれ独立な条件を1つずつ除くため、一般ファイバーの次元は $2k$ である。

各 $u_\alpha,f_s$ はCasimirであるため、Hamiltonベクトル場が張る分布は $\ker d\Phi$ に含まれるが、定理\ref{thm:pg-normal}により両者の次元はいずれも $2k$ なので一致する。最大階数部分での逆2-形式はJacobi恒等式により閉形式であり、よって滑らかなファイバーの連結成分がsymplectic leafをなす。また $\Phi$ の不定点集合は $\{u=0,\ f_1=\cdots=f_k=0\}$ であるが、その上では、ブロック勾配が1つでも零ならアフィン階数が高々 $2k-2$、全て非零ならEuler方向が像に入るため補題\ref{lem:pg-projective-rank}より射影階数が $2k-2$ となる。ゆえに最大階数 $2k$ のsymplectic leafは不定点集合と交わらない。
\end{proof}

各ブロックの3変数制限 $f_s(x_s,0)$ が滑らかな平面3次曲線を定めるとき、一般のsymplectic leafのザリスキー閉包は美しい完全交叉となる。

\begin{proposition}\label{prop:pg-leaf-closure}
各 $f_{s,0}(x_s):=f_s(x_s,0)$（$1\le s\le k$）が滑らかな平面3次曲線を定めると仮定する。このとき、$\PP(V)$ 上の $\pi$ と $X$ 上の $\pi_X$ の双方について、一般のsymplectic leafの閉包は
\begin{itemize}
  \item $m>0$ のとき：$\PP^{3k}$ 内の型 $(3^k)$ の既約な完全交叉、
  \item $m=0$ のとき：$\PP^{3k-1}$ 内の型 $(3^{k-1})$ の既約な完全交叉
\end{itemize}
であり、いずれの場合も次元は $2k$ である。
\end{proposition}
\begin{proof}
$m>0$ の場合と $m=0$ の場合に分けて証明する。

\paragraph{場合1：$m>0$}
滑らかな開集合 $U_{\PP}=\PP(V)\cap\{u_1\ne0\}$ と $U_X=X\cap\{u_1\ne0\}$ 上で、第一積分が定める射 $\psi_{\PP}:U_{\PP}\to\mathbb A^{m-1}\times\mathbb A^k$ および $\psi_X:U_X\to\mathbb A^{m-1}\times\mathbb A^{k-1}$ を
\begin{equation}\label{eq:pg-affine-first-integrals}
\begin{aligned}
\psi_{\PP}([x,u])&=
\left(\frac{u_2}{u_1},\ldots,\frac{u_m}{u_1},
  \frac{f_1}{u_1^3},\ldots,\frac{f_k}{u_1^3}\right),\\
\psi_X([x,u])&=
\left(\frac{u_2}{u_1},\ldots,\frac{u_m}{u_1},
  \frac{f_1}{u_1^3},\ldots,\frac{f_{k-1}}{u_1^3}\right)
\end{aligned}
\end{equation}
とおく。命題\ref{prop:pg-first-integrals}の微分の独立性から両射は支配的であり、標数0における一般的滑らかさの定理により、一般の底点に対するファイバーは滑らかである。一般の底点の値を $c_\alpha=u_\alpha/u_1$（$c_1=1$）、$d_s=f_s/u_1^3$ と書く（$\psi_X$ では $d_k=-\sum_{s<k}d_s$ であり、特に $k=1$ では後半の座標がなく $d_1=0$ である）。

$c=(c_1,\ldots,c_m)$、$t=u_1$ とおき、部分射影空間 $\{u=ct\}\simeq\PP^{3k}$ の中で $k$ 本の3次方程式
\begin{equation}\label{eq:pg-leaf-closure-equations}
  f_s(x_s,ct)-d_st^3=0\qquad(1\leq s\leq k)
\end{equation}
が定める閉部分スキーム $Y$ を考える。無限遠超平面 $t=0$ との交わりは独立な式 $f_{s,0}(x_s)=0$ で定まるため次元 $2k-1$ であり、一方 $t\ne0$ のアフィン部分（$t=1$ と規格化）は $k$ 個のアフィン3次曲面 $\{f_s(x_s,c)=d_s\}\subset\mathbb A^3$ の直積である。従って $Y$ は純次元 $2k$ の型 $(3^k)$ 完全交叉であり、無限遠 $t=0$ に含まれる成分を持たない。一般の底点に対するアフィン部分の滑らかさから各アフィン因子も滑らかであり、さらに各因子の射影閉包 $\{f_s(x_s,ct)=d_st^3\}\subset\PP^3$ は無限遠 $t=0$ でも仮定 $\nabla f_{s,0}\ne0$ により滑らか（ゆえに連結かつ既約な3次曲面）である。従って各アフィン因子およびその直積は既約である。完全交叉 $Y$ はCohen--Macaulayで埋込み成分を持たず、全ての最小成分が既約かつ滑らかな開集合 $Y\cap\{t\ne0\}$ と交わるため、$Y$ は被約・既約であり、求める葉の閉包に一致する。

\paragraph{場合2：$m=0$}
$k\ge2$ である。$f_1\ne0$ の一般底点に対して $c_1=1$、$f_s=c_s f_1$（$s\ge2$）と書く（$X$ 上では $\sum_{s=1}^k c_s=0$ を課すが、$k\ge2$ より全ての $c_s$ が非零となる開集合は空でない）。これら $k-1$ 本の3次方程式が $\PP^{3k-1}$ 内に定める閉部分スキームを考える。$\{f_1=0\}$ の部分は全 $f_s=0$ で定まり次元 $2k-1$ である。一方 $\{f_1\ne0\}$ の開部分は、アフィン超曲面の直積 $\{f_1=1\}\times\prod_{s=2}^k\{f_s=c_s\}$ を $\mu_3$ の同時スカラー作用で割った像である。$c_s\ne0$ の一般点では各因子が滑らかな3次曲面 $\{f_s(x_s)=c_s t^3\}$ のアフィン開集合なので既約であり、全てのブロック勾配が非零なので方程式の微分は1次独立でファイバーは滑らかである。場合1と同じ純次元性・Cohen--Macaulay性により、この型 $(3^{k-1})$ の完全交叉が被約・既約な葉の閉包を与える。
\end{proof}

\begin{remark}\label{rem:pg-zero-restriction-leaf}
3変数制限 $f_{s,0}$ の滑らかさを仮定しなくても、命題\ref{prop:pg-first-integrals}はそのまま成り立つ。$m>0$ の場合に葉の閉包を得るには、方程式\eqref{eq:pg-leaf-closure-equations}のイデアルを $t$ に関して飽和させ、必要に応じてファイバーの既約成分を取ればよい。例えば $k=1$ かつ $f_{1,0}\equiv0$ のとき、$F(x,ct)=tQ_c(x,t)$ と分解する。ここで $Q_\alpha(x)=F_{u_\alpha}(x,0)$ は $\PP_x^2$ 上に共通零点を持たない（共通零点があれば $X$ の特異点になるため）。Bertiniの定理により一般の $\sum_\alpha c_\alpha Q_\alpha$ は滑らかな平面2次曲線を定めるため、$Q_c(x,t)$ の階数は少なくとも3で $t$ で割り切れず既約である。従って葉の閉包は、固定平面 $t=0$ を除いた残余2次曲面 $Q_c=0$ となる。
\end{remark}


\section{全退化スキームの方程式}\label{sec:pg-degeneracy}

本節では、正規形\eqref{eq:pg-normal}を用いて、Poisson構造の階数が $2j-2$ 以下となる全ての退化スキーム $D_{2j-2}(\pi)=Z(\pi^j)$（$\PP(V)$ 上）および $D_{2j-2}(\pi_X)$（$X$ 上）の定義イデアルを明示的に求める。さらに、各3変数制限 $f_s(x_s,0)$ が滑らかな平面3次曲線 $E_s$ を定める場合には、退化集合が楕円曲線のjoinと臨界点の有限平坦被覆に分解し、特に零点スキーム $D_0(\pi_X)$ が $\coprod_s E_s$ と完全交叉 $Z$ の非交和となることを示す。

\subsection{任意のブロック表示に対する定義イデアルの生成元}
各ブロック $s\in\{1,\ldots,k\}$ について、$q_s=\iota_{df_s}(\partial_{x_{s1}}\wedge\partial_{x_{s2}}\wedge\partial_{x_{s3}})$ とし、ブロック方向の偏微分およびそのブロックEuler和を
\begin{equation}
A_{s,i}=\frac{\partial f_s}{\partial x_{si}},\qquad
p_s=\sum_{i=1}^3 x_{si}A_{s,i}=3f_s-\sum_{\alpha=1}^m u_\alpha\frac{\partial f_s}{\partial u_\alpha}
\end{equation}
とおく。以下の命題は、3変数制限 $f_s(x_s,0)$ の滑らかさを一切仮定せずに成り立つ。

\begin{proposition}\label{prop:pg-deg-ideal}
$1\leq j\leq k$ に対し、$\PP(V)$ 上の退化スキーム $D_{2j-2}(\pi)=Z(\pi^j)$ のイデアル層は、大きさ $j$ の全ての部分集合 $I\subset\{1,\ldots,k\}$ に対する次の3種類の斉次多項式
\begin{equation}
u_\alpha\prod_{s\in I}A_{s,i_s},\qquad
p_t\prod_{s\in I\setminus\{t\}}A_{s,i_s}\ (t\in I),\qquad
x_{t,l}\prod_{s\in I}A_{s,i_s}\ (t\notin I)\label{eq:pg-deg-generators}
\end{equation}
（各添字 $\alpha,i_s,l$ は可能な範囲をすべて動く）によって生成される。また、$X$ 上の退化スキーム $D_{2j-2}(\pi_X)$ のイデアル層は、これらに $F$ を加えて得られる。各生成元の斉次次数はいずれも $2j+1$ である。
\end{proposition}
\begin{proof}
アフィン錐の原点以外では、Euler場 $E$ が張る直線を局所的な直和因子に選べるため、商束の外積からの写像 $E\wedge-$ は局所分裂単射となる。従って射影化された $\pi^j$ の零点イデアルは、アフィン錐上の多重ベクトル場 $E\wedge q^j$ の係数イデアルに等しい。

$q_s^2=0$ より $q^j=j!\sum_{|I|=j}\bigwedge_{s\in I}q_s$ と展開される。ここにEuler場 $E=\sum_\alpha u_\alpha\partial_{u_\alpha}+\sum_{t,l}x_{tl}\partial_{x_{tl}}$ を外積すると、各項は次の3つの型に分かれる：
\begin{itemize}
  \item $E$ の中心方向成分 $\sum_\alpha u_\alpha\partial_{u_\alpha}$ は第1の型 $u_\alpha\prod_{s\in I}A_{s,i_s}$ を与える。
  \item 選択したブロック $t\in I$ の成分は、恒等式 $(\sum_l x_{tl}\partial_{x_{tl}})\wedge q_t=p_t\,\partial_{x_{t1}}\wedge\partial_{x_{t2}}\wedge\partial_{x_{t3}}$ により第2の型 $p_t\prod_{s\in I\setminus\{t\}}A_{s,i_s}$ を与える。
  \item 選択しなかったブロック $t\notin I$ の成分は第3の型 $x_{t,l}\prod_{s\in I}A_{s,i_s}$ を与える。
\end{itemize}
これら3つの型は、各空間（中心空間、選択ブロック、非選択ブロック）における外積次数がそれぞれ $(1,2,\ldots,2,0)$、$(0,3,2,\ldots,2,0)$、$(0,2,\ldots,2,1)$ と互いに異なるため、異なる型の間で相殺は起こらない。最後に $T_X\subset T_{\PP(V)}|_X$ は局所的に直和因子として分裂するため、外積束への包含は切断の零点イデアルを変えず、$X$ 上のイデアルは $F$ を加えるだけで得られる。
\end{proof}

\subsection{臨界点の有限平坦被覆}
本小節および続く2小節では、各ブロックの3変数制限について次の滑らかさの条件を仮定する：
\begin{equation}
f_{s,0}=f_s(x_s,0)\text{ が滑らかな平面3次曲線 }E_s\text{ を定める}\quad(1\le s\le k).\label{eq:pg-smooth-block}
\end{equation}
各楕円曲線 $E_s$ は $\{u=0,\ x_t=0\ (t\neq s)\}$ によって $X$ に埋め込まれており、互いに交わらない。

\begin{lemma}\label{lem:pg-critical-cover}
条件\eqref{eq:pg-smooth-block}の下で、$m>0$ ならば各ブロックの臨界点スキームおよび共通臨界点スキーム
\begin{equation}
K_s=V(A_{s,1},A_{s,2},A_{s,3})\subset\PP(x_s,u)=\PP^{m+2},\qquad
\Delta=V(A_{s,i}\mid s,i)\subset\PP(V)
\end{equation}
は超平面族 $\{u=0\}$ と交わらない。中心座標への射影は $K_s\to\PP_u^{m-1}$ を次数 $8$ の有限平坦被覆、$\Delta\to\PP_u^{m-1}$ を次数 $8^k$ の有限平坦被覆とし、ファイバー積分解
\[
  \Delta\simeq K_1\times_{\PP_u^{m-1}}\cdots\times_{\PP_u^{m-1}}K_k
\]
が成り立つ。さらに $m\geq0$ の場合を含め、空でない部分集合 $I\subset\{1,\ldots,k\}$ に対する部分臨界点スキーム $\Gamma_I=V(A_{s,i}\mid s\in I,\ 1\leq i\leq3)$ は、非空である限り型 $(2^{3|I|})$ の完全交叉であり、
\begin{equation}
\dim\Gamma_I=N-1-3|I|,\qquad\deg\Gamma_I=8^{|I|}
\end{equation}
を満たす。特に $m>0$ では $\dim\Delta=m-1$、$\deg\Delta=8^k$、$\omega_\Delta\simeq\OO_\Delta(3k-m)$ である（$m=0$ では $\Delta=\varnothing$）。
\end{lemma}
\begin{proof}
$R=\CC[u_1,\ldots,u_m]$、$Q_s=R[x_s]/(A_{s,1},A_{s,2},A_{s,3})$ とおく。$u=0$ でのファイバー $Q_s/(u)$ は滑らかな平面3次式 $f_{s,0}$ のJacobian環であり、そのHilbert級数は $(1+t)^3$ である。$Q_s/(u)$ の斉次基底をなす8個の元を $Q_s$ へ持ち上げれば、次数に関する帰納法から $Q_s$ は有限生成 $R$-加群となる。一方、$R[x_s]$ は次元 $m+3$ で方程式は3個なので $\dim Q_s\ge m$ であり、有限性と合わせて $\dim Q_s=m$ を得る。従って $A_{s,1},A_{s,2},A_{s,3}$ は正則列をなし、$Q_s$ はCohen--Macaulayである。

さらに $u_1,\ldots,u_m$ は $Q_s$ の斉次パラメータ系であるため $Q_s$-正則列となり、Koszul複体の完全性から $\operatorname{Tor}_1^R(Q_s,R/(u))=0$ が従う。上の8個の持上げが定める全射
\begin{equation}
R\oplus R(-1)^{\oplus3}\oplus R(-2)^{\oplus3}\oplus R(-3)\longrightarrow Q_s
\end{equation}
の核を $(u)$ で割ると $\operatorname{Tor}_1^R(Q_s,R/(u))=0$ により零となるため、次数付きNakayamaの補題から核自体が零、すなわちこの全射は同型である。複数ブロック $s\in I$ を合わせた座標環は $R$ 上のテンソル積なので階数 $8^{|I|}$ の有限自由 $R$-加群となり、選ばなかったブロックの変数を付け加えることで $\Gamma_I$ の高さと完全交叉性が従う（$m=0$ かつ $I=\{1,\ldots,k\}$ のときのみ共通零点がアフィン原点に限られ、射影スキームは空となる）。

$u=0$ 上では仮定\eqref{eq:pg-smooth-block}より各ブロック勾配の共通零点は原点のみなので、$K_s$ と $\Delta$ は $\{u=0\}$ と交わらない。各アフィン開集合 $\{u_\alpha\ne0\}$ で全ブロックを同時に $u_\alpha=1$ と規格化すれば、上で示した $R$ 上の有限自由性から被覆の有限平坦性とファイバー積表示が従い、重なりでも同時規格化により自然に貼り合わさる。次数および双対化層 $\omega_\Delta$ は、型 $(2^{3k})$ 完全交叉のHilbert級数と随伴公式から直ちに得られる。
\end{proof}

\subsection{退化集合の台と零点スキームの分解}
空でない部分集合 $L\subset\{1,\ldots,k\}$ に対し、楕円曲線 $E_s$（$s\in L$）の射影的joinを
\begin{equation}
J_L=V\bigl(u,\ x_t\ (t\notin L),\ f_{s,0}\ (s\in L)\bigr)
\end{equation}
とおく。$J_L$ は独立なアフィン錐の直積の射影化なので被約・既約であり、次元は $2|L|-1$ である。

\begin{proposition}\label{prop:pg-deg-support}
条件\eqref{eq:pg-smooth-block}の下で、$0\leq r<k$ を満たす各 $r$ について退化スキームの台は
\begin{equation}
|D_{2r}(\pi)|=\bigcup_{|I|=k-r}\Gamma_I\ \cup\ \bigcup_{|L|=r+1}J_L
\end{equation}
で与えられる（$X$ 上の台 $|D_{2r}(\pi_X)|$ は右辺と $X$ との交叉である）。特に零点スキーム（$r=0$）は、スキーム構造も含めて非交和
\begin{equation}
D_0(\pi)=\left(\coprod_{s=1}^k E_s\right)\sqcup\Delta,
\qquad D_0(\pi_X)=\left(\coprod_{s=1}^k E_s\right)\sqcup Z,
\qquad Z=\Delta\cap X
\end{equation}
に分解する。
\end{proposition}
\begin{proof}
アフィン錐上の点で勾配 $\nabla_{x_s}f_s$ が非零となるブロックの個数を $\ell$ とすると、錐上の双ベクトル場 $q$ の階数は $2\ell$ である。
\begin{itemize}
  \item $u\neq0$ の点では、Euler場 $E$ が中心方向の非零成分を持つため像 $\operatorname{im}q^\sharp$ に入らず、補題\ref{lem:pg-projective-rank}より射影階数は $2\ell$ のままである。従って $\operatorname{rank}\pi\le 2r\iff \ell\le r$、すなわち少なくとも $k-r$ 個のブロックで勾配が消える点（$\bigcup_{|I|=k-r}\Gamma_I$）である。
  \item $u=0$ の点では、滑らかさ\eqref{eq:pg-smooth-block}により $\nabla_{x_s}f_s=0\iff x_s=0$ である。また非零ブロックの像は $\ker df_{s,0}$ なので、$E\in\operatorname{im}q^\sharp$ となる条件は各非零ブロックで $3f_{s,0}=0$ が成り立つことである。従って射影階数が $2r$ 以下になるのは、$\ell\le r$（$\Gamma_I$ に属する）か、あるいは $\ell=r+1$ かつ全非零ブロックで $f_{s,0}=0$（すなわち $J_L$ に属する）のときにちょうど限る。
\end{itemize}
これにより退化集合の台の表示を得る。

次に零点スキームのスキーム構造を命題\ref{prop:pg-deg-ideal}（$j=1$）から確認する。開集合 $\{u_\alpha\ne0\}$ 上では、生成元 $u_\alpha A_{s,i}$ が全てのブロック偏微分 $A_{s,i}$ を生成し、残りの生成元も $(A_{s,i})$ に含まれるため、零点スキームは $\Delta$ に厳密に一致する。一方、$E_s$ の任意の点の近傍では、$E_s$ の滑らかさによりある $A_{s,i}$ が単元となる。すると式\eqref{eq:pg-deg-generators}の生成元は局所的に $u_\alpha$、$x_{t,l}$（$t\ne s$）、および $p_s$ を生成し、これらを法として $p_s\equiv 3f_{s,0}$ かつ他の生成元はすべて消える。従って零点イデアルは厳密に $(u,x_t\ (t\ne s),f_{s,0})$ に一致する。各 $E_s$ は被約であり、補題\ref{lem:pg-critical-cover}より $\Delta$ と交わらないため、求めるスキーム論的非交和分解を得る。
\end{proof}

なお、$m=0$ における $\Gamma_I$ は座標線形空間を台とするが、そのJacobianイデアルは一般に非被約な厚化を定める。

\subsection{残余零点スキーム $Z$ の完全交叉性と一般の滑らかさ}
\begin{proposition}\label{prop:pg-zero-ci}
条件\eqref{eq:pg-smooth-block}の下で、残余零点スキーム $Z=\Delta\cap X$ について次が成り立つ：
\begin{itemize}
  \item $m=1$ ならば $Z=\varnothing$ である。
  \item $m\geq2$ ならば $Z$ は型 $(2^{3k},3)$ の完全交叉であり、
  \begin{equation}
  \dim Z=m-2,\qquad \deg Z=3\cdot8^k,\qquad
  \omega_Z\simeq\OO_Z(3k+3-m)
  \end{equation}
  を満たす。
\end{itemize}
さらに、滑らかな3変数制限 $f_{s,0}$ を固定して中心変数 $u$ を含む残りの係数を一般に選ぶと、$K_s,\Delta$ および（$m\geq2$ における）$Z$ はすべて滑らかである。またこれらが正次元ならば、滑らかな場合には自動的に既約である。
\end{proposition}
\begin{proof}
\paragraph{ステップ1：完全交叉性と数値不変量}
$\Delta$ の斉次座標環を $Q$ とする。補題\ref{lem:pg-critical-cover}の証明より、$Q$ は $R=\CC[u]$ 上有限自由なCohen--Macaulay環であり、各最小素イデアルは $R$ の零イデアル $(0)$ へ縮約する（各成分は純次元性より次元 $m$ であり、有限性から $R$ への像も次元 $m$ だからである）。

背理法により、$F$ が $Q$ のある最小素イデアル $\mathfrak p$ に属すると仮定し、その関数体を $L=\operatorname{Frac}(Q/\mathfrak p)$ とおく。$L/\CC(u)$ は標数0の有限拡大ゆえ有限分離拡大であるため、$du_1,\ldots,du_m$ はK\"ahler微分加群 $\Omega^1_{L/\CC}$ の基底をなす。$\Delta$ 上では全てのブロック偏微分 $A_{s,i}$ が消えているため、$L$ 上で
\[
  0=dF=\sum_{\alpha=1}^m F_{u_\alpha}\,du_\alpha
\]
となり、全ての中心偏微分 $F_{u_\alpha}$ も $\mathfrak p$ 上で消える。ところがこの成分は $u$-空間を支配するためアフィン原点以外の点を持ち、そこで $F$ の全偏微分が消えて $X$ の滑らかさに反する。ゆえに $F$ は $Q$ のどの最小素イデアルにも属さず、Cohen--Macaulay性から $Q$ 上の非零因子である。これが $m\ge2$ での型 $(2^{3k},3)$ 完全交叉性と次元・次数・双対化層を与える。また $m=1$ のとき、もし $\Delta\cap X$ に点があれば $u_1\ne0$ かつ $3F=u_1 F_{u_1}=0$ より全偏微分が零となって $X$ の滑らかさに反するため、$Z=\varnothing$ である。

\paragraph{ステップ2：一般の係数に対する滑らかさと既約性}
$u$ を含む係数を動かすパラメータ空間を考える。$u_\alpha\ne0$ の任意の点において、単項式 $x_{si}u_\alpha^2$ の係数は偏微分 $A_{s,i}$ の値を互いに独立に動かすため、臨界点の族の全空間は滑らかである。$u=0$ には解がないため、この族はパラメータ空間と射影空間の積の閉部分多様体となり、パラメータ空間への射影は固有である。各ファイバーは非空な有限被覆ゆえ支配的であり、標数0における一般的滑らかさの定理から、一般の係数に対して $K_s$ および $\Delta$ は滑らかである。

そのような滑らかな $\Delta$ を固定し、$F$ に中心変数のみの3次式 $c(u)$ を加えてもブロック偏微分 $A_{s,i}$ は一切変わらない。切断 $c(u)|_\Delta$ の線形系は（$\Delta\cap\{u=0\}=\varnothing$ より）基点を持たないため、Bertiniの定理から一般の $c(u)$ に対して $Z=\Delta\cap X$ は滑らかとなる。また $X$ の滑らかさもFermat型 $\sum_s f_{s,0}+\sum_\alpha u_\alpha^3$ を含む非空開条件なので、これらすべてを同時に満たす一般の係数が存在する。最後に、正次元の射影完全交叉に対してはKoszul分解と線束の中間コホモロジー消滅から $H^0(\OO)=\CC$ が従うため連結であり、滑らかならば既約である。
\end{proof}

\begin{example}[Fermat 3次三様体の葉と零点スキーム]\label{ex:pg-fermat-threefold-leaves-zeros}
以上の結果を最も簡単な具体例 $n=3$、$k=1$、$m=2$ のFermat 3次三様体で確認しよう。
\begin{equation}\label{eq:pg-fermat-threefold}
F=\frac{x_1^3+x_2^3+x_3^3+u_1^3+u_2^3}{3},
\qquad
P=\partial_{x_1}\wedge\partial_{x_2}\wedge\partial_{x_3}
\end{equation}
が定める $X=V(F)\subset\PP^4$ 上のPoisson構造 $\pi_X$ を考える。第一積分は $[u_1:u_2]\in\PP^1$ であり、$u_2=cu_1$、$t=u_1$ とおけば、命題\ref{prop:pg-leaf-closure}により一般のsymplectic leafの閉包は
\begin{equation}\label{eq:pg-fermat-threefold-leaf}
V\bigl(x_1^3+x_2^3+x_3^3+(1+c^3)t^3\bigr)\subset\PP^3
\end{equation}
となる。$1+c^3\neq0$ ならば全偏微分の共通零点はアフィン原点だけなので、これは滑らかな3次曲面である。

一方、ブロック変数の偏微分は $A_{1,i}=x_i^2$ であるから、命題\ref{prop:pg-deg-support}より零点スキームは
\begin{equation}\label{eq:pg-fermat-threefold-zeros}
\begin{aligned}
D_0(\pi_X)&=E\sqcup Z,\\
E&=V(u_1,u_2,x_1^3+x_2^3+x_3^3),\\
Z&=V(x_1^2,x_2^2,x_3^2,u_1^3+u_2^3)\subset\PP^4
\end{aligned}
\end{equation}
に分解する。$E$ は滑らかな平面楕円曲線であり、残余スキーム $Z$ の台は中心直線 $\{x_1=x_2=x_3=0\}$ 上の3点である。各点 $p\in|Z|$ では $u_1\neq0$ かつ $1+(u_2/u_1)^3$ の根は単純なので、局所座標 $\xi_i=x_i/u_1$ を取ると局所環は
\begin{equation}\label{eq:pg-fermat-threefold-zero-ring}
\OO_{Z,p}\simeq
\CC[\xi_1,\xi_2,\xi_3]/(\xi_1^2,\xi_2^2,\xi_3^2)
\qquad(p\in|Z|)
\end{equation}
となる。従って各点の長さは $2^3=8$、$Z$ 全体の長さは $3\times8=24$ であり、命題\ref{prop:pg-zero-ci}の次数公式 $3\cdot8^k$ と厳密に一致する（一般のFermat型における厚化構造は命題\ref{prop:pg-fermat-thickening}で詳しく記述する）。
\end{example}

\section{3次四様体と持上げ障害}\label{sec:pg-fourfold}

本節では次元 $n=4$（$N=6$）の滑らかな3次四様体 $X\subset\PP^5$ に焦点を当て、一般点で非退化（一般階数 $4$）となるPoisson構造を完全に分類する。本節では記号を簡略化し、超曲面 $X$ 上のPoisson構造を単に $\pi$ と書く。

ここでの主結果は以下の3点である：(1) 一般階数4のPoisson構造は $(3,3)$ 分離型の3次四様体に限られ、すべて purely elliptic log symplectic となること（命題\ref{prop:pg-split-elliptic}）、(2) この条件が周囲の $\PP^5$ のHessian階数軌跡が平面を含むことと同値であること（系\ref{cor:pg-cubic-fourfold-elliptic}）、および(3) 非連結な零点曲線 $E_f\sqcup E_g$ に沿う3重消失を保った持上げの障害類を計算し、PymのarXiv v1の非存在主張との不整合を明確に説明すること（命題\ref{prop:pg-triple-lifting-obstruction}、系\ref{cor:pg-split-lifting-obstruction}、注意\ref{rem:pg-split-elliptic-lift-obstruction}）。

\subsection{分離型3次四様体の純楕円対数symplectic構造}
\label{subsec:pg-split-elliptic}
偶数次元 $2k$ の多様体上で一般点非退化なPoisson構造 $\pi$ に対し、そのPfaffian（最高外積 $\frac{1}{k!}\pi^{\wedge k}$）の零因子 $D$ が被約であるとき、逆形式 $\omega=\pi^{-1}$ を\emph{対数symplectic形式}（log symplectic form）と呼ぶ。局所体積形式 $\nu$ に対するmodularベクトル場を $Z_\nu=\delta_\nu\pi$ と書くと、$\pi$ の零点上における $Z_\nu$ の値は $\nu$ の選び方に依存しない（体積形式の変更による差がHamiltonベクトル場となり、$\pi=0$ の点で消えるため）。

4次元において、退化因子 $D$ が正規であり、かつ特異集合 $\operatorname{Sing}D$ 上の全ての点で $\pi=0$ および $Z_\nu\neq0$ が成り立つことが、PymのarXiv初版 \cite[Definition~4.2, Remark~4.4]{PPymElliptic} における \emph{purely elliptic} 条件である。以下では $(3,3)$ 分離型の3次四様体に対してこの条件を直接計算で確認する。

\begin{proposition}[分離型3次四様体の対数symplectic構造]
\label{prop:pg-split-elliptic}
$f\in\CC[x_0,x_1,x_2]_3$ および $g\in\CC[y_0,y_1,y_2]_3$ がともに滑らかな平面3次曲線を定めるとし、
\begin{equation}
\begin{aligned}
X&=V(f+g)\subset\PP^5,\\
E_f&=V(f,y_0,y_1,y_2),\qquad
E_g=V(x_0,x_1,x_2,g)
\end{aligned}
\label{eq:pg-split-fourfold}
\end{equation}
とおく。このとき $X$ は滑らかな3次四様体であり、2つの3変数Jacobian双ベクトル場
\[
  q_f=\iota_{df}(\partial_{x_0}\wedge\partial_{x_1}\wedge\partial_{x_2}),\qquad
  q_g=\iota_{dg}(\partial_{y_0}\wedge\partial_{y_1}\wedge\partial_{y_2})
\]
の和 $q_f+q_g$ は $X$ 上の一般階数 $4$ のPoisson構造 $\pi$ へ降下する。さらに次が成り立つ：
\begin{enumerate}
  \item 退化因子は2つの楕円曲線のjoin
  \begin{equation}
  D=V(f,g)=\operatorname{Join}(E_f,E_g)\subset X,
  \qquad
  (\operatorname{Sing}D)_{\mathrm{red}}=E_f\sqcup E_g
  \label{eq:pg-split-elliptic-divisor}
  \end{equation}
  であり、$D$ は被約・既約・正規な反標準因子で、各楕円曲線 $E_f,E_g$ に沿って重複度 $3$ を持つ。
  \item Poisson零点スキームは被約な非交和 $E_f\sqcup E_g$ に一致し、その各点でmodularベクトル場は零でない。
\end{enumerate}
従って $\pi$ はPym \cite[Definition~4.2]{PPymElliptic} の意味で \emph{purely elliptic log symplectic} である。
\end{proposition}
\begin{proof}
\paragraph{ステップ1：滑らかさ・降下・退化因子の正規性}
$f$ と $g$ の勾配はそれぞれ原点以外では同時に消えないため、$f+g$ の全偏微分の共通零点は原点のみであり、$X$ は滑らかである。$q_f$ と $q_g$ は変数が分離した2つのJacobian Poisson構造なので互いに可換であり、$f+g$ はそのCasimirかつ係数は2次斉次であるから、$q_f+q_g$ は $X$ 上のPoisson構造 $\pi$ へ降下する。

$D=V(f,g)\subset\PP^5$ のアフィン錐は2つの既約な平面3次錐の直積であるため、$D$ は既約・被約な型 $(3,3)$ の完全交叉であり、特にCohen--Macaulay（従って $S_2$）である。$x\ne0$ かつ $y\ne0$ の点では $(df,dg)$ の階数は $2$ であり、階数が $1$ 以下に下がるのは $x=0$ または $y=0$ のとき、すなわち $E_f\sqcup E_g$ 上のみである。$(\operatorname{Sing}D)_{\mathrm{red}}=E_f\sqcup E_g$ は $D$（次元3）の中で余次元 $2$ であるため、Serreの正規性判定法（$R_1$ かつ $S_2$）により $D$ は正規である。

\paragraph{ステップ2：$E_f$ 近傍での局所計算（Pfaffianとmodularベクトル場）}
$E_f$ の任意の点の近傍で局所形を求める（$E_g$ については $x$ と $y$ の役割を交換すれば同じである）。$x$ 側に行列式 $1$ の線形変換を施して、考えている点で $x_0\ne0$ かつ $f_{x_2}\ne0$ となるように取る。アフィンチャート $x_0=1$ 上で局所座標 $w=x_1/x_0$、$b=x_2/x_0$、$v_i=y_i/x_0$（$i=0,1,2$）を導入し、$X$ の方程式 $f(1,w,b)+g(v)=0$ を陰関数定理で $b=b(w,v)$ と解く。このとき $A:=f_{x_2}(1,w,b(w,v))$ は局所単元である。

次数0関数のPoisson括弧を式\eqref{eq:pg-projective-ratio-bracket}で直接計算すると、
\begin{equation}
\begin{aligned}
\{w,v_i\}&=A v_i,&
\{v_0,v_1\}&=g_{v_2},&
\{v_1,v_2\}&=g_{v_0},&
\{v_2,v_0\}&=g_{v_1},\\
\pi&=A\,\partial_w\wedge E_v+q_g(v),&
E_v&=\sum_{i=0}^2v_i\partial_{v_i}
\end{aligned}
\label{eq:pg-split-local-poisson}
\end{equation}
を得る。局所体積形式を $\nu=dw\wedge dv_0\wedge dv_1\wedge dv_2$ とすれば、Euler等式 $\sum_{i=0}^2 v_i g_{v_i}=3g(v)$ によりPfaffianは
\begin{equation}
\frac12\pi^{\wedge2}=3A\,g(v)\,\nu^{-1}
\label{eq:pg-split-pfaffian}
\end{equation}
と求まる。一方、$x\ne0$ かつ $y\ne0$ の点では補題\ref{lem:pg-projective-rank}により階数が4から2へ下がる条件はちょうど $f=g=0$ である。従って $D=V(f,g)$ が退化因子全体であり、式\eqref{eq:pg-split-pfaffian}より $D$ は重複度 $1$（被約）である。随伴公式より $K_X^{-1}\simeq\OO_X(3)$ であり、$D$ はその切断 $g|_X$ の零因子なので反標準因子である。また $D$ の局所方程式は3次式 $g(v)=0$ なので、$E_f=(v_0,v_1,v_2)$ に沿う重複度は $3$ である。4次元の歪対称行列の逆行列公式から $\omega=\pi^{-1}$ の極は高々1位（単純極）であり、Jacobi恒等式より $d\omega=0$、かつ最高外積が $K_X(D)$ の非零切断を与えるため、$\omega$ は対数symplectic形式である。

最後に零点スキームとmodularベクトル場を確認する。$A$ が単元であるため、式\eqref{eq:pg-split-local-poisson}の係数全体が生成するイデアルは $(v_0,v_1,v_2)$ に一致し、零点スキームは $E_f$ 上で被約である。また $A=f_{x_2}(1,w,b(w,v))$ の $v$ 依存性は $b(w,v)$ のみを経由し、$f(1,w,b)+g(v)=0$ の微分から $d_v b=-d_v g/A$ である。ゆえに $q_g$ と $d_v A$ の縮約は $\iota_{d_v g}q_g=0$ により消滅する。これを用いて $\pi$ の成分を微分すると、$\nu$ に関するmodularベクトル場 $Z_\nu=\delta_\nu\pi$ は
\begin{equation}
Z_\nu=\frac{\iota_{d(Ag)}\pi}{Ag}
=\iota_{d\log A}\pi-3A\,\partial_w,
\qquad
Z_\nu|_{E_f}=-3A(w,0)\partial_w\ne0
\label{eq:pg-split-modular}
\end{equation}
となる。$E_g$ 上でも同様に非零である。各特異点を通るsymplectic leafは0次元（$X$ 内の余次元4）であり、$D$ の正規性と合わせて purely elliptic log symplectic の全条件が満たされる。
\end{proof}

\begin{remark}[退化因子の特異スキームの厚み]
\label{rem:pg-split-singular-schemes}
退化因子 $D$ の特異スキーム（Jacobianイデアルが定めるスキーム）は、$E_f$ 近傍では $(g_{v_0},g_{v_1},g_{v_2})$ で定まり、滑らかな平面3次式の3つの2次偏微分が正則列をなすことから横断方向の長さは $2^3=8$ である（$E_g$ 側も同様）。これは被約であるPoisson零点スキームとは区別される。
\end{remark}

\Needspace{16\baselineskip}
\subsection{一般階数4とHessian軌跡の同値性}\label{subsec:pg-fourfold-hessian}
正規形定理とBricalli--Favale--Pirolaの結果を合わせると、滑らかな3次四様体が一般階数4のPoisson構造を持つ条件が完全に特徴づけられる。

\begin{corollary}[3次四様体の一般階数4とHessian軌跡]
\label{cor:pg-cubic-fourfold-elliptic}
滑らかな3次四様体 $X=V(F)\subset\PP(V)$（$\dim V=6$）に対し、周囲の $\PP(V)$ におけるHessianの階数 $3$ 以下の軌跡を
\begin{equation}\label{eq:pg-fourfold-hessian-rank-three}
\mathscr H_{\le3}(F)=
\{[v]\in\PP(V):\operatorname{rank}(\operatorname{Hess}F(v))\le3\}
\end{equation}
とおく。このとき次の3条件は同値である：
\begin{enumerate}
\item $X$ は一般階数 $4$ の正則Poisson構造を持つ。
\item 適切な射影座標を選んで $F=f(x_0,x_1,x_2)+g(y_0,y_1,y_2)$ と分離型に書ける（ここで $f,g$ はいずれも滑らかな平面3次曲線を定める）。
\item 軌跡 $\mathscr H_{\le3}(F)$ は線形平面 $\PP^2$ を含む。
\end{enumerate}
さらに、一般階数 $4$ の全てのPoisson構造は、適切な第2項の表示の下で命題\ref{prop:pg-split-elliptic}の形となり、例外なく purely elliptic log symplectic である。
\end{corollary}
\begin{proof}
$(1)\Rightarrow(2)$：定理\ref{thm:pg-normal}において $N=6$、一般階数 $2k=4$ とすると $k=2$、$m=0$ であり、$F=f(x_0,x_1,x_2)+g(y_0,y_1,y_2)$ と分離する。もし一方の平面3次曲線に特異点があれば、他方のブロック座標を零と置くことで $X$ の特異点が得られるため、$f,g$ はともに滑らかでなければならない。

$(2)\Rightarrow(1)$ および最後の主張は命題\ref{prop:pg-split-elliptic}そのものである。また $(2)\Leftrightarrow(3)$ は、Bricalli--Favale--Pirola \cite[Theorem~3.4, pp.~10--12]{PBricalliFavalePirola} において同論文の記法で $n=5$、$k=2$ とした場合の結果にほかならない。
\end{proof}

\begin{remark}\label{rem:pg-fourfold-ts-hessian-scope}
第3項の線形平面 $\PP^2$ は周囲の $\PP(V)$ 内にあり、分離表示では座標平面 $\PP^2_x$ と $\PP^2_y$ がこれに当たる（それらと $X$ との交叉が楕円曲線 $E_f,E_g$ である）。なお同文献のTheorem~Aは、滑らかな3次四様体のHessian超曲面が既約かつ正規であることと、任意のサイズのThom--Sebastiani分解を持たないことが同値であることを示している。従ってHessian超曲面の可約性や非正規性は、$(3,3)$ 型だけでなく $(1,5)$ 型や $(2,4)$ 型の分離からも生じ得る。
\end{remark}

実際に、Hessian超曲面が可約であっても一般階数4のPoisson構造を持たない3次四様体が存在することを具体例で確認しておこう。

\begin{example}[Hessianの可約性だけでは階数4を判定できない]
\label{ex:pg-reducible-hessian-no-rank-four}
$s=z_1+z_2+z_3+z_4$ とおき、$(1,1,4)$ 分離型の3次式
\begin{equation}\label{eq:pg-hessian-counterexample-cubic}
F=x_0^3+x_1^3+h(z),\qquad
h(z)=\sum_{i=1}^4z_i^3-s^3
\end{equation}
を考える。このとき $X=V(F)\subset\PP^5$ は滑らかでHessian超曲面は可約であるが、$X$ 上には一般階数 $4$ の正則Poisson構造は存在しない。

\paragraph{滑らかさとHessianの可約性}
$h$ の勾配が消えるとすると $z_i^2=s^2$（$1\le i\le4$）である。$s=0$ ならすべての $z_i=0$ となり、$s\ne0$ なら $z_i=\epsilon_i s$（$\epsilon_i\in\{1,-1\}$）から和を取って $\sum_{i=1}^4\epsilon_i=1$ を得るが、左辺は偶数なので矛盾する。よって $V(h)\subset\PP^3$ は滑らかな3次曲面であり、$X=V(F)$ も滑らかである。全成分が $1$ の $4\times4$ 行列を $J$ とすると、
\begin{equation}\label{eq:pg-hessian-counterexample-matrix}
\frac16\operatorname{Hess}F
=\operatorname{diag}\bigl(x_0,x_1,
        \operatorname{diag}(z_1,z_2,z_3,z_4)-sJ\bigr)
\end{equation}
である。右下の $4\times4$ ブロックの行列式は恒等的に零ではない（例えば $z_1=\cdots=z_4=1$ で値 $-15$ を取る）ため、$\det\operatorname{Hess}F$ は1次因子 $x_0 x_1$ を持ち可約である。

\paragraph{$(3,3)$ 分離の不可能性}
全ての点 $v\in\CC^6$ について $\operatorname{Hess}F(v)$ に関して自己共役となる定数行列の空間
\begin{equation}\label{eq:pg-hessian-counterexample-symmetrizer}
E_F=\{T\in\operatorname{Mat}_6(\CC):
        \operatorname{Hess}F(v)T=T^{\mathsf t}\operatorname{Hess}F(v)
        \text{ for all }v\}
\end{equation}
を計算する。$x_0,x_1$ の係数を比較すれば、$T$ の第1・第2行の非対角成分は零である。また $z_i$ の係数行列 $E_{ii}-J$（$1\le i\le4$）の共通核は零であるため（$(E_{ii}-J)v=0$ の $i$ 以外の成分から $\sum_j v_j=0$、さらに $i$ 成分から $v_i=0$ となる）、第1・第2列の非対角成分も零である。残る $4\times4$ ブロックを $T_4$、その第 $j$ 列の和を $c_j$ とおく。$(E_{ii}-J)T_4$ の対称性を $i$ と異なる添字 $k,l$ で比較すると $c_k=c_l$ を得るので、すべての $c_j$ は等しい。次に $(i,l)$ 成分と $(l,i)$ 成分を比較すると $(T_4)_{il}=0$（$l\ne i$）が従う。よって $T_4$ はスカラー行列であり、
\begin{equation}\label{eq:pg-hessian-counterexample-algebra}
E_F=\{\operatorname{diag}(a,b,cI_4):a,b,c\in\CC\},
\qquad
\{\operatorname{rank}T:T\in E_F,\ T^2=T\}
=\{0,1,2,4,5,6\}
\end{equation}
となる。線形座標変換 $x=Py$ に対して $\operatorname{Hess}(F\circ P)(y)=P^{\mathsf t}\operatorname{Hess}F(Py)P$ より $E_{F\circ P}=P^{-1}E_FP$ であるため、$E_F$ 内の冪等元の階数集合は座標変換で不変である。もし $F$ が $(3,3)$ 分離型に書けるならば、一方の3変数ブロックへの射影が $E_F$ の階数 $3$ の冪等元を与えるはずであるが、上の集合に $3$ は含まれない。従って系\ref{cor:pg-cubic-fourfold-elliptic}により、$X$ には一般階数 $4$ のPoisson構造が存在しない。
\end{example}

\Needspace{10\baselineskip}
\subsection{高次消失を保つ持上げの障害}\label{subsec:pg-lifting-obstruction}
命題\ref{prop:pg-split-elliptic}で構成した分離型3次四様体の零点曲線 $C=E_f\sqcup E_g$ は2つの連結成分を持つ。退化因子の方程式 $s=g|_X$ は $C$ に沿って3重に消えるが、この3重消失を保ったまま周囲の $\PP^5$ 上の3次多項式へ持ち上げようとすると、$C$ の非連結性に由来するコホモロジー障害が生じる。この障害こそが、PymのarXiv v1における3次四様体上の非存在主張と命題\ref{prop:pg-split-elliptic}との不整合の正確な原因である（注意\ref{rem:pg-split-elliptic-lift-obstruction}）。

まず一般の超曲面 $X\subset\PP^M$ と任意の消失次数 $r\ge2$ に対してこの持上げ障害群を決定し、続いて分離型3次四様体の障害類を具体的に計算する。

\begin{proposition}[高次消失を保つ持上げと連結成分の障害]
\label{prop:pg-triple-lifting-obstruction}
$X=V(F)\subset\mathbb P^M_{\mathbb C}$ を滑らかな次数 $d$ の超曲面とし、$i:X\hookrightarrow\mathbb P^M$ を閉埋込みとする。$C\subset X$ を空でない滑らかな閉部分多様体で、その全ての連結成分が正次元であるものとする。$\mathcal I\subset\mathcal O_{\mathbb P^M}$ および $\mathcal J\subset\mathcal O_X$ をそれぞれ $\mathbb P^M$ および $X$ における $C$ のイデアル層とする。このとき任意の整数 $r\ge2$ に対して、自然な同型
\begin{equation}\label{eq:pg-split-disconnected-cohomology}
 H^1(\mathbb P^M,\mathcal I^{r-1})
 \simeq \mathbb C^{\pi_0(C)}/\mathbb C\mathbf1
\end{equation}
および完全列
\begin{equation}\label{eq:pg-triple-lifting-sections}
 0\longrightarrow H^0(\mathbb P^M,\mathcal I^r(d))
 \longrightarrow H^0(X,\mathcal J^r(d))
 \xrightarrow{\ \delta_r\ }
 \mathbb C^{\pi_0(C)}/\mathbb C\mathbf1
\end{equation}
が存在する。従って、切断 $s\in H^0(X,\mathcal J^r(d))$ が $C$ に沿う $r$ 重消失を保って周囲の次数 $d$ の斉次多項式へ持ち上がることと $\delta_r(s)=0$ は同値であり、持上げが存在すれば一意である。特に $C$ が連結ならば、すべての $s$ が一意に持ち上がる。以下、$r=3$ のとき $\delta_3$ を単に $\delta$ と書く。
\end{proposition}

\begin{proof}
\paragraph{ステップ1：イデアル冪の短完全列}
$C$ と $X$ はともに滑らかなので、$C$ の各点の近傍で $F$ の局所方程式を $\mathcal I$ を生成する正則パラメータ系の一部として選べる。すると随伴次数環 $\operatorname{gr}_{\mathcal I}\mathcal O_{\mathbb P^M}\simeq\operatorname{Sym}_{\mathcal O_C}(\mathcal I/\mathcal I^2)$ において $F$ の次数1記号は局所的に自由生成元となり、その乗法は単射である。従って $(F)\cap\mathcal I^r=F\mathcal I^{r-1}$ が局所的に成り立ち、$\mathcal J=\mathcal I\mathcal O_X$ と合わせて短完全列
\begin{equation}\label{eq:pg-split-lifting-sequence}
 0\longrightarrow\mathcal I^{r-1}
 \xrightarrow{\ \cdot F\ }\mathcal I^r(d)
 \longrightarrow i_*\mathcal J^r(d)\longrightarrow0
\end{equation}
を得る。

\paragraph{ステップ2：厚化上の大域関数と $H^1(\mathbb P^M,\mathcal I^{r-1})$ の計算}
余法完全列と射影空間のEuler列を合成すると、ベクトル束の単射
\begin{equation}
 \mathcal I/\mathcal I^2
 \hookrightarrow\Omega^1_{\mathbb P^M}|_C
 \hookrightarrow\mathcal O_C(-1)^{\oplus(M+1)}
\end{equation}
を得る。これらの商はいずれも局所自由なので合成も部分束の埋込みであり、任意の $j\ge1$ について対称冪を取っても単射性が保たれる。$C$ の各連結成分が正次元の射影多様体であることから $H^0(C,\mathcal O_C(-j))=0$ であり、従って
\begin{equation}\label{eq:pg-conormal-powers-no-sections}
 \mathcal I^j/\mathcal I^{j+1}
 \simeq\operatorname{Sym}^j(\mathcal I/\mathcal I^2)
 \hookrightarrow\mathcal O_C(-j)^{\oplus\binom{M+j}{j}},
 \qquad H^0(C,\mathcal I^j/\mathcal I^{j+1})=0
\end{equation}
が成り立つ。

$q\ge0$ に対して厚化スキーム $C_{(q)}=V(\mathcal I^{q+1})$ を考えると、各 $q\ge1$ に対する短完全列
\[
  0\longrightarrow\mathcal I^q/\mathcal I^{q+1}
  \longrightarrow\mathcal O_{C_{(q)}}
  \longrightarrow\mathcal O_{C_{(q-1)}}\longrightarrow0
\]
と式\eqref{eq:pg-conormal-powers-no-sections}から、大域切断の制限写像 $H^0(C_{(q)},\mathcal O_{C_{(q)}})\to H^0(C_{(q-1)},\mathcal O_{C_{(q-1)}})$ は単射である。一方、全ての $C_{(q)}$ の基礎位相空間は $C$ に一致するため、各連結成分上の定数関数はそのまま厚化 $C_{(q)}$ 上の定数関数として延長し、この制限写像は全射でもある。帰納的に、全ての $q\ge0$ に対して同型
\begin{equation}\label{eq:pg-thickening-global-functions}
 H^0(C_{(q)},\mathcal O_{C_{(q)}})
 \xrightarrow{\sim} H^0(C,\mathcal O_C)
 =\mathbb C^{\pi_0(C)}\qquad(q\ge0)
\end{equation}
を得る。

$q=r-2$ として短完全列
\begin{equation}
 0\longrightarrow\mathcal I^{r-1}\longrightarrow\mathcal O_{\mathbb P^M}
 \longrightarrow\mathcal O_{C_{(r-2)}}\longrightarrow0
\end{equation}
のコホモロジー長完全列を取れば、$H^1(\mathbb P^M,\mathcal O_{\mathbb P^M})=0$ により境界写像が同型
\begin{equation}
 \partial:\mathbb C^{\pi_0(C)}/\mathbb C\mathbf1
 \xrightarrow{\sim}H^1(\mathbb P^M,\mathcal I^{r-1})
\end{equation}
を誘導する。同型\eqref{eq:pg-split-disconnected-cohomology}として $-\partial^{-1}$ を採用し、完全列\eqref{eq:pg-split-lifting-sequence}の連結準同型との合成を $\delta_r$ と定める。

\paragraph{ステップ3：障害類 $\delta_r(s)$ の明示的記述}
この符号規約の下では、$s\in H^0(X,\mathcal J^r(d))$ の周囲の $\PP^M$ への大域持上げ $h\in H^0(\mathbb P^M,\mathcal O(d))$（$H^1(\mathbb P^M,\mathcal O)=0$ により必ず存在する）と、$C$ の各連結成分近傍での局所持上げ $h_a\in\mathcal I^r(d)$ を取ったとき、商 $(h-h_a)/F\bmod\mathcal I^{r-1}$ が各連結成分上で取る定数値の組がちょうど $\delta_r(s)$ を表す（大域持上げ $h$ を $h+\lambda F$ に取り替えると全成分定数に同じ定数 $\lambda$ が加わるため、商空間 $\mathbb C^{\pi_0(C)}/\mathbb C\mathbf1$ の元としてwell-definedである）。最後に $C\ne\varnothing$ より $H^0(\mathbb P^M,\mathcal I^{r-1})=0$ なので、完全列\eqref{eq:pg-split-lifting-sequence}のコホモロジー長完全列から完全列\eqref{eq:pg-triple-lifting-sections}を得る。
\end{proof}

\begin{corollary}[分離型3次四様体の非零障害]
\label{cor:pg-split-lifting-obstruction}
$F=f(x_0,x_1,x_2)+g(y_0,y_1,y_2)$（$E_f=V(f)\subset\mathbb P^2_x$ および $E_g=V(g)\subset\mathbb P^2_y$ は滑らかな平面3次曲線）とし、$X=V(F)\subset\mathbb P^5$、$C=E_f\sqcup E_g\subset X$ とおく。このとき退化因子を定める切断 $s=g|_X\in H^0(X,\mathcal J^3(3))$ の障害類は
\begin{equation}
\label{eq:pg-split-obstruction-class}\delta(s)=[(0,1)]\ne0\quad\text{in }\mathbb C^2/\mathbb C(1,1)
\end{equation}
である。特に、両連結成分 $E_f,E_g$ に沿う3重消失を同時に保つような $\PP^5$ 上の3次持上げは存在しない。
\end{corollary}

\begin{proof}
$E_f$ の近傍では $s$ は $y$ の3次式 $g$ で表され、$E_g$ の近傍では $F=f+g=0$ より $s=-f|_X$（$x$ の3次式）と表されるため、$s\in H^0(X,\mathcal J^3(3))$ である。大域持上げとして $h=g$ を固定すると、$E_f$ 側では局所持上げ $h_a=g$、$E_g$ 側では局所持上げ $h_a=-f$ を取れる。従って $(h-h_a)/F$ の値は $E_f$ 上で $(g-g)/F=0$、$E_g$ 上で $(g-(-f))/F=F/F=1$ となり、$\delta(s)=[(0,1)]\ne0$ を得る。

この障害は極めて初等的に見ることもできる。実際、$s=g|_X$ の $\PP^5$ への大域的な3次持上げ全体は1パラメータ族
\begin{equation}\label{eq:pg-split-ambient-lift}
 H_\lambda=g+\lambda F=\lambda f+(1+\lambda)g\qquad(\lambda\in\CC)
\end{equation}
で尽くされる。$E_f$ に沿う3重消失は $\lambda=0$ を要求し、$E_g$ に沿う3重消失は $1+\lambda=0$（すなわち $\lambda=-1$）を要求するため、両者は決して両立しない。
\end{proof}

\begin{remark}[PymのarXiv v1との比較]
\label{rem:pg-split-elliptic-lift-obstruction}
Pym~\cite{PPymElliptic} のarXiv:1507.05668v1 のDefinition~4.2, p.~17 における purely elliptic 条件には、零点曲線 $C$ の連結性は含まれていない。命題\ref{prop:pg-split-elliptic}で構成した分離型3次四様体のPoisson構造はこの条件をすべて満たしており、同版Theorem~1.3, p.~4 および Proposition~5.7, p.~30 に述べられた3次四様体に関する非存在・非正規性の記述とは両立しない。

持上げを扱う同版Lemma~5.5, pp.~28--29 では、滑らかな中心 $C$ に連結性を仮定せず、$H^1(\PP^5,\mathcal I_C)=0$ を用いている。しかし非連結曲線 $C=E_f\sqcup E_g$ に対してはこのコホモロジー群は $H^1(\PP^5,\mathcal I_C)\simeq\CC\ne0$ であり、さらに命題\ref{prop:pg-triple-lifting-obstruction}と系\ref{cor:pg-split-lifting-obstruction}が示す通り $H^1(\PP^5,\mathcal I_C^2)\simeq\CC$ に属する非零の持上げ障害が現れる。従って同版の持上げ補題をこの非連結曲線に適用することはできない。なお、刊行版本文および訂正文における該当部分の扱いについては本稿では未確認である。
\end{remark}

% !TeX root = ../../../poisson-cubic.tex
\section{Poisson正規形から従うBondal予想}
\label{sec:pg-bondal}

Poisson構造$\pi$の退化スキームを
\begin{equation}
 D_{2j}(\pi)=V(\pi^{j+1})
\end{equation}
と書く。Bondal予想は、Fano多様体上では$0\le2j<\dim X$の各$j$について、
このスキームが次元$2j+1$以上の既約成分を持つと主張する
\cite[Conjecture~1.1]{BondalGualtieriPym}。

\begin{theorem}[滑らかな3次超曲面に対するBondal予想]
\label{thm:pg-bondal}
$X\subset\mathbb P^{n+1}_{\mathbb C}$を$n\ge3$の滑らかな3次超曲面、
$\pi$をそのPoisson構造とする。$0\le2j<n$ならば
\begin{equation}
 \dim D_{2j}(\pi)\ge2j+1.
\end{equation}
さらに、$\pi\ne0$の一般階数を$2k$とすると、$0\le j<k$の各$j$について、
$D_{2j}(\pi)$は次元$2j+1$の既約なjoin部分多様体を含む。
\end{theorem}

\begin{proof}
$\pi=0$、または$j\ge k$なら$D_{2j}(\pi)=X$なのでよい。
以下、定理\ref{thm:pg-normal}の正規形を用い、$0\le j<k$とする。
$r=j+1$個のブロックを選び、添字を$1,\ldots,r$とする。
$g_s(x_s)=f_s(x_s,0)$が非零なら、$V(g_s)\subset\mathbb P(W_s)$の
既約曲線成分$C_s$を選ぶ。$g_s=0$なら$C_s$を任意の直線とする。
独立な座標空間にあるこれらの曲線のjoin
\begin{equation}
 J_r=\operatorname{Join}(C_1,\ldots,C_r)
 \subset\mathbb P(W_1\oplus\cdots\oplus W_r)
\end{equation}
を考える。そのアフィン錐は$C_s$の錐の直積であるから、$J_r$は既約で
$\dim J_r=2r-1=2j+1$である。
その上では$u=0$、選ばなかったブロックの座標は零、各$g_s$も零である。
従って$F=\sum_sf_s=0$であり、$J_r\subset X$となる。

非零の錐上の点$z=(x_1,\ldots,x_r,0,\ldots,0;0)$で階数を調べる。
この点では$q=\sum_{s=1}^rq_s$で、各$q_s$の階数は$0$または$2$である。
いずれかが零なら$\operatorname{rank}q(z)\le2r-2$である。
全てが階数$2$なら
\begin{equation}
 \operatorname{im}q_s(z)^\sharp=\ker dg_s(x_s),
 \qquad dg_s(x_s)(x_s)=3g_s(x_s)=0.
\end{equation}
従ってEulerベクトル$E(z)=\sum_sx_s$は$q(z)^\sharp$の像に属する。
$F$はCasimirなので$q(z)\in\bigwedge^2T_zV(F)$であり、
$T_zV(F)/\mathbb CE(z)=T_{[z]}X$へ移ると、
補題\ref{lem:pg-projective-rank}により階数は$2$下がる。
いずれの場合も$\operatorname{rank}\pi_{[z]}\le2r-2=2j$となり、
$J_r\subset D_{2j}(\pi)$を得る。
$J_r$を含む既約成分の次元は少なくとも$2j+1$である。
\end{proof}

\begin{remark}\label{rem:pg-bondal-scope}
この証明は3変数制限の滑らかさを用いない。
滑らかな場合には$C_s$を平面3次曲線そのものに取れる。
退化スキームのイデアルは命題\ref{prop:pg-deg-ideal}で与えられる。
Fano4次元の場合は\cite[Theorem~7.13]{BondalGualtieriPym}、
5次元の場合は\cite[Corollary~1.3]{BondalDruelPereiraPymTouzet}にも含まれる。
\end{remark}

% !TeX root = ../../../poisson-cubic.tex
\section{Poisson構造の空間とモジュライスタック}
\label{sec:pg-poisson-stacks}

全ての対象を$\mathbb C$上で考える。Poisson構造のスカラーは割らず、零構造も含める。
まず超曲面を固定し、次に超曲面を動かす場合を扱う。

\subsection{固定した3次超曲面}
\begin{definition}\label{def:pg-pois-space}
滑らかな3次超曲面$X$のPoissonスキームを、2次写像の零スキーム
\begin{equation}\label{eq:pg-pois-space}
 \operatorname{Pois}(X)=Z\left(
 H^0(X,\bigwedge^2T_X)\longrightarrow H^0(X,\bigwedge^3T_X),
 \quad\pi\longmapsto[\pi,\pi]\right)
\end{equation}
と定める。
\end{definition}

\begin{definition}\label{def:pg-fixed-pois-stack}
$\operatorname{Aut}(X)$を$X$の自己同型群スキームとし、
$\phi\cdot\pi=\phi_*\pi$によって$\operatorname{Pois}(X)$に作用させる。
固定モジュライスタックを
\begin{equation}\label{eq:pg-fixed-pois-stack}
 \mathfrak{Pois}(X)=[\operatorname{Pois}(X)/\operatorname{Aut}(X)]
\end{equation}
と定める。点$\pi$の安定化子は
\begin{equation}\label{eq:pg-poisson-aut}
 \operatorname{Aut}(X,\pi)
 =\{\phi\in\operatorname{Aut}(X):\phi_*\pi=\pi\}
\end{equation}
である。
\end{definition}

\begin{proposition}\label{prop:pg-fixed-stack-algebraic}
$n\ge3$の滑らかな3次超曲面$X$について、
$\operatorname{Aut}(X)$は有限エタールであり、
$\mathfrak{Pois}(X)$は有限型Deligne--Mumfordスタックである。また
\begin{equation}\label{eq:pg-fixed-stack-components}
 \operatorname{Irr}(\mathfrak{Pois}(X))
 \simeq\operatorname{Irr}(\operatorname{Pois}(X))/\operatorname{Aut}(X).
\end{equation}
\end{proposition}

\begin{proof}
自己同型群に関する主張はMatsumura--Monsky\cite{PMatsumuraMonsky}の定理の特例である。
ここではgradientのsyzygyから証明する。
Grothendieck--Lefschetzの定理\cite{PGLefschetz}より
$\operatorname{Pic}(X)=\mathbb Z[\mathcal O_X(1)]$である。
従って自己同型は$\mathcal O_X(1)$を保ち、
$H^0(\mathbb P(V),\mathcal O(1))\simeq H^0(X,\mathcal O_X(1))$を通じて
射影線形変換に延長する。
ゆえに$\operatorname{Aut}(X)$は$\mathrm{PGL}(V)$内の$[F]$の閉安定化子である。

そのLie代数を計算する。Euler列を$-3$で捻ると
$H^0(T_{\mathbb P(V)}(-3))=H^1(T_{\mathbb P(V)}(-3))=0$なので、
$H^0(T_{\mathbb P(V)})\simeq H^0(T_{\mathbb P(V)}|_X)$である。
$X$に接する無限小射影変換は、1次式$\ell_i$と定数$c$により
\begin{equation}\label{eq:pg-stack-infinitesimal}
 \sum_i\ell_iF_i=cF
\end{equation}
と表される。Euler等式を引けば
$\sum_i(\ell_i-cz_i/3)F_i=0$となる。
補題\ref{lem:pg-gradient}による1次係数のsyzygyの消滅から
$\ell_i=cz_i/3$であり、このベクトル場は射影化すると零になる。
従って$\operatorname{Lie}\operatorname{Aut}(X)=0$である。
標数0の有限型群スキームは滑らかなので、
$\operatorname{Aut}(X)$は0次元のアフィン群スキーム、従って有限エタールである。

有限型アフィンスキーム$\operatorname{Pois}(X)$の有限エタール群による商は
有限型Deligne--Mumfordスタックである
\cite[Theorem~97.17.2, Tag~06FI]{PStacks}。
商へのアトラスは有限エタール全射であり、
一つの成分の像の逆像はその群軌道に属する成分の和である。
従って相異なる成分の像が一致するのは同じ群軌道に属するときに限り、
\eqref{eq:pg-fixed-stack-components}を得る。
\end{proof}

\subsection{超曲面を動かす相対モジュライ}
$N=n+2$、$V=\mathbb C^N$とし、
$U_N\subset\mathbb P(\operatorname{Sym}^3V^*)$を滑らかな3次式の開集合とする。
普遍超曲面を$p:\mathscr X_N\to U_N$、tautological lineを
$L=\mathcal O_{U_N}(-1)$と書く。
命題\ref{prop:pg-represent}の相対形は
\begin{equation}\label{eq:pg-relative-pois-bundle}
 \mathscr E_N:=p_*\bigwedge^2T_{\mathscr X_N/U_N}
 \simeq\bigwedge^3V\otimes L
\end{equation}
である。$[F]$上では$P\otimes F$を$\iota_{dF}P$に送る。
局所代表を$F$から$aF$へ替えると$P$は$a^{-1}P$へ替わるので、
この対応は重なりで一致する。

この表示とJacobi条件が非被約な基底への変更とも両立することを確認する。
$U_N$のアフィン開集合上で$F$の代表を選ぶと、各次数のgradient Koszul複体は
有限自由加群の複体となる。全ての幾何学的ファイバーで偏微分は2次正則列をなし、
Jacobian環のHilbert級数は
\begin{equation}\label{eq:pg-stack-gradient-hilbert}
 \operatorname{Hilb}_{\mathbb C[z]/(F_1,\ldots,F_N)}(t)=(1+t)^N
\end{equation}
である。従って正のKoszulホモロジーは全ファイバーで零で、
次数0のホモロジーの次元も一定である。
各複体の右端から順に次元を数えると全微分の階数が一定となる。
$U_N$は被約なので、微分を表す行列の所要の小行列式を可逆にした開集合では
像と核が直和因子となる。よって各次数の複体は局所的に分裂し、
任意の底変更後も完全性を保つ。
射影空間のEuler列に用いる線束のコホモロジーも任意の底変更と両立する。
従って命題\ref{prop:pg-represent}の表示、および
命題\ref{prop:pg-jacobi-extension}のEuler・divergence・低次数Koszul完全性による
Jacobi持上げは、任意の$U_N$-スキーム上で成立する。
ここで$F$の各次数の乗法写像も全ファイバーで単射なので局所分裂し、
任意の底変更後に単射である。
従って$F$による約分は可能であり、
式\eqref{eq:pg-euler-jacobi-divisibility}は
$0\to\OO(1)\xrightarrow{F}\OO(4)\to\OO_X(4)\to0$
の大域切断から得られる。
後者で現れる非零整数による除法は$\mathbb C$上で可逆である。

そこで、$\operatorname{Tot}(\mathscr E_N)$上の相対Jacobi零スキームを
\begin{equation}\label{eq:pg-relative-pois-space}
 \mathscr P_n=
 \{([F],\pi):[\pi,\pi]=0\}
 \subset\operatorname{Tot}(\mathscr E_N)
\end{equation}
と定める。局所代表$F$を選べば、これは
$[\iota_{dF}P,\iota_{dF}P]=0$の多項式係数によって定まり、
$P$について2次の有限個の方程式である。
上の底変更の議論により、各ファイバーは非被約構造も含めて
$\operatorname{Pois}(X_F)$と一致する。
代表の変更で括弧は変わらないため、これらの閉スキームは貼り合わさる。

\begin{definition}\label{def:pg-relative-pois-stack}
相対モジュライスタックを
\begin{equation}\label{eq:pg-relative-pois-stack}
 \mathfrak{Pois}^{\mathrm{rel}}_n
 =[\mathscr P_n/\mathrm{PGL}(V)]
\end{equation}
と定める。ここで$\mathrm{GL}(V)$のスカラー行列は
$\bigwedge^3V$と$L$に逆の重みで作用するため、作用は
$\mathrm{PGL}(V)$へ降りる。
このスタックは滑らかな3次超曲面とPoisson構造の組を、
Poisson構造を保つ同型で分類する。
\end{definition}

\begin{theorem}\label{thm:pg-relative-stack-algebraic}
$\mathfrak{Pois}^{\mathrm{rel}}_n$は有限型Deligne--Mumfordスタックであり、忘却射
\begin{equation}\label{eq:pg-relative-forget}
 \mathfrak{Pois}^{\mathrm{rel}}_n\longrightarrow
 \mathfrak{Cub}_n:=[U_N/\mathrm{PGL}(V)]
\end{equation}
は有限型アフィン射である。
$X$を指定した点$\operatorname{Spec}\mathbb C\to\mathfrak{Cub}_n$に沿う
二ファイバーは$\operatorname{Pois}(X)$であり、
その同型類の残余gerbeに沿うファイバーは$\mathfrak{Pois}(X)$である。
\end{theorem}

\begin{proof}
$\mathscr P_n$は有限型スキーム、$\mathrm{PGL}(V)$は滑らかな有限型代数群なので、
商は有限型代数的スタックである
\cite[Theorem~97.17.2, Tag~06FI]{PStacks}。
幾何学的点$(X,\pi)$の安定化子は有限エタール群
$\operatorname{Aut}(X)$の閉部分群であり、標数0なので有限エタールである。
従って対角は不分岐で、商はDeligne--Mumfordである。
忘却射をアトラス$U_N\to\mathfrak{Cub}_n$で底変更すると
$\mathscr P_n\to U_N$を得る。
この射はベクトル束の閉部分スキームの射影であるから有限型アフィンであり、
同じ性質が忘却射へ降下する。
二つのファイバーの記述は、前者では$X$への同型を選び、
後者ではその選択を$\operatorname{Aut}(X)$で割ることによる。
\end{proof}

\begin{remark}\label{rem:pg-stack-component-warning}
連結群$\mathrm{PGL}(V)$は$\mathscr P_n$の各既約成分を保つので、
$\mathfrak{Pois}^{\mathrm{rel}}_n$と$\mathscr P_n$の既約成分は一対一に対応する。
固定した$X$の有限群$\operatorname{Aut}(X)$による商では、
式\eqref{eq:pg-fixed-stack-components}のように成分が同一視され得る。
\end{remark}

% !TeX root = ../../../poisson-cubic.tex
\section{Fermat型と相対モジュライの既約成分}
\label{sec:pg-fermat-relative-components}

本節でいう既約成分は基礎位相空間の既約成分である。
$N=n+2\ge5$、$V=\mathbb C^N$とし、
\begin{equation}
 F_N=\sum_{i=1}^Nz_i^3,\qquad
 X_N=V(F_N),\qquad
 \mathfrak P_N^{\mathrm{Fer}}=\operatorname{Pois}(X_N)
\end{equation}
と置く。前節の相対Poissonスキームを$\mathscr P_n$、その非零部分を
$\mathscr P_n^\times$と書く。
$\left\{\begin{smallmatrix}N\\k\end{smallmatrix}\right\}_{\ge3}$は、
$\{1,\ldots,N\}$を大きさ$3$以上の$k$個の部分に分ける非順序集合分割の数とする。

\begin{theorem}\label{thm:pg-fermat-relative-count}
次が成り立つ。
\begin{enumerate}
\item $\mathfrak P_N^{\mathrm{Fer}}$の既約成分は、
各部分の大きさが$3$以上の集合分割で添字付けられる。$k$個の部分を持つ成分は
次元$3N-8k$、一般Poisson階数$2k$であり、
\begin{equation}\label{eq:pg-fermat-count}
 \#\operatorname{Irr}(\mathfrak P_N^{\mathrm{Fer}})
 =\sum_{k=1}^{\lfloor N/3\rfloor}
 \left\{\begin{smallmatrix}N\\k\end{smallmatrix}\right\}_{\ge3}.
\end{equation}
\item 固定モジュライスタック$\mathfrak{Pois}(X_N)$の既約成分は、
$N$の各部分が$3$以上の整数分割で添字付けられる。従って
\begin{equation}\label{eq:pg-fermat-stack-count}
 \#\operatorname{Irr}(\mathfrak{Pois}(X_N))
 =p_{\ge3}(N)=[q^N]\prod_{b\ge3}(1-q^b)^{-1}.
\end{equation}
\item 相対モジュライスタックには一般Poisson階数$2k$ごとに一つの既約成分があり、
\begin{equation}\label{eq:pg-relative-count}
 \#\operatorname{Irr}(\mathfrak{Pois}^{\mathrm{rel}}_n)=\lfloor N/3\rfloor.
\end{equation}
その次元は、$m=N-3k$として
\begin{equation}\label{eq:pg-relative-dim}
 d_{n,k}=k(3N-8)+\binom{m+2}{3}
 +k\left(10+6m+3\binom{m+1}{2}\right)-N^2
\end{equation}
である。
\end{enumerate}
零Poisson構造を除いても、これらの成分数は変わらない。
\end{theorem}

\subsection{固定Fermat超曲面}
\begin{proof}[第一項の証明]
Fermat式のHessianは$6\operatorname{diag}(z_1,\ldots,z_N)$であるから、
補題\ref{lem:pg-linearize}よりPoisson条件は
\begin{equation}\label{eq:pg-fermat-contractions}
 (\iota_{dz_i}P)\wedge(\iota_{dz_i}P)=0
 \qquad(1\le i\le N)
\end{equation}
を含む。定理\ref{thm:pg-normal}による分解
$P=\sum_{s=1}^kP_s$、$0\ne P_s\in\bigwedge^3W_s$を取る。
異なる$s,t$の積
$(\iota_{dz_i}P_s)\wedge(\iota_{dz_i}P_t)$は
$\bigwedge^2W_s\otimes\bigwedge^2W_t$という相異なる直和成分に属する。
従って\eqref{eq:pg-fermat-contractions}から、
各$dz_i$が非零となる$W_s$は高々一つである。
集合
\begin{equation}
 B_s=\{i:dz_i|_{W_s}\ne0\},\qquad
 V_B=\langle e_i:i\in B\rangle
\end{equation}
は互いに素で、$|B_s|\ge3$、$W_s\subset V_{B_s}$を満たす。
未使用の添字をいずれかの$B_s$に加えれば、$\{1,\ldots,N\}$の分割を得る。

逆に、そのような分割$\mathcal B$に対して
\begin{equation}
 Z_{\mathcal B}=\sum_{B\in\mathcal B}
 \widehat{\operatorname{Gr}}(3,V_B)\subset\bigwedge^3V
\end{equation}
と置く。各錐は独立な線形部分空間$\bigwedge^3V_B$に入るので、
この和は錐の直積であり閉既約である。
各点$P=\sum_BP_B$について、非零$P_B$の支持を$V_B$の基底に延長する。
この座標では$F_N=\sum_B\sum_{i\in B}z_i^3$は異なるブロックを混ぜず、
正規形定理の逆向きにより$P$はPoissonである。
従って有限個の$Z_{\mathcal B}$が$\mathfrak P_N^{\mathrm{Fer}}$の被約化を覆う。
各錐の次元は$3(|B|-3)+1$なので
\begin{equation}
 \dim Z_{\mathcal B}=3N-8|\mathcal B|.
\end{equation}

これらの閉既約部分集合の間に真の包含がないことを示す。
縮約写像$\bigwedge^2V^*\to V$の一般階数は
$Z_{\mathcal B}$上で$3|\mathcal B|$である。
階数が一定値以下となる条件は閉じているので、部分数が少ないものは多いものを含まない。
逆向きの包含は上の次元式から不可能である。
部分数が等しい二つの分割を比較するには、各支持$W_s$が
$B_s$の全座標方向に非零射影を持つ一般点を取る。
補題\ref{ne:blocks}により$P$から$W_s$の組が復元され、
従ってその座標支持$B_s$も復元される。
同じ部分数の別の分割へ入る余地はない。
従って各$Z_{\mathcal B}$が一つの既約成分である。
一般点のPoisson階数は正規形定理より$2|\mathcal B|$となる。
全ての錐は原点を含むので、原点は追加の成分を与えない。
\end{proof}

\begin{proof}[第二項の証明]
命題\ref{prop:pg-fixed-stack-algebraic}により成分の自己同型群軌道を数えればよい。
自己同型は$F_N(Az)=cF_N(z)$を満たす射影線形変換である。
Hessian行列式と連鎖律から
\begin{equation}\label{eq:pg-fermat-hessian-discriminant}
 \det\operatorname{Hess}(F_N)=6^N\prod_i z_i,
\end{equation}
\begin{equation}\label{eq:pg-fermat-hessian-transform}
 \prod_i(Az)_i=c^N\det(A)^{-2}\prod_i z_i
\end{equation}
を得る。一意分解性から$A$はmonomial行列であり、
$F_N(Az)=cF_N(z)$へ戻せば、スカラー行列を除いた対角成分は3乗根である。
従って
\begin{equation}\label{eq:pg-fermat-automorphism-group}
 \operatorname{Aut}(X_N)\simeq(\mu_3^N/\mu_3)\rtimes\mathfrak S_N.
\end{equation}
対角部分は各座標支持を保ち、$\mathfrak S_N$はその添字を置換する。
よって二つの集合分割が同じ軌道に属する条件は、部分の大きさの多重集合が等しいことである。
これは各部分が$3$以上の整数分割にほかならず、第二項を得る。
\end{proof}

\subsection{相対モジュライ}
\begin{proof}[第三項の証明]
一般Poisson階数$2k$の被約層を$\mathscr P_{n,k}$と書く。
正規形定理により、これは縮約写像$\bigwedge^2V^*\to V$の階数が$3k$となる
局所閉条件で定まる。まず3次式の非零代表を選んだ空間に移り、
\begin{equation}
 P_k=e_{123}+e_{456}+\cdots+e_{3k-2,3k-1,3k},\qquad
 \mathcal O_{N,k}=\mathrm{GL}(V)\cdot P_k
\end{equation}
と置く。$P_k$の支持を$W=\bigoplus_sW_s$、$V=W\oplus U$、$\dim U=m$とする。
補題\ref{ne:blocks}により、その安定化子$H_k$は$W_s$を置換する。
単位成分の元は各$W_s$上で$\mathrm{SL}(W_s)$に属し、
$U$上の可逆変換と$U\to W$の写像は任意である。
従って
\begin{equation}
 \dim H_k=8k+m^2+3km,\qquad
 \dim\mathcal O_{N,k}=k(3N-8).
\end{equation}

$P=P_k$に対するHessian条件\eqref{eq:pg-K}は、
異なる$W_s,W_t$間の混合2階微分の消滅と同値である。
実際、二つずつの共変ベクトルを各ブロックから選ぶとその混合成分が得られ、
逆向きはブロックごとのJacobian括弧から従う。
従って許される3次式の線形空間$\mathcal F_{N,k}$は
\begin{equation}
 F=c(u)+\sum_{s=1}^k
 \bigl(f_{s,3}(x_s)+f_{s,2}(x_s,u)+f_{s,1}(x_s,u)\bigr)
\end{equation}
で与えられる。ここで$f_{s,j}$の$x_s$次数は$j$、$u$次数は$3-j$であり、
\begin{equation}
 \ell_{N,k}:=\dim\mathcal F_{N,k}
 =\binom{m+2}{3}+k\left(10+6m+3\binom{m+1}{2}\right).
\end{equation}
滑らかな3次式の部分$\mathcal F_{N,k}^{\mathrm{sm}}$はFermat式を含む非空開集合である。
軌道上でHessian条件の線形写像は一定階数を持つので、その核は
$\mathcal O_{N,k}$上の階数$\ell_{N,k}$のベクトル束となる。
ここで$g\cdot F=F\circ g^{-1}$、$g\cdot P=(\bigwedge^3g)P$とする。
随伴束では$(gh,F_0)\sim(g,h\cdot F_0)$、$h\in H_k$と同一視する。
3次代表を選んだ被約階数層との比較射は
\begin{equation}
 \mathrm{GL}(V)\times^{H_k}\mathcal F_{N,k}^{\mathrm{sm}}
 \longrightarrow\widetilde{\mathscr P}_{N,k},\qquad
 [g,F_0]\longmapsto(g\cdot F_0,g\cdot P_k)
\end{equation}
である。$\mathrm{GL}(V)\to\mathcal O_{N,k}$のエタール局所での持上げを選べば逆射が得られ、
二つの持上げの差は$H_k$の作用なので逆射は降下する。
従ってこれは上の核ベクトル束の開部分との同型であり、
階数層は既約で次元$k(3N-8)+\ell_{N,k}$を持つ。

3次代表の変更は自由な作用
\begin{equation}
 a\cdot(F,P)=(aF,a^{-1}P),\qquad a\in\mathbb G_m,
\end{equation}
であり、$P\otimes F\in\bigwedge^3V\otimes L$を不変に保つ。
従って$\mathscr P_{n,k}$は既約で
$\dim\mathscr P_{n,k}=k(3N-8)+\ell_{N,k}-1$である。
$\mathrm{PGL}(V)$による商で$N^2-1$を引くと
\eqref{eq:pg-relative-dim}を得る。

各階数層の閉包が成分であることを示す。
縮約写像の階数条件から、小さい$k$の層の閉包は大きい$k$の層を含まない。
一方、隣り合う層の次元差は
\begin{equation}
 \dim\mathscr P_{n,k}-\dim\mathscr P_{n,k+1}
 =3(3k-1)(N-3k)+8>0
\end{equation}
なので、逆向きの包含もない。
零切断は、各$F$上の分解可能3重ベクトルをスカラー倍することで
$\mathscr P_{n,1}$の閉包に含まれる。
従って階数層の閉包が全ての既約成分を与え、その数は$\lfloor N/3\rfloor$である。
連結群$\mathrm{PGL}(V)$の商でも成分数は変わらない。
\end{proof}

\subsection{成分数の比較}
集合分割の指数母関数は
\begin{equation}\label{eq:pg-fermat-egf}
 \sum_{N\ge0}\frac{t^N}{N!}
 \left\{\begin{smallmatrix}N\\k\end{smallmatrix}\right\}_{\ge3}
 =\frac{(e^t-1-t-t^2/2)^k}{k!}.
\end{equation}
これと\eqref{eq:pg-fermat-stack-count}から次の値を得る。
\begin{center}
\small
\begin{tabular}{r|rrrr|rrr}
\toprule
$n$ & 階数$2$ & 階数$4$ & 階数$6$ & 階数$8$
& 固定空間 & 固定スタック & 相対スタック \\
\midrule
3 & 1 & -- & -- & -- & 1 & 1 & 1 \\
4 & 1 & 10 & -- & -- & 11 & 2 & 2 \\
5 & 1 & 35 & -- & -- & 36 & 2 & 2 \\
6 & 1 & 91 & -- & -- & 92 & 3 & 2 \\
7 & 1 & 210 & 280 & -- & 491 & 4 & 3 \\
8 & 1 & 456 & 2100 & -- & 2557 & 5 & 3 \\
9 & 1 & 957 & 10395 & -- & 11353 & 6 & 3 \\
10 & 1 & 1969 & 42735 & 15400 & 60105 & 9 & 4 \\
\bottomrule
\end{tabular}
\end{center}
中央の四列は固定Fermat空間の階数別成分数である。
例えば$N=9$では、階数$4$の$210$成分は$(3,6)$型の$84$個と
$(4,5)$型の$126$個からなり、固定スタックでは二つの成分になる。
一般の滑らかな3次超曲面の固定空間と固定スタックは、
系\ref{cor:pg-generic}によりともに既約である。


% !TeX root = ../../../poisson-cubic.tex
\section{非楕円型ブロックを持つ固定Poisson成分}\label{ne:examples}
3変数制限が特異でも、中心変数との混合項により超曲面全体は滑らかになり得る。
ここでは、節点型または尖点型の3変数制限を一般点に持つ固定Poisson成分を構成する。
座標$(x,y,z,a,b,c,u)$に対応するベクトル基底を$(e_0,\ldots,e_6)$とする。
\begin{equation}\label{ne:F}
\begin{aligned}
 f_{\mathrm N}&=xyz+y^3+z^3,& f_{\mathrm C}&=xz^2+y^3,\\
 F_\epsilon&=f_\epsilon+a^3+b^3+c^3+u(x^2+a^2+b^2+c^2)+u^3,
 &X_\epsilon&=V(F_\epsilon)\subset\PP^6
 \quad(\epsilon\in\{\mathrm N,\mathrm C\}).
\end{aligned}
\end{equation}
$P_1=e_0\wedge e_1\wedge e_2$、$P_2=e_3\wedge e_4\wedge e_5$、
$L=\CC P_1\oplus\CC P_2$とおく。

\begin{theorem}[非楕円型を一般点に持つ既約成分]\label{ne:main}
$X_{\mathrm N}$、$X_{\mathrm C}$は5次元の滑らかな3次超曲面である。
$P\mapsto\pi_P=(\iota_{dF_\epsilon}P)|_{X_\epsilon}$の下で
$L\simeq\mathbb A^2$は$\operatorname{Pois}(X_\epsilon)$の1つの既約成分の被約構造であり、
$P_0=P_1+P_2$でPoissonスキームそのものは滑らかである。
各成分の一般Poisson階数は4である。

$\lambda\mu\ne0$の点$\lambda P_1+\mu P_2$の3変数ブロック分解は順序を除き内在的である。
一方のブロックの平面3次曲線は、$X_{\mathrm N}$では節点型、$X_{\mathrm C}$では尖点型であり、
他方は滑らかなFermat3次曲線である。
特異なブロックを座標変換で滑らかな楕円ブロックへ取り替えることはできない。
\end{theorem}

\subsection{超曲面の滑らかさ}
\begin{lemma}\label{ne:smooth}
\eqref{ne:F}の全偏微分の共通零点は原点だけである。
\end{lemma}
\begin{proof}
$u=0$では$F_a=F_b=F_c=0$より$a=b=c=0$。
節点型の勾配$(yz,xz+3y^2,xy+3z^2)$が零なら$y=z=0$であり、
尖点型の勾配$(z^2,3y^2,2xz)$が零でも同様である。
残る$F_u=x^2$から$x=0$となる。

$u\ne0$では$u=1$に規格化する。
$a,b,c$はそれぞれ$0$または$-2/3$であり、その非零個数を$r$とすると
$a^2+b^2+c^2=4r/9$（$0\le r\le3$）。
節点型の残る3式は
\begin{equation}
 yz+2x=0,\qquad xz+3y^2=0,\qquad xy+3z^2=0.
\end{equation}
$yz=0$なら$x=y=z=0$。
$yz\ne0$なら後二式を掛けて$x^2=9yz=-18x$、従って$x=-18$。
よって$F_u=0$は、$3+4r/9=0$または$327+4r/9=0$を要求し、いずれも不可能である。
尖点型では$z^2+2x=0$、$3y^2=0$、$2xz=0$から$x=y=z=0$で、
$F_u=3+4r/9\ne0$となる。
\end{proof}

\subsection{Poissonスキームの接空間}
命題\ref{prop:pg-represent}により全双ベクトル空間は$\Lambda^3\CC^7$、次元35である。
命題\ref{prop:pg-jacobi-extension}のJacobi持上げと、勾配の1次syzygyの消滅は、
固定した$F_\epsilon$に対する線形写像と完全列を用いて示された。
これらは双対数環へのテンソル後も成立するため、Jacobi条件の線形化は
$\operatorname{Pois}(X_\epsilon)$の接空間を与える。
$Q=\sum_{i<j<k}q_{ijk}e_i\wedge e_j\wedge e_k$とする。
$Q^{ijk}$と$P_0^{ijk}$は、それぞれ$Q$と$P_0$の係数を添字の交代性で拡張したものとし、
$(F_\epsilon)_{br\ell}=\partial_b\partial_r\partial_\ell F_\epsilon$とおく。
Jacobi条件を$P_0$で線形化した整数行列$M_\epsilon$は
\begin{equation}\label{ne:M}
\begin{split}
 (M_\epsilon Q)_{ijk\alpha\ell}=\sum_{r,b}\bigl(&Q^{ir\alpha}P_0^{jkb}+P_0^{ir\alpha}Q^{jkb}
  +Q^{jr\alpha}P_0^{kib}+P_0^{jr\alpha}Q^{kib}\\
 &+Q^{kr\alpha}P_0^{ijb}+P_0^{kr\alpha}Q^{ijb}\bigr)(F_\epsilon)_{br\ell}
\end{split}
\end{equation}
で与えられる。$i<j<k$、$\alpha,\ell\in\{0,\ldots,6\}$であり、
行$(i,j,k,\alpha,\ell)$と列$(i,j,k)$を辞書式順で並べる。
零行を除くと節点型408行、尖点型312行となる。

\begin{lemma}[Jacobi方程式の線形化の階数]\label{ne:rank}
両行列の階数は33、核は$L$である。
列$(012)$と$(345)$を除いた行列には、次の行列式を持つ33次小行列
$M_{\mathrm N}'$、$M_{\mathrm C}'$が存在する：
\begin{equation}
 \det M_{\mathrm N}'=2^{24}3^{12}=8916100448256,\qquad
 \det M_{\mathrm C}'=-2^{33}3^{12}=-4565043429507072.
\end{equation}
\end{lemma}
\begin{proof}
$\lambda P_1+\mu P_2$は、中心変数$u$を共有する2つの3変数Jacobian括弧を与える。
ブロック間の括弧は零で、各ブロックの係数は他方の変数に依存しない。
よってJacobi恒等式が成り立ち、$(012)$、$(345)$の2列は零である。
従って階数は高々33であり、$L$は核に含まれる。

下界は整数上の小行列式によって得られる。
式\eqref{ne:M}から生成した$M_\epsilon$の転置を有理数上で行簡約し、
そのピボット列に対応する$M_\epsilon$の33行を選ぶ。
列$(012)$と$(345)$を除いた33次小行列の行列式を整数上で計算すると、
表示した2つの値となる。いずれも非零なので、階数は33、核は$L$である。
\end{proof}

選択した行・列の添字と小行列の全成分は、
\href{https://github.com/uedakazushi/poisson-cubic/tree/1f8d4c305a9c5c4eff4c21657254255a1a6ebf51/computations}{補足計算資料}の
\path{computations/nonelliptic_components_certificate.json}に収録している。
検算コード\path{computations/verify_nonelliptic_components.py}は
式\eqref{ne:M}から行列を再生成し、このデータと行列式の値を照合する。
この整数計算はPythonとSymPyを用い、配布ソースのルートで次を実行すれば再現できる。
\begin{quote}
\small\verb|python3 computations/verify_nonelliptic_components.py|
\end{quote}

\begin{proof}[定理\ref{ne:main}の成分性と階数]
$L$は2次元の閉既約部分スキームであり、その全点がPoissonである。
$P_0$での接空間は2次元なので、Poissonスキームの局所環はKrull次元2、
埋込み次元2の正則局所環となる。
これは$\CC$上有限型なので、$P_0$での滑らかさを意味する。
同じ次元を持つ閉部分$L$はその唯一の局所既約成分であり、局所的にスキームとして一致する。
従って$L$は大域的な既約成分で、一般点で被約である。

点$[0:1:0:-1:0:0:0]$は両$X_\epsilon$に属する。
$\lambda\mu\ne0$のとき、そこでアフィン双ベクトルの像は
$\langle e_0,e_2,e_4,e_5\rangle$で階数4。
Eulerベクトル$e_1-e_3$はその像に含まれないため、射影化しても階数4である。
ブロックの内在性には補題\ref{ne:blocks}を用いる。
\end{proof}

\begin{proof}[定理\ref{ne:main}のブロックに関する主張]
$\lambda\mu\ne0$では支持は$\langle e_0,e_1,e_2\rangle\oplus\langle e_3,e_4,e_5\rangle$。
補題\ref{ne:blocks}により2ブロックは一意である。
$F_\epsilon$の対応する射影平面への制限は$f_\epsilon$および$a^3+b^3+c^3$である。
$f_{\mathrm N}$は$[1:0:0]$近傍で$yz+y^3+z^3$となり節点、
$f_{\mathrm C}$は$z^2+y^3$となり尖点を持つ。
各式は$x$について1次で、その係数と定数項が互いに素なので既約である。
この特異性は内在的ブロックへの制限であり、座標の選択では除去できない。
\end{proof}

\subsection{特異ブロックと相対成分}
上の平面3次曲線の一般化Jacobianは次のようになる。
節点3次曲線では$\Pic^0\simeq\mathbb G_m$、尖点型では$\Pic^0\simeq\mathbb G_a$。
前者は正規化$\PP^1$上の2点のファイバーを非零スカラーで接着することで、
後者は$\CC[[t^2,t^3]]\subset\CC[[t]]$の単元商の1次係数$t$により得られる。

定理\ref{thm:pg-fermat-relative-count}により、$\lambda\mu\ne0$の点は
一般階数4の唯一の相対成分上にある。
その階数層では、ブロック制限の滑らかさはFermat式を含む非空開条件である。
従って本節の固定底の成分は、相対成分内の特異ブロックをもつ部分に位置する。
これらのブロックの量子化と行列式曲線は\cite{CubicQuantization}で扱う。

% !TeX root = ../../../poisson-cubic.tex
\section{特殊なブロック制限とFermat型}\label{sec:pg-special}

本節では、正規形\eqref{eq:pg-normal}において典型的な「滑らかな平面3次曲線」の仮定が成り立たない、あるいは退化が極端に集中する2つの重要な特殊ケースを調べる。第一に、あるブロックの3変数制限 $f_s(x_s,0)$ が恒等的に零となる場合に必要な中心変数の個数を決定する。第二に、Fermat型3次超曲面において零点スキームの残余部分 $Z$ がどのように非被約な厚化として現れるかを明示する。

\subsection{3変数制限が恒等的に零となる場合}
前節では3変数制限 $f_s(x_s,0)$ が特異な平面3次曲線（節点型・尖点型）になる例を構成した。さらに極端な状況として、$f_s(x_s,0)$ が恒等的に零になることも、中心変数の個数 $m$ が十分大きければ起こり得る。

\begin{proposition}\label{prop:pg-zero-restriction}
正規形\eqref{eq:pg-normal}で与えられる3次超曲面 $X$ が滑らかであり、あるブロック $s$ について3変数制限 $f_s(x_s,0)$ が恒等的に零であるとする。このとき $m\geq3$ でなければならない。逆に、$m=3$ の場合にはこのような滑らかな3次超曲面が存在する。
\end{proposition}
\begin{proof}
まず必要条件 $m\geq3$ を示す。$f_s(x_s,0)\equiv 0$ ならば、対応する座標平面 $\PP_s^2=\{u=0,\ x_t=0\ (t\neq s)\}$ 上で $F$ の全てのブロック変数に関する偏微分が恒等的に消える。残る偏微分は中心変数に関する高々 $m$ 個の2次斉次式 $F_{u_\alpha}(x_s,0)$（$1\le\alpha\le m$）だけである。もし $m\leq2$ ならば、$\PP_s^2$ 上でこれら高々2個の2次式の共通零点（すなわち3変数多項式環で高さ高々2のイデアルの非零零点）が必ず存在する。そのような点では $F$ の全偏微分が消えるため、$X$ の滑らかさに反する。ゆえに $m\geq3$ である。

次に $m=3$ での存在を具体例で示す。9変数 $a_1,a_2,a_3,b_1,b_2,b_3,u_1,u_2,u_3$ の3次式
\begin{equation}\label{eq:pg-zero-restriction-example}
F=\sum_{i=1}^3(u_i a_i^2+u_i^3)+b_1^3+b_2^3+b_3^3
\end{equation}
を考える。全偏微分が同時に零になるとすると、まず $F_{a_i}=2u_ia_i=0$ および $F_{u_i}=a_i^2+3u_i^2=0$ から各 $i$ について $a_i=u_i=0$ が従い、さらに $F_{b_i}=3b_i^2=0$ から $b_i=0$ となる。従って共通零点はアフィン原点だけであり、$X=V(F)\subset\PP^8$ は滑らかである。これは $k=2$、$m=3$ の正規形であり、$a$-ブロックの $u=0$ 制限は恒等的に零となっている。
\end{proof}

なお、例\eqref{eq:pg-zero-restriction-example}の退化スキームは、3変数制限の滑らかさを仮定しない一般的な生成元表示（式\eqref{eq:pg-deg-generators}）によってそのまま記述される。

\subsection{Fermat型における非被約な零点スキーム}
命題\ref{prop:pg-zero-ci}では、一般の係数に対して零点スキームの残余部分 $Z=\Delta\cap X$ が滑らかになることを見た。対照的に、Fermat型3次超曲面では臨界点被覆 $\Delta$ の全ての分岐が中心線形空間上に集中し、$Z$ は一様な厚みを持つ非被約スキームとなる。

正規形の座標を用いて、Fermat型の3次式を
\begin{equation}\label{eq:pg-fermat}
F=\frac13\left(\sum_{\alpha=1}^m u_\alpha^3+
\sum_{s=1}^k\sum_{i=1}^3x_{si}^3\right)
\end{equation}
と書く。
\begin{proposition}\label{prop:pg-fermat-thickening}
Fermat型3次式\eqref{eq:pg-fermat}において、$m>0$ ならば臨界点スキーム $\Delta=V(x_{si}^2\mid s,i)$ は中心線形空間 $\PP_u^{m-1}=\{x_{si}=0\}$ を台とし、横断方向の長さ $8^k$ を持つ厚化である。さらに $m\geq2$ ならば、残余零点スキーム $Z=\Delta\cap X$ は $\PP_u^{m-1}$ 内のFermat超曲面 $V(\sum_{\alpha=1}^m u_\alpha^3)$ の同じ横断長さ $8^k$ の厚化となる。一方、$m=0$ の場合には $\Delta=\varnothing$ である。
\end{proposition}
\begin{proof}
式\eqref{eq:pg-fermat}のブロック変数に関する偏微分は $A_{s,i}=x_{si}^2$ である。従って $\Delta$ の斉次座標環は
\begin{equation}
\CC[u]\otimes_\CC\bigotimes_{s=1}^k
\CC[x_{s1},x_{s2},x_{s3}]/(x_{s1}^2,x_{s2}^2,x_{s3}^2)
\end{equation}
と同型になる。各ブロックに対応するArtin局所環 $\CC[x_{s1},x_{s2},x_{s3}]/(x_{s1}^2,x_{s2}^2,x_{s3}^2)$ の長さは $2^3=8$ であるため、$k$ 個のテンソル積は横断長さ $8^k$ の厚化を与える。さらに、この商環では全ての $x_{si}^3$ が零になるため、$F$ の像は $\frac13\sum_{\alpha=1}^m u_\alpha^3$ に一致する。これが $m\ge2$ における $Z=\Delta\cap X$ の厚化構造を与える。最後に $m=0$ のとき、偏微分 $x_{si}^2=0$ の共通零点はアフィン原点のみであるため、射影スキーム $\Delta$ は空である。
\end{proof}

\appendix
\renewcommand{\theequation}{\Alph{section}.\arabic{equation}}
\renewcommand{\thetheorem}{\Alph{section}.\arabic{theorem}}
\renewcommand{\thesubsection}{\Alph{section}.\arabic{subsection}}
% !TeX root = ../../../poisson-cubic.tex
\section{射影完全交叉と斉次多様体のPoisson構造}
\label{sec:general-poisson-existence}

本付録では、序論で述べた「次元3以上の滑らかな射影完全交叉のうち、非零Poisson構造を持つものは $\PP^n$、2次超曲面、3次超曲面、$(2,2)$型完全交叉の4種類にちょうど限られる」という事実を自足的に証明し、あわせてトーリック多様体および滑らかな射影斉次多様体における非零Poisson構造の存在をまとめる。

完全交叉における非存在部分の要となる微分形式の低次数消滅は、Br\"uckmannの定理（Fu--Hwang \cite[Lemma~4.1]{FuHwangCI} に引用）の一部である。以下では、アフィン錐上の余法完全列の外積分解を用いる初等的かつ直接的な証明を与える。

\subsection{滑らかな射影完全交叉の存在判定}
まず、大域双ベクトル場の非存在を導くための微分形式の消滅補題を示す。

\begin{lemma}[微分形式の低次数消滅]\label{lem:ci-low-twist-vanishing}
$X\subset\PP^M_{\CC}$ を次元 $n\ge2$ の滑らかな射影完全交叉とする。このとき
\begin{equation}
H^0(X,\Omega_X^p(t))=0
\qquad(1\le p\le n-1,\ t\le p)
\label{eq:ci-vanishing}
\end{equation}
が成り立つ。
\end{lemma}
\begin{proof}
余次元を $c=M-n$ とし、$X$ を定める斉次正則列を $F_1,\ldots,F_c$ とする。$X$ の斉次座標環および頂点を除いたアフィン錐を
\begin{equation}
S=\CC[z_0,\ldots,z_M],\qquad
R=S/(F_1,\ldots,F_c),\qquad
U=\Spec R\setminus\{R_+\}
\label{eq:ci-punctured-cone}
\end{equation}
とおく。

\paragraph{ステップ1：錐 $U$ の構造層コホモロジー}
射影空間 $\PP^M$ 上の線束の中間コホモロジー消滅 \cite[III, Theorem~5.1]{PHartshorne} を $X$ のKoszul分解へ適用すると、すべての $\ell\in\mathbb Z$ に対して
\begin{equation}
H^0(X,\OO_X(\ell))=R_\ell,\qquad
H^i(X,\OO_X(\ell))=0\quad(1\le i\le n-1)
\label{eq:ci-line-bundle-cohomology}
\end{equation}
を得る（ただし $\ell<0$ では $R_\ell=0$ と規約する）。実際、長さ $c$ のKoszul分解のスペクトル系列で現れる中間コホモロジーの次数は $1$ 以上 $n-1+c=M-1$ 以下であり、すべて消滅する。射影 $q:U\to X$ はアフィンな $\mathbb G_m$-torsor であり、$q_*\OO_U=\bigoplus_{\ell\in\mathbb Z}\OO_X(\ell)$ が成り立つから、
\begin{equation}
\Gamma(U,\OO_U)=R,\qquad
H^i(U,\OO_U)=0\quad(1\le i\le n-1)
\label{eq:ci-cone-cohomology}
\end{equation}
を得る。

\paragraph{ステップ2：余法完全列の外積分解}
$U$ は滑らかな完全交叉であるため、余法完全列はベクトル束の短完全列
\begin{equation}
0\longrightarrow\mathcal N:=\bigoplus_{j=1}^c\OO_U\,dF_j
\longrightarrow\mathcal G:=\bigoplus_{a=0}^{M}\OO_U\,dz_a
\longrightarrow\Omega_U^1\longrightarrow0
\label{eq:ci-conormal-sequence}
\end{equation}
をなす。これに付随する $p$ 次外積の局所自由分解は
\begin{equation}
\begin{gathered}
\mathcal C_j=\operatorname{Sym}^j\mathcal N
        \otimes\bigwedge^{p-j}\mathcal G\quad(0\le j\le p),\\
0\longrightarrow\mathcal C_p\longrightarrow\cdots
\longrightarrow\mathcal C_0\longrightarrow\Omega_U^p\longrightarrow0
\end{gathered}
\label{eq:ci-exterior-conormal-resolution}
\end{equation}
である。各 $\mathcal C_j$ は同変構造を忘れれば $\OO_U$ の有限直和であるため、仮定 $p\le n-1$ と式\eqref{eq:ci-cone-cohomology}より $1\le i\le p$ に対して $H^i(U,\mathcal C_j)=0$ が成り立つ。従って\eqref{eq:ci-exterior-conormal-resolution}を短完全列に切ってコホモロジー長完全列を順に追えば、制限写像 $\Gamma(U,\bigwedge^p\mathcal G)\to\Gamma(U,\Omega_U^p)$ は全射である。この全射は錐のスカラー作用と可換であり、各 $dz_a$ の重みは $1$ なので、重み $t$ の成分を取り出すと全射
\begin{equation}
R_{t-p}\otimes\bigwedge^p(\CC^{M+1})^{\vee}
\twoheadrightarrow\Gamma(U,\Omega_U^p)_t
\label{eq:ci-cone-form-surjection}
\end{equation}
を得る。

\paragraph{ステップ3：水平形式への制限とCartanの公式}
$H^0(X,\Omega_X^p(t))$ の大域切断は、$U$ 上の重み $t$ の水平 $p$-形式 $\alpha$、すなわちEulerベクトル場 $E=\sum_{a=0}^M z_a\partial_{z_a}$ に対して $\iota_E\alpha=0$ かつ $\mathcal L_E\alpha=t\alpha$ を満たす切断に対応する。
\begin{itemize}
  \item $t<p$ のとき：$R_{t-p}=0$ なので式\eqref{eq:ci-cone-form-surjection}の左辺が零となり、$\alpha=0$ である。
  \item $t=p$ のとき：$R_0=\CC$ なので、\eqref{eq:ci-cone-form-surjection}により $\alpha$ は定数係数 $p$-形式の $U$ への制限であり、特に閉形式 $d\alpha=0$ である。よってCartanの魔法公式から
  \begin{equation}
    p\alpha=\mathcal L_E\alpha
    =d(\iota_E\alpha)+\iota_E(d\alpha)=0
    \label{eq:ci-cartan-vanishing}
  \end{equation}
  を得る。$p>0$ かつ標数 $0$ なので $\alpha=0$ である。
\end{itemize}
以上より主張が従う。
\end{proof}

この消滅補題と一般の線形射影を組み合わせることで、完全交叉のPoisson構造の存在判定が完成する。

\begin{theorem}\label{thm:ci-poisson-existence}
複素数体上の次元 $n\ge3$ の滑らかな射影完全交叉
\begin{equation}
X=X_{d_1,\ldots,d_c}\subset\PP^{n+c},
\qquad n\ge3,\quad d_j\ge2
\label{eq:ci-poisson-setup}
\end{equation}
（ただし1次方程式はあらかじめ除いている）について、次の3条件は同値である。
\begin{enumerate}[label=(\arabic*)]
\item $H^0(X,\bigwedge^2T_X)\ne0$。
\item $X$ は非零の正則Poisson構造を持つ。
\item 多重次数 $(d_1,\ldots,d_c)$ は空列（すなわち $X=\PP^n$）、$(2)$、$(3)$、$(2,2)$ のいずれかである。
\end{enumerate}
\end{theorem}
\begin{proof}
$(2)\Rightarrow(1)$ は定義より明らかである。

\paragraph{$(1)\Rightarrow(3)$ の証明}
$s=\sum_{j=1}^c(d_j-1)$ とおくと、随伴公式により
\begin{equation}
\omega_X\simeq\OO_X(-(n+1-s)),\qquad
\bigwedge^2T_X\simeq\Omega_X^{n-2}\otimes\omega_X^{-1}\simeq\Omega_X^{n-2}(n+1-s)
\label{eq:ci-forms}
\end{equation}
が成り立つ。$n\ge3$ なので $p=n-2$ は $1\le p\le n-1$ を満たす。もし $s\ge3$ ならば $t=n+1-s\le n-2=p$ となり、補題\ref{lem:ci-low-twist-vanishing}から $H^0(X,\bigwedge^2T_X)=0$ が従う。従って(1)が成り立つには $s\le2$ が必要であり、各 $d_j\ge2$ と合わせて多重次数は空列、$(2)$、$(3)$、$(2,2)$ の4種類に限られる。

\paragraph{$(3)\Rightarrow(2)$ の証明}
$s\le2$ と仮定し、一般の1次独立な1次形式 $L_0,\ldots,L_{n-2}$ を選ぶ。Euler場 $E=\sum_i z_i\partial_{z_i}$ を用いて、水平な $(n-2)$-形式
\begin{equation}
\alpha=\iota_E(dL_0\wedge\cdots\wedge dL_{n-2})
  \in H^0(X,\Omega_X^{n-2}(n-1))
\label{eq:ci-projective-form}
\end{equation}
を定義する。一般の選択に対して $\alpha|_X\ne0$ である。実際、滑らかな点 $x\in X$ と $L_0(x)\ne0$ を固定すると、射影埋込みの微分は単射なので $T_xX$ 上で1次独立な $n-2$ 個の微分 $d(L_i/L_0)$ を選ぶことができ、これは1次形式の組に関する非空開条件である。開集合 $\{L_0\ne0\}$ 上では
\begin{equation}
\alpha=L_0^{n-1}\,
d(L_1/L_0)\wedge\cdots\wedge d(L_{n-2}/L_0)
\label{eq:ci-local-form}
\end{equation}
と表され、一般点における $\alpha$ の核は線形射影 $X\dashrightarrow\PP^{n-2}$ のファイバーに接する階数2の可積分分布を定める。

ここで式\eqref{eq:ci-forms}のtwistの差 $2-s\ge0$ を補うため、$s=2$ ならそのまま $\alpha$ を、$s=1$ なら非零の1次切断を掛けたものを、$s=0$ なら非零の2次切断を掛けたものを、同型\eqref{eq:ci-forms}によって大域双ベクトル場 $\pi\in H^0(X,\bigwedge^2T_X)$ へ移す。いずれの場合も $\pi$ は一般階数2で非零であり、稠密開集合上では局所的に $\pi=a\,D_1\wedge D_2$ と書ける（$D_1,D_2$ は上の可積分分布の局所基底）。分布の可積分性から $[D_1,D_2]\in\langle D_1,D_2\rangle$ であるため、
\[
  [\pi,\pi]=2a^2\,D_1\wedge D_2\wedge[D_1,D_2]=0
\]
がこの稠密開集合上で成り立つ。Schouten括弧は大域的正則切断なので $X$ 全体で $[\pi,\pi]=0$ となり、非零Poisson構造を得る。
\end{proof}

\begin{remark}[低次元：曲線と曲面の場合]\label{rem:ci-small-dimension}
曲線（$n=1$）では $\bigwedge^2T_X=0$ であり、非零Poisson構造は存在しない。曲面（$n=2$）では $\bigwedge^2T_X=\omega_X^{-1}$ かつSchouten条件 $[\pi,\pi]=0$ が自動的に成立する。従って滑らかな完全交叉曲面が非零Poisson構造を持つ条件は $s\le3$ に緩み、上の4型に加えて $(4)$、$(2,3)$、$(2,2,2)$ のK3完全交叉曲面が現れる。
\end{remark}

\subsection{トーリック多様体}
完全交叉から離れて対称性の高い多様体に目を向けると、2次元以上の可換なトーラス作用を持つ多様体では、可換な基本ベクトル場の外積によって直ちに非零Poisson構造が構成できる。

\begin{proposition}\label{prop:toric-poisson}
次元 $2$ 以上の複素正規トーリック多様体 $X$ には非零の代数的Poisson括弧が存在する。特に $X$ が滑らかならば、それは大域双ベクトル場 $\pi\in H^0(X,\bigwedge^2T_X)$ で与えられる。
\end{proposition}
\begin{proof}
$X$ に含まれる稠密トーラス $T\simeq(\CC^*)^n$（$n\ge2$）のLie代数から1次独立な2元を取り、それらが誘導する大域導分を $D_1,D_2$ とする。$T$ は可換なので $[D_1,D_2]=0$ であり、$\pi=D_1\wedge D_2$ は自動的にJacobi条件 $[\pi,\pi]=0$ を満たす。$X$ を覆う各 $T$-不変アフィン開集合上で括弧 $\{f,g\}=D_1(f)D_2(g)-D_2(f)D_1(g)$ は正則かつ貼り合わさり、稠密トーラス上では $D_1,D_2$ が1次独立なので非零である。
\end{proof}

\subsection{滑らかな射影斉次多様体}
同じ可換ベクトル場の構成は、一般の滑らかな射影斉次多様体にも適用できる。

\begin{theorem}\label{thm:homogeneous-poisson}
$\CC$ 上の次元 $2$ 以上の滑らかな連結射影斉次多様体 $X$ は、常に非零の正則Poisson構造を持つ。
\end{theorem}
\begin{proof}
Borel--Remmertの定理 \cite[Theorem~I]{BorelRemmert} により、$X$ はアーベル多様体 $A$ と半単純代数群の旗多様体の直積 $G/P$ の積
\[
  X\simeq A\times G/P
\]
に分解される。$A$ の並進ベクトル場同士、および $G/P$ の最大トーラスが誘導する基本ベクトル場同士は互いに可換であり、異なる因子間のベクトル場も可換である。従って一般点で1次独立な2本の大域可換ベクトル場 $D_1,D_2$ を選べれば、$\pi=D_1\wedge D_2$ が求める非零正則Poisson構造を与える。その存在を因子ごとに確認しよう。
\begin{itemize}
  \item $G/P$ に階数 $2$ 以上の非自明な単純因子が含まれる場合：$P$ に含まれない単純根 $\alpha$ と、Dynkin図でそれに隣接する単純根 $\beta$ を選ぶと、$\alpha$ と $\alpha+\beta$ は接空間に現れる1次独立な2つのトーラス重みを与える。反対のBorel部分群の開軌道上では、根部分群が与える座標に関してトーラス作用がこれらの重みで対角化されるため、最大トーラスの一般軌道は2次元以上となり、独立な2本の基本ベクトル場が得られる。
  \item $G/P$ の非自明な単純因子がすべて階数 $1$（すなわち $\PP^1$）の場合：そのような因子が2個以上あれば、それぞれのトーラス作用から独立な2方向が得られる。
  \item 残る場合（$G/P$ が点または単一の $\PP^1$ のとき）：$\dim X\ge2$ より $\dim A\ge1$（$G/P$ が点なら $\dim A\ge2$）であるため、$A$ の並進ベクトル場から必要な数だけ独立な方向を補える。
\end{itemize}
いずれの場合も一般点で1次独立な可換大域ベクトル場 $D_1,D_2$ が存在し、$\pi=D_1\wedge D_2$ は非零かつ $[\pi,\pi]=0$ を満たす。
\end{proof}

% !TeX root = ../../poisson-cubic.tex
% Bibliography for the Poisson manuscript.
\begin{thebibliography}{99}
\addcontentsline{toc}{section}{参考文献}
\small
\bibitem{PHartshorne}
R.~Hartshorne,
\emph{Algebraic Geometry}, Graduate Texts in Mathematics 52,
Springer, 1977.
\href{https://doi.org/10.1007/978-1-4757-3849-0}{doi:10.1007/978-1-4757-3849-0}.

\bibitem{PEtingofLie}
P.~Etingof,
\emph{Lie Groups and Lie Algebras}, MIT lecture notes, 2024 course version.
\href{https://live.ocw.mit.edu/courses/18-755-lie-groups-and-lie-algebras-ii-spring-2024/mit18_755_s24_lec_full.pdf}{MIT OpenCourseWare 公開原稿}.
本文ではTheorem~11.16, p.~66、Proposition~16.9, p.~88、
Lemma~19.16, p.~104、Corollary~19.18とTheorem~19.19, pp.~105--106、
Definition~21.13, p.~114、および\S32.2を用いる。

\bibitem{PDerdzinskiGal}
A.~Derdzinski and \'{S}.~R.~Gal,
\emph{Curvature spectra of simple Lie groups},
Abhandlungen aus dem Mathematischen Seminar der Universit\"at Hamburg
\textbf{83} (2013), no.~2, 219--230.
\url{https://doi.org/10.1007/s12188-013-0085-z};
\href{https://arxiv.org/abs/1304.2801}{arXiv:1304.2801}.
参照箇所は著者最終原稿の(1.4), p.~2、Theorem~B, p.~3、
Lemma~7.1(a), p.~10。
\url{https://people.math.osu.edu/derdzinski.1/preprints/amsuh14.pdf}.

\bibitem{PBondal93}
A.~I.~Bondal,
\emph{Non-commutative deformations and Poisson brackets on projective spaces},
Max-Planck-Institut f\"ur Mathematik Preprint 93-67 (1993).
\S4、印刷pp.~11--12のEuler--divergence恒等式と無番号Theoremを参照。
\url{https://archive.mpim-bonn.mpg.de/id/eprint/1398/1/preprint_1993_67.pdf}.

\bibitem{PPymElliptic}
B.~Pym,
\emph{Elliptic singularities on log symplectic manifolds and
Feigin--Odesskii Poisson brackets},
Compositio Mathematica \textbf{153} (2017), no.~4, 717--744.
\url{https://doi.org/10.1112/S0010437X16008174};
\href{https://arxiv.org/abs/1507.05668v1}{arXiv:1507.05668v1}.
本文の参照版はarXiv v1：
Definition~4.2, p.~17、Theorem~1.3, p.~4、
Lemma~5.5, pp.~28--29、Proposition~5.7, p.~30。

\bibitem{PBricalliFavalePirola}
D.~Bricalli, F.~F.~Favale and G.~P.~Pirola,
\emph{On the irreducibility of Hessian loci of cubic hypersurfaces},
Forum of Mathematics, Sigma \textbf{13} (2025), e80, 1--36.
\href{https://doi.org/10.1017/fms.2025.36}{doi:10.1017/fms.2025.36};
\href{https://arxiv.org/abs/2406.12024}{arXiv:2406.12024}.
参照箇所は刊行版Theorem~A, p.~2、Theorem~B, p.~3、
Theorem~3.4とその証明, pp.~10--12。

\bibitem{PMayanskiy}
E.~Mayanskiy,
Poisson cohomology of two Fano threefolds,
\emph{Journal of Algebra} \textbf{424} (2015), 21--45.
\href{https://doi.org/10.1016/j.jalgebra.2014.08.049}{doi:10.1016/j.jalgebra.2014.08.049};
\href{https://arxiv.org/abs/1303.5916}{arXiv:1303.5916}.
本文で用いるGrassmann記述はTheorem~1（原稿p.~2）。
同定理および原稿p.~3は分類をLoray--Pereira--Touzetに帰属させている。

\bibitem{BondalGualtieriPym}
M.~Gualtieri and B.~Pym,
\emph{Poisson modules and degeneracy loci},
Proc. Lond. Math. Soc. \textbf{107} (2013), 627--654.
\href{https://doi.org/10.1112/plms/pds090}{doi:10.1112/plms/pds090};
\href{https://arxiv.org/abs/1203.4293}{arXiv:1203.4293}.
Bondal予想の定式化、四次元Fanoの場合の証明。

\bibitem{BondalDruelPereiraPymTouzet}
S.~Druel, J.~V.~Pereira, B.~Pym and F.~Touzet,
\emph{Bondal's conjecture in dimension five},
\href{https://arxiv.org/abs/2606.03424}{arXiv:2606.03424} (2026).
定理番号はarXiv v1による：Theorem~1.2（奇数次元のsubmaximal退化軌跡）、
Corollary~1.3（Fano五次元）。
\url{https://druel.perso.math.cnrs.fr/textes/bondal5d.pdf}.

\bibitem{FuHwangCI}
B. Fu and J.-M. Hwang,
\emph{Isotrivial VMRT-structures of complete intersection type},
Asian Journal of Mathematics \textbf{22} (2018), 333--354.
Lemma~4.1（Br\"uckmannの消滅定理を引用）。
\url{https://arxiv.org/abs/1608.00846v2}.
\url{https://doi.org/10.4310/AJM.2018.v22.n2.a10}.

\bibitem{BorelRemmert}
A. Borel and R. Remmert,
\emph{\"Uber kompakte homogene K\"ahlersche Mannigfaltigkeiten},
Mathematische Annalen \textbf{145} (1962), 429--439, Theorem I.
英訳：\url{https://wglobke.github.io/publications/BorelRemmert1962.pdf}.

\bibitem{LPTPoissonCubic}
F.~Loray, J.~V.~Pereira and F.~Touzet,
\emph{Foliations with trivial canonical bundle on Fano 3-folds},
Math. Nachr. \textbf{286} (2013), 921--940.
\href{https://arxiv.org/abs/1112.3871}{arXiv:1112.3871}.

\bibitem{PStacks}
The Stacks Project Authors,
\emph{The Stacks Project},
Cohen--Macaulay環・Koszul複体・射影空間のコホモロジー・商スタック。
参照タグ：
\href{https://stacks.math.columbia.edu/tag/00ND}{00ND},
\href{https://stacks.math.columbia.edu/tag/02JN}{02JN},
\href{https://stacks.math.columbia.edu/tag/062F}{062F},
\href{https://stacks.math.columbia.edu/tag/01XT}{01XT},
\href{https://stacks.math.columbia.edu/tag/06FI}{06FI}.

\bibitem{PGLefschetz}
A.~Grothendieck,
\emph{Cohomologie locale des faisceaux coh\'erents et th\'eor\`emes
de Lefschetz locaux et globaux (SGA~2)},
Expos\'e~XII, Corollaire~3.7 (Picard groups of sufficiently
high-dimensional smooth projective complete intersections).
\url{https://grothendiecksga.com/read/sga2/en/XII-3.html}.

\bibitem{PMatsumuraMonsky}
H.~Matsumura and P.~Monsky,
\emph{On the automorphisms of hypersurfaces},
J. Math. Kyoto Univ. \textbf{3} (1963/64), 347--361.
\url{https://doi.org/10.1215/kjm/1250524785}.

\bibitem{CubicQuantization}
\emph{滑らかな三次超曲面の代数的量子化},
研究原稿、2026年11月21日。
\href{https://github.com/uedakazushi/cubic-quantization/blob/main/papers/cubic-quantization.tex}{GitHub公開ソース（主ファイル）}.
本稿のPoisson正規形を入力とする代数的中心元・PBW性・同変降下を扱う独立論文。
\end{thebibliography}

\end{document}
